\documentclass[12pt,a4paper]{article}
\usepackage[utf8]{inputenc}
\usepackage[T1]{fontenc}
\usepackage[italian]{babel}
\usepackage{geometry}
\usepackage{listings}
\usepackage{xcolor}
\usepackage{amsmath}
\usepackage{amssymb}
\usepackage{hyperref}
\usepackage{titlesec}
\usepackage{parskip}
\usepackage{booktabs}
\usepackage{array}
\usepackage{enumitem}
\usepackage{fancyhdr}
\usepackage{tcolorbox}
\usepackage{mdframed}
\tcbuselibrary{skins, breakable}
\usepackage{amsthm}
\usepackage{hyperref}
\usepackage{tikz-cd}

\geometry{top=2.5cm, bottom=2.5cm, left=3cm, right=3cm}

% Stile per il codice Python
\definecolor{codegreen}{rgb}{0,0.6,0}
\definecolor{codegray}{rgb}{0.5,0.5,0.5}
\definecolor{codepurple}{rgb}{0.58,0,0.82}
\definecolor{backcolour}{rgb}{0.96,0.96,0.96}
\definecolor{keywordcolor}{rgb}{0.0,0.3,0.7}
\definecolor{stringcolor}{rgb}{0.7,0.1,0.1}

\lstdefinestyle{pythonstyle}{
    backgroundcolor=\color{backcolour},
    commentstyle=\color{codegreen}\itshape,
    keywordstyle=\color{keywordcolor}\bfseries,
    numberstyle=\tiny\color{codegray},
    stringstyle=\color{stringcolor},
    basicstyle=\ttfamily\small,
    breakatwhitespace=false,
    breaklines=true,
    captionpos=b,
    keepspaces=true,
    numbers=left,
    numbersep=5pt,
    showspaces=false,
    showstringspaces=false,
    showtabs=false,
    tabsize=2,
    frame=single,
    rulecolor=\color{gray!40},
    framesep=5pt
}
\lstset{style=pythonstyle}

% Box per note e avvisi
\newtcolorbox{notebox}{
  colback=yellow!10,
  colframe=yellow!60!black,
  fonttitle=\bfseries,
  title=Nota,
  breakable
}

\newtcolorbox{warningbox}{
  colback=orange!10,
  colframe=orange!80!black,
  fonttitle=\bfseries,
  title=Attenzione,
  breakable
}

\newtcolorbox{tipbox}{
  colback=green!8,
  colframe=green!60!black,
  fonttitle=\bfseries,
  title=Suggerimento,
  breakable
}

\pagestyle{fancy}
\fancyhf{}
\rhead{\textcolor{black}{Scienze matematiche per l'intelligenza artificiale}}
\lhead{\textcolor{black}{Alessandro Lo Curcio}}
\cfoot{\thepage}
\renewcommand{\headrulewidth}{0.4pt}

% Definizione del nuovo colore per la copertina (blu scuro)
\definecolor{coverdarkblue}{RGB}{20, 40, 70}

\begin{document}

% ── PAGINA DI COPERTINA PER CALCOLO SCIENTIFICO (PDE) ────────────────────────
\begin{titlepage}
  \pagecolor{coverdarkblue}
  \color{white}
  \vspace*{1.5cm}
  \begin{center}
    {\large\scshape Università degli Studi La Sapienza, Roma}\\[0.6cm]
    {\small Alessandro Lo Curcio \quad $\bullet$ \quad \today}\\[1.5cm]

    % Disegno con elementi iconici del Calcolo Scientifico (PDE, griglie, domini)
    \begin{tikzpicture}[scale=1.2]
  % Riquadro unico semplice
  \draw[white!70, thick] (0,0) rectangle (8,5);
  
  % Linea di base (dominio)
  \draw[white!50, thick] (1,1) -- (7,1);
  
  % Nodi della discretizzazione
  \foreach \x in {1.5,2.5,3.5,4.5,5.5,6.5}
    \fill[white] (\x,1) circle (2pt);
  
  % Linee verticali (griglia)
  \foreach \x in {1.5,2.5,3.5,4.5,5.5,6.5}
    \draw[white!40, dashed] (\x,1) -- (\x,3.5);
  
  % Soluzione numerica approssimata (spezzata)
  \draw[white, thick]
    (1.5,1.8) --
    (2.5,2.6) --
    (3.5,2.1) --
    (4.5,3.2) --
    (5.5,2.4) --
    (6.5,3.0);
  
  % Punti della soluzione
  \foreach \p in {(1.5,1.8),(2.5,2.6),(3.5,2.1),(4.5,3.2),(5.5,2.4),(6.5,3.0)}
    \fill[white] \p circle (2pt);
  
  % Equazione numerica
  \node[white, font=\large] at (6,4.2) {$x_{k+1}=x_k-\frac{f(x_k)}{f'(x_k)}$};
  
  % Notazione del dominio
  \node[white, font=\tiny] at (1,1.2) {$\Omega$};
  
\end{tikzpicture}

    \vspace{1.2cm}
    {\fontsize{32}{38}\selectfont\bfseries Metodi Numerici}\\[0.3cm]
    {\Large\itshape Parte 1}\\[2cm]

    \begin{tikzpicture}
      \draw[white!40, line width=0.5pt] (0,0) -- (10,0);
    \end{tikzpicture}\\[0.5cm]

    {\normalsize\itshape
      ``The purpose of computing is insight, not numbers.''}\\[0.4cm]
    {\small --- Richard Hamming, 1962}
  \end{center}
  \vfill
  
  \begin{center}
    {\large Corso di Laurea in Scienze Matematiche per l'Intelligenza Artificiale}\\
    {\large Anno Accademico 2024--2025}
  \end{center}
\end{titlepage}
\nopagecolor
\clearpage
\tableofcontents

\newpage
\section*{Richiami di Algebra Lineare}
\addcontentsline{toc}{section}{Richiami di Algebra Lineare}

\subsection*{Numeri complessi}
\addcontentsline{toc}{subsection}{Numeri complessi}
Un numero complesso \( z \in \mathbb{C} \) è definito come \( z = a + ib \), con \( a, b \in \mathbb{R} \), \( i^2 = -1 \).

\begin{itemize}
  \item Coniugato: \( \overline{z} = a - ib \)
  \item Modulo: \( |z| = \sqrt{a^2 + b^2} \)
  \item Prodotto scalare tra complessi: \( (z_1, z_2) = \overline{z_1} z_2 \in\mathbb{R}^+\)
  \item Forma esponenziale: \( z = re^{i\theta} = r(\cos\theta + i\sin\theta) \)
\end{itemize}

\subsection*{Vettori e spazi vettoriali}
\addcontentsline{toc}{subsection}{Vettori e spazi vettoriali}

Un \textbf{vettore} in \( \mathbb{R}^n \) o \( \mathbb{C}^n \) è una tupla ordinata di \( n \) numeri.

Uno \textbf{spazio vettoriale} \( V \) su un campo \( \mathbb{K} = \mathbb{R} \text{ o } \mathbb{C} \) è un insieme chiuso rispetto a due operazioni:  
la somma vettoriale e la moltiplicazione per uno scalare in \( \mathbb{K} \).

\vspace{0.3cm}

\textbf{Base}: un insieme di vettori \( \mathcal{B} = \{v_1, \dots, v_k\} \subseteq V \) tale che:  
1. i vettori sono linearmente indipendenti;  
2. generano tutto lo spazio, cioè ogni \( v \in V \) si scrive come  
\[
v = \sum_{i=1}^k \alpha_i v_i, \quad \alpha_i \in \mathbb{K}.
\]

\vspace{0.3cm}

\textbf{Span}: dato un insieme di vettori \( \{v_1, \dots, v_k\} \), lo span è  
\[
\text{span}(v_1, \dots, v_k) = \left\{ \sum_{i=1}^k \alpha_i v_i \mid \alpha_i \in \mathbb{K} \right\}.
\]

\vspace{0.3cm}

\textbf{Dimensione}: il numero di vettori in una base di \( V \), indicato con  
\[
\dim(V) = k.
\]

\subsection*{Matrici e trasformazioni lineari}
\addcontentsline{toc}{subsection}{Matrici e trasformazioni lineari}

Una \textbf{matrice} \( A \in \mathbb{K}^{n \times m} \) rappresenta una trasformazione lineare  
\[
T: \mathbb{K}^m \to \mathbb{K}^n, \quad T(x) = Ax,
\]
dove \( x \in \mathbb{K}^m \) è un vettore colonna.

\vspace{0.3cm}

\textbf{Prodotto di matrici}:  
Se \( A \in \mathbb{K}^{n \times m} \) e \( B \in \mathbb{K}^{m \times p} \), allora  
\[
C = AB \in \mathbb{K}^{n \times p}, \quad C_{ij} = \sum_{k=1}^m A_{ik} B_{kj}.
\]

\vspace{0.3cm}

\textbf{Nucleo} (kernel) di $\underbrace{T}_{A}$:  
\[
\underbrace{\ker(T)}_{ker(A)} = \{ x \in \mathbb{K}^m \mid \underbrace{T(x)}_{Ax} = 0 \}.
\]

\vspace{0.3cm}

\textbf{Immagine} (range) di $\underbrace{T}_{A}$:  
\[
\underbrace{\operatorname{Im}(T)}_{\operatorname{Im}(A)} = \{ \underbrace{T(x)}_{Ax} \mid x \in \mathbb{K}^m \} \subseteq \mathbb{K}^n.
\]


\vspace{0.3cm}

\textbf{Rango} di $\underbrace{T}_{A}$:  
\[
\underbrace{\operatorname{rank}(T)}_{\operatorname{rank}(A)} = \dim \left( \underbrace{\operatorname{Im}(T)}_{\operatorname{Im}(A)} \right)
\]
\text{O equivalentemente data A}$\in\mathbb{K}^{n\times m}$
\[
\operatorname{rank}(A) = \max \left\{ k \in \mathbb{N} \mid \exists S \subseteq \{1, \ldots, m\}, |S| = k, \left( \sum_{j \in S} \alpha_j A_j = 0 \implies \forall j, \alpha_j = 0 \right) \right\}
\]

\vspace{0.3cm}

\textbf{Teorema fondamentale della dimensione}:  
Per \( A \in \mathbb{K}^{n \times m} \) vale  
\[
\operatorname{rank}(A) + \dim(\ker(A)) = m,
\]
dove \( m \) è il numero di colonne.
\paragraph{} Siano \(U, V\) sottospazi vettoriali di uno spazio vettoriale \(W\). Allora vale la formula di Grassmann:
\[
\dim(U + V) = \dim(U) + \dim(V) - \dim(U \cap V).
\]


\subsection*{Riduzione di Gauss e rango}
La \textbf{riduzione di Gauss} trasforma \( A \) in forma a scala tramite operazioni elementari sulle righe, che preservano il rango.

Una matrice si dice \textbf{ridotta a scala} se:  
- il primo elemento non nullo (pivot) di ogni riga è a destra di quello della riga precedente;  
- le righe nulle sono in fondo.

Le operazioni elementari sono:  
1. scambio di righe;  
2. moltiplicazione di una riga per uno scalare non nullo;  
3. somma di un multiplo di una riga a un'altra.

Ogni operazione corrisponde a moltiplicare \( A \) a sinistra per una matrice elementare invertibile \( E_i \).

Così si ha
\[
E_k \cdots E_1 A = R,
\]
dove \( R \) è in forma a scala e
\[
\operatorname{rank}(A) = \operatorname{rank}(R) = \text{numero di pivot in } R.
\]

\subsubsection*{Calcolo della matrice inversa}
Per \( A \in \mathbb{R}^{n \times n} \), si affianca \( I_n \) e si applicano pivotaggi:
\[
\left[
\begin{array}{c|c}
A & I_n
\end{array}
\right]
\quad \xrightarrow{\text{pivotaggi}} \quad
\left[
\begin{array}{c|c}
I_n & B
\end{array}
\right].
\]

Allora i passaggi che abbiamo fatto sono stati equivalenti a moltiplicare entrambe le matrici $I_n$ e $A$ per \( B = A^{-1} \) poiché
\[
AB = BA = I_n.
\]

\subsection*{Determinante}
\addcontentsline{toc}{subsection}{Determinante}
Definito solo per matrici n$\times$ n. Proprietà principali:
\begin{itemize}
  \item \( \det(AB) = \det(A)\det(B) \)
  \item \( \det(A^T) = \det(A) \)
  \item \( \det(A) = 0 \iff A \) non è invertibile ovvero $\not\exists B \in \mathbb{R}^{n\times n} \text{ tale che }  AB=BA=I_n$
\end{itemize}
Se $\det(
A) \ne 0$ con $A=\begin{pmatrix}
\vdots \\
a_{i1} \cdots a_{in} \\
\vdots \\
a_{j1} \cdots a_{jn} \\
\vdots
\end{pmatrix} = 
\begin{pmatrix}
\vdots \\
R_i \\
\vdots \\
R_j \\
\vdots
\end{pmatrix}, \quad R_i = \begin{bmatrix}
a_{i1} & \cdots & a_{in}
\end{bmatrix}\Rightarrow$
\[
\begin{cases}
    \exists B \in \mathbb{R}^{n \times n} \text{ tale che } AB = BA = I_n \quad (\text{cioè } B = A^{-1}) \\
    \sum_{i=1}^n \alpha_i A_i = 0 \Rightarrow \alpha_1 = \dots = \alpha_n = 0 \quad \text{(colonne linearmente indipendenti)}
 \\
     \sum_{i=1}^n \alpha_i (A^\top)_i = 0 \Rightarrow \alpha_1 = \dots = \alpha_n = 0 \quad \text{(righe linearmente indipendenti)}
 \\
    Ax = 0 \Rightarrow x = \vec{0} \text{ (A ha ker banale cioè ker(A)}\textbf{=}\text{\{$\vec{0}$\})}\\
    \operatorname{rank}(A) = n \\
    Ax = \lambda x \Rightarrow \lambda \ne 0 \\
    \forall y \in \mathbb{R}^n \quad \exists ! x \in \mathbb{R}^n \text{ t.c. } Ax = y \text{ (Im(A)}\textbf{ = }\text{$\mathbb{R}^n)$}\\

A' = 
\begin{pmatrix}
\vdots \\
a_{j1} \cdots a_{jn} \\
\vdots \\
a_{i1} \cdots a_{in} \\
\vdots
\end{pmatrix}
\quad
\Rightarrow
\det(A') = -\det(A)\\

A' = 
\begin{pmatrix}
\vdots \\
\lambda a_{i1} \cdots \lambda a_{in} \\
\vdots
\end{pmatrix}

\Rightarrow
\det(A') = \lambda \det(A)
\\R_i = R_i^{(1)} + R_i^{(2)}
\Rightarrow
\det\begin{pmatrix}
\vdots \\
R_i \\
\vdots
\end{pmatrix}
=
\det\begin{pmatrix}
\vdots \\
R_i^{(1)} \\
\vdots
\end{pmatrix}
+
\det\begin{pmatrix}
\vdots \\
R_i^{(2)} \\
\vdots
\end{pmatrix}\\
A' = 
\begin{pmatrix}
\vdots \\
R_i + \lambda R_j \\
\vdots \\
R_j \\
\vdots
\end{pmatrix}
\Rightarrow
\det(A') = \det(A)\\
A = (a_{ij}) \quad \text{con} \quad a_{ij} = 0 \quad \forall i > j
\Rightarrow
\det(A) = \prod_{i=1}^n a_{ii}
\end{cases}
\]



\subsection*{Cambio di base e matrici simili}
\addcontentsline{toc}{subsection}{Cambio di base e matrici simili}
Sia \( V \) uno spazio vettoriale di dimensione \( n \), con due basi
\[
E_1 = \{ w_1, \ldots, w_n \}, \quad E = \{ v_1, \ldots, v_n \}
\]
e sia \( T: V \to V \) un endomorfismo.

\vspace{0.3cm}
\textbf{Definizione di \( P \):} la matrice di cambio base \( P \) trasforma coordinate da base \( E \) a base \( E_1 \), cioè
\[
[v]_{E_1} = P [v]_{E} \quad \forall v \in V.
\]

\vspace{0.3cm}
\textbf{Obiettivo:} collegare le matrici di \( T \) in basi diverse:
\[
[T]_{E_1} \quad \text{e} \quad [T]_{E}.
\]

\vspace{0.3cm}
\[
\begin{cases}
x = [v]_{E} \\[6pt]
P x = [v]_{E_1} \\[6pt]
[T]_{E_1} (P x) = [T(v)]_{E_1} \\[6pt]
P^{-1} [T]_{E_1} (P x) = [T(v)]_{E} \\[6pt]
\end{cases}
\Rightarrow
[T]_{E} x = P^{-1} [T]_{E_1} P x.
\]

Essendo vero per ogni \( x \), si conclude che:
\[
\boxed{
[T]_{E} = P^{-1} [T]_{E_1} P.
}
\]

Quindi il cambio di base è semplicemente un “passaggio di coordinate” tramite \( P \) e \( P^{-1} \), che cambia la rappresentazione di \( T \) senza alterarne l’azione su \( V \).

Due matrici \( A, B \in M_n(\mathbb{K}) \) si dicono \emph{simili} se esiste una matrice invertibile \( P \in M_n(\mathbb{K}) \) tale che
\[
B = P^{-1} A P.
\]

\vspace{1em}
\hline
\vspace{1em}

Ora mi chiedo: chi è la matrice \( P \) di passaggio dalla base \(\bm{E}\) alla base \(\bm{E_1}\)?
Denotiamo questa matrice con
\[
P^{\bm{E} \to \bm{E_1}}.
\]

Questa è una matrice quadrata \( n \times n \), quindi dobbiamo esprimere i vettori nuovi \( w_1, w_2, \ldots, w_n \) come combinazioni lineari di quelli della vecchia base \( v_1, v_2, \ldots, v_n \):
\[
\begin{cases}
w_1 = a_{11} v_1 + a_{21} v_2 + \cdots + a_{n1} v_n \\
w_2 = a_{12} v_1 + a_{22} v_2 + \cdots + a_{n2} v_n \\
\vdots \\
w_n = a_{1n} v_1 + a_{2n} v_2 + \cdots + a_{nn} v_n
\end{cases}
\]

Quindi la matrice \( P \) è la matrice formata dai coefficienti \( (a_{ij}) \), cioè
\[
P^{\bm{E} \to \bm{E_1}} = (a_{ij})_{i,j=1}^n.
\]

Questi coefficienti si trovano risolvendo \( n \) sistemi lineari, uno per ogni vettore \( w_j \), esprimendo \( w_j \) come combinazione lineare dei vettori \( v_i \).

\vspace{0.3cm}
\textbf{Esempio numerico:}

Consideriamo \( V = \mathbb{R}^2 \) con basi  
\[
\bm{E} = \left\{ v_1 = \begin{pmatrix} 1 \\ 0 \end{pmatrix}, \quad v_2 = \begin{pmatrix} 0 \\ 1 \end{pmatrix} \right\}
,\quad\bm{E_1} = \left\{ w_1 = \begin{pmatrix} 1 \\ 1 \end{pmatrix}, \quad w_2 = \begin{pmatrix} -1 \\ 2 \end{pmatrix} \right\}.
\]

Per trovare i coefficienti \( a_{ij} \) per il primo vettore \( w_1 \), risolviamo il sistema:
\[
w_1 = a_{11} v_1 + a_{21} v_2 \quad \Longrightarrow \quad
\begin{pmatrix} 1 \\ 1 \end{pmatrix} = a_{11} \begin{pmatrix} 1 \\ 0 \end{pmatrix} + a_{21} \begin{pmatrix} 0 \\ 1 \end{pmatrix} = \begin{pmatrix} a_{11} \\ a_{21} \end{pmatrix},
\]
che dà  
\[
a_{11} = 1, \quad a_{21} = 1.
\]

Per il secondo vettore \( w_2 \):
\[
w_2 = a_{12} v_1 + a_{22} v_2 \quad \Longrightarrow \quad
\begin{pmatrix} -1 \\ 2 \end{pmatrix} = \begin{pmatrix} a_{12} \\ a_{22} \end{pmatrix},
\]
quindi
\[
a_{12} = -1, \quad a_{22} = 2.
\]

Da cui la matrice di passaggio è
\[
P^{\bm{E} \to \bm{E_1}} = \begin{pmatrix}
1 & -1 \\
1 & 2
\end{pmatrix}.
\]

\vspace{0.3cm}
\textbf{Interpretazione:} la matrice \( P^{\bm{E} \to \bm{E_1}} \) rappresenta la matrice associata alla trasformazione identica
\[
\mathrm{id} : V \to V
\]
\textbf{da base \(\bm{E_1}\) a base \(\bm{E}\)}, cioè  
\[
P^{\bm{E} \to \bm{E_1}} = M^{\bm{E_1}, \bm{E}}(\mathrm{id}).
\]

Questo perché, dato un vettore \( v \in V \), le sue coordinate in \(\bm{E_1}\) si trasformano in coordinate in \(\bm{E}\) tramite
\[
[v]_{\bm{E}} = P^{\bm{E} \to \bm{E_1}} [v]_{\bm{E_1}}.
\]
Ricordiamo le basi di \( V = \mathbb{R}^2 \):
\[
\bm{E} = \left\{ v_1 = \begin{pmatrix} 1 \\ 0 \end{pmatrix}, \quad v_2 = \begin{pmatrix} 0 \\ 1 \end{pmatrix} \right\}, \quad
\bm{E_1} = \left\{ w_1 = \begin{pmatrix} 1 \\ 1 \end{pmatrix}, \quad w_2 = \begin{pmatrix} -1 \\ 2 \end{pmatrix} \right\}.
\]

La matrice di passaggio calcolata è
\[
P^{\bm{E} \to \bm{E_1}} = \begin{pmatrix} 1 & -1 \\ 1 & 2 \end{pmatrix}.
\]

\vspace{0.3cm}
\textbf{Obiettivo:} dimostrare che
\[
P^{\bm{E} \to \bm{E_1}} = M^{\bm{E_1}, \bm{E}}(\mathrm{id}),
\]
cioè la matrice associata alla trasformazione identica
\[
\mathrm{id} : V \to V,
\]
considerata come applicazione da base \(\bm{E_1}\) a base \(\bm{E}\).

\vspace{0.3cm}
Per definizione, la matrice \( M^{\bm{E_1}, \bm{E}}(\mathrm{id}) \) si ottiene esprimendo ogni vettore della base \(\bm{E_1}\) come combinazione lineare dei vettori della base \(\bm{E}\).

Per \( w_1 \), riscriviamo:
\[
w_1 = \begin{pmatrix} 1 \\ 1 \end{pmatrix} = 1 \cdot v_1 + 1 \cdot v_2,
\]
quindi la colonna corrispondente è
\[
\begin{pmatrix} 1 \\ 1 \end{pmatrix}.
\]

Per \( w_2 \), abbiamo:
\[
w_2 = \begin{pmatrix} -1 \\ 2 \end{pmatrix} = (-1) \cdot v_1 + 2 \cdot v_2,
\]
quindi la seconda colonna è
\[
\begin{pmatrix} -1 \\ 2 \end{pmatrix}.
\]

\vspace{0.3cm}
Assemblando le colonne:
\[
M^{\bm{E_1}, \bm{E}}(\mathrm{id}) = \begin{pmatrix} 1 & -1 \\ 1 & 2 \end{pmatrix} = P^{\bm{E} \to \bm{E_1}}.
\]

\vspace{0.3cm}
\
P^{\bm{E} \to \bm{E_1}} = M^{\bm{E_1}, \bm{E}}(\mathrm{id}),

\]
cioè la matrice di passaggio dalla base \(\bm{E_1}\) alla base \(\bm{E}\) è proprio la matrice associata alla trasformazione identica in quelle basi
\vspace{1em}
\hline
\vspace{1em}
Sia \( V \) uno spazio vettoriale con \(\dim(V) = n\), consideriamo un endomorfismo
\[
\phi : V \to V.
\]

Siano \( \bm{E} \) e \( \bm{E_1} \) due basi di \( V \), con
\[
\bm{E} = \{ v_1, \ldots, v_n \}, \quad \bm{E_1} = \{ w_1, \ldots, w_n \}.
\]

Definiamo la matrice di passaggio
\[
P^{\bm{E} \to \bm{E_1}} = M_1^{\bm{E_1} \to \bm{E}}(\mathrm{id}_V),
\]
cioè la matrice associata alla trasformazione identica da base \(\bm{E_1}\) a base \(\bm{E}\).

Analogamente, per lo spazio \( W = V \) (essendo endomorfismo), abbiamo
\[
Q^{\bm{E} \to \bm{E_1}} = M_2^{\bm{E_1} \to \bm{E}}(\mathrm{id}_V).
\]

Poiché \( V = W \) e \( \bm{F} = \bm{E} \), \( \bm{F_1} = \bm{E_1} \), si ha che
\[
P = Q,
\]
cioè la matrice di cambio base nel dominio è la stessa del codominio.

Sia
\[
A = M^{\bm{E} \to \bm{E}}(\phi)
\]
la matrice di \(\phi\) rispetto alla base \(\bm{E}\).

Se vogliamo trovare la matrice di \(\phi\) rispetto alla base \(\bm{E_1}\), cioè
\[
B = M^{\bm{E_1} \to \bm{E_1}}(\phi),
\]
allora vale la formula del cambio di base:
\[
\boxed{
B = Q^{-1} A P.
}
\]

Ma, come detto, \( P = Q \), quindi
\[
\boxed{
B = P^{-1} A P.
}
\]

\subsubsection*{Esercizio: cambio di base per \( T : V \to W \)}
\addcontentsline{toc}{subsubsection}{Esercizio: cambio di base per \( T : V \to W \)}

Siano
\[
V = \mathbb{R}^2, \quad W = \mathbb{R}^3,
\quad
T : V \to W, \quad
T\left(\begin{pmatrix} x \\ y \end{pmatrix}\right) = \begin{pmatrix} x + y \\ 2x - y \\ x \end{pmatrix}.
\]

Basi:
\[
\bm{B}_V = \left\{ \begin{pmatrix} 1 \\ 0 \end{pmatrix}, \begin{pmatrix} 0 \\ 1 \end{pmatrix} \right\}, \quad
\bm{B}_V' = \left\{ \begin{pmatrix} 1 \\ 1 \end{pmatrix}, \begin{pmatrix} 1 \\ -1 \end{pmatrix} \right\}
\]
\[
\bm{B}_W = \left\{ \begin{pmatrix} 1 \\ 0 \\ 0 \end{pmatrix}, \begin{pmatrix} 0 \\ 1 \\ 0 \end{pmatrix}, \begin{pmatrix} 0 \\ 0 \\ 1 \end{pmatrix} \right\}, \quad
\bm{B}_W' = \left\{ \begin{pmatrix} 1 \\ 1 \\ 0 \end{pmatrix}, \begin{pmatrix} 0 \\ 1 \\ 1 \end{pmatrix}, \begin{pmatrix} 1 \\ 0 \\ 1 \end{pmatrix} \right\}
\]

1. Trova \( [T]_{\bm{B}_W, \bm{B}_V} \)

Applichiamo \( T \) ai vettori della base \( \bm{B}_V \):
\[
T\left( \begin{pmatrix} 1 \\ 0 \end{pmatrix} \right) = \begin{pmatrix} 1 \\ 2 \\ 1 \end{pmatrix}, \quad
T\left( \begin{pmatrix} 0 \\ 1 \end{pmatrix} \right) = \begin{pmatrix} 1 \\ -1 \\ 0 \end{pmatrix}
\]

Quindi
\[
[T]_{\bm{B}_W, \bm{B}_V} =
\begin{pmatrix}
1 & 1 \\
2 & -1 \\
1 & 0
\end{pmatrix}
\]

2. Trova \( [T]_{\bm{B}_W', \bm{B}_V'} \)

Serve:
\[
[T]_{\bm{B}_W', \bm{B}_V'} = P_W^{-1} [T]_{\bm{B}_W, \bm{B}_V} P_V
\]

Matrice di cambio base \( P_V \): colonne = \( \bm{B}_V' \) scritti in \( \bm{B}_V \)
\[
P_V = \begin{pmatrix}
1 & 1 \\
1 & -1
\end{pmatrix}
\]

Matrice \( P_W \): colonne = \( \bm{B}_W' \) scritti in \( \bm{B}_W \)
\[
P_W =
\begin{pmatrix}
1 & 0 & 1 \\
1 & 1 & 0 \\
0 & 1 & 1
\end{pmatrix}
\Rightarrow
P_W^{-1} = 
\begin{pmatrix}
1 & 0 & -1 \\
-1 & 1 & 1 \\
1 & -1 & 0
\end{pmatrix}
\]

Ora calcolo:
\[
[T]_{\bm{B}_W', \bm{B}_V'} =
\begin{pmatrix}
1 & 0 & -1 \\
-1 & 1 & 1 \\
1 & -1 & 0
\end{pmatrix}
\begin{pmatrix}
1 & 1 \\
2 & -1 \\
1 & 0
\end{pmatrix}
\begin{pmatrix}
1 & 1 \\
1 & -1
\end{pmatrix}
=
\boxed{
\begin{pmatrix}
0 & 1 \\
4 & -2 \\
-2 & 0
\end{pmatrix}
}
\]
\subsection*{Autovalori e autovettori}
\addcontentsline{toc}{subsection}{Autovalori e autovettori}
Dato \( A \in \mathbb{R}^{n \times n} \), \( \lambda \in \mathbb{C} \) è autovalore se \( \exists x \ne 0 \) tale che:
\[ Ax = \lambda x \iff (A - \lambda I)x = 0 \]
Gli autovalori \( \lambda_i \) sono radici del polinomio caratteristico \( \det(A - \lambda I) = 0 \)
\text{, i corrispondenti autovettori} si trovano risolvendo \( (A - \lambda_i I)x = 0 \) per x

\subsection*{Matrici diagonalizzabili}
\addcontentsline{toc}{subsection}{Matrici diagonalizzabili}


\[
A \in \mathbb{R}^{n \times n} \text{ è diagonalizzabile} \iff \text{gli autovettori di } A \text{ formano una base di } \mathbb{R}^n
\]

Questo è equivalente a dire che per ogni autovalore \( \lambda \) di \( A \), la \textbf{molteplicità geometrica} coincide con la \textbf{molteplicità algebrica}:

\[
\forall \lambda \in \sigma(A), \quad \dim(E_\lambda) = \text{molteplicità algebrica di } \lambda
\]

Dove \( E_\lambda = \ker(A - \lambda I) \) è lo \textbf{spazio degli autovettori} relativi a \( \lambda \), anche detto \textbf{autospazio}.

\vspace{1em}

\textbf{Molteplicità algebrica} di un autovalore \( \lambda \):

\[
\text{MA}(\lambda) := \text{molteplicità della radice } \lambda \text{ del polinomio caratteristico } \chi_A(\lambda) = \det(A - \lambda I)
\]
La molteplicità di una radice a di un polinomio p(x)
 è il massimo numero k tale che $(x-a)^k$ divide p(x).

\textbf{Molteplicità geometrica} di un autovalore \( \lambda \):

\[
\text{MG}(\lambda) := \dim(\ker(A - \lambda I)) = n - \operatorname{rank}(A - \lambda I)
\]

\vspace{1em}

\textbf{Osservazioni}:
\begin{itemize}
  \item Sempre: \( 1 \le \text{MG}(\lambda) \le \text{MA}(\lambda) \)
  \item \( A \) è diagonalizzabile \( \iff \text{MG}(\lambda) = \text{MA}(\lambda) \) per ogni \( \lambda \)
  \item In tal caso, \( A = PDP^{-1} \), con \( D \) diagonale e \( P \) contenente una base di autovettori di \( A \)
\end{itemize}


\subsection*{Matrici simmetriche e teorema spettrale}
\addcontentsline{toc}{subsection}{Matrici simmetriche e teorema spettrale}

Sia \( A \in \mathbb{R}^{n \times n} \) simmetrica, cioè \( A = A^T \). Allora:

\begin{itemize}
  \item Gli autovalori di \( A \) sono reali
  \item Autovettori associati ad autovalori distinti sono ortogonali
  \item \( A \) è diagonalizzabile con una matrice ortogonale: \( A = Q D Q^T \)
\end{itemize}

\vspace{0.5em}
\textbf{Dimostrazione (1): gli autovalori sono reali}

Sia \( Ax = \lambda x \), con \( x \ne 0 \). Allora:
\[
x^T A x = \lambda x^T x \qquad (\text{perché } Ax = \lambda x)
\]
\[
x^T A x = (Ax)^T x = (\lambda x)^T x = \overline{\lambda} x^T x
\Rightarrow \lambda x^T x = \overline{\lambda} x^T x
\Rightarrow \lambda = \overline{\lambda} \in \mathbb{R}
\]

\vspace{0.5em}
\textbf{Dimostrazione (2): ortogonalità di autovettori associati a autovalori distinti}

Siano \( Ax = \lambda x \), \( Ay = \mu y \), con \( \lambda \ne \mu \).

\[
x^T A y = \mu x^T y \qquad (\text{perché } Ay = \mu y)
\]
\[
x^T A y = (A x)^T y = (\lambda x)^T y = \lambda x^T y
\Rightarrow \lambda x^T y = \mu x^T y
\Rightarrow (\lambda - \mu) x^T y = 0
\Rightarrow x^T y = 0
\Rightarrow x \perp y
\]

\vspace{0.5em}
\textbf{Dimostrazione (3): esistenza di una base ortonormale di autovettori}

Sia \( A \in \mathbb{R}^{n \times n} \) simmetrica. Allora:

\begin{itemize}
  \item \( A \) ha \( n \) autovalori reali (contati con molteplicità algebrica) per il teorema fondamentale dell'algebra: il polinomio caratteristico \( \det(A - \lambda I) \in \mathbb{R}[\lambda] \) è di grado $n$ quindi ha \( n \) radici (che rappresentano gli autovalori).
  
  \item Per ogni autovalore \( \lambda \), l'insieme:
  \[
  E_\lambda := \ker(A - \lambda I) = \{x \in \mathbb{R}^n \mid Ax = \lambda x\}
  \]
  è chiamato \textbf{autospazio} associato a \( \lambda \). È un sottospazio vettoriale di \( \mathbb{R}^n \).

  \item Ogni autospazio è \textbf{\( A \)-invariante}: se \( x \in E_\lambda \), allora \( Ax = \lambda x \in E_\lambda \), cioè \( A \) manda vettori di \( E_\lambda \) dentro \( E_\lambda \).

  \item Ogni sottospazio vettoriale di \( \mathbb{R}^n \) ammette una \textbf{base ortonormale} (per il prodotto scalare standard) per il teorema di Gram-Schmidt.

  \item Inoltre, dalla (2), autospazi associati ad autovalori distinti sono ortogonali: se \( x \in E_\lambda \), \( y \in E_\mu \), con \( \lambda \ne \mu \), allora \( x \perp y \).
\end{itemize}

\noindent Quindi possiamo prendere basi ortonormali \( \mathcal{B}_1, \dots, \mathcal{B}_k \) per gli autospazi distinti \( E_{\lambda_1}, \dots, E_{\lambda_k} \), e metterle insieme per ottenere una base ortonormale \( \mathcal{B} \) di \( \mathbb{R}^n \) formata tutta da autovettori.

Allora, detta \( Q \) la matrice che ha per colonne i vettori di \( \mathcal{B} \), abbiamo:

\[
Q \text{ ortogonale}, \quad D = \operatorname{diag}(\lambda_1, \dots, \lambda_n), \quad A = Q D Q^T
\]

\hfill $\square$













\newpage

\section*{Richiami di Analisi 1}
\addcontentsline{toc}{section}{Richiami di Analisi 1}

\subsection*{Intorni, Punti di accumulazione}
\addcontentsline{toc}{subsection}{Intorni, Punti di accumulazione}
\begin{itemize}
    \item \textbf{Intorno:}  
    Sia \( x_0 \in \underbrace{\bar{\mathbb{R}}}_{\mathbb{R}\cup\{-\infty,+\infty\}}
    }} \). Un insieme \( V \subseteq \bar{\mathbb{R}} \) è un \emph{intorno} di \( x_0 \) se  
    \[
    \exists \delta > 0 : (x_0 - \delta, x_0 + \delta) \subseteq V.
    \]
    
    \item \textbf{Punto di accumulazione:}  
    Un punto \( x_0 \in \bar{\mathbb{R}} \) è un \emph{punto di accumulazione} di un insieme \( A \subseteq \bar{\mathbb{R}} \) se  
    \[
    \forall V \text{ intorno di } x_0, \quad (V \setminus \{x_0\}) \cap A \neq \emptyset.
    \]
    (cioè ogni intorno di \( x_0 \) contiene punti di \( A \) diversi da \( x_0 \)).
\end{itemize}

\subsection*{Funzioni reali di variabile reale}
\addcontentsline{toc}{subsection}{Funzioni reali di variabile reale}

Siano \( A, B \subseteq \mathbb{R} \) due sottoinsiemi dei numeri reali.  
Una \textbf{funzione reale di variabile reale} è un'applicazione
\[
f : A \to B, \quad x \mapsto f(x),
\]
che associa a ogni elemento \( x \in A \) un unico elemento \( f(x) \in B \).

\begin{itemize}
    \item L'insieme \( A \) è detto \textbf{dominio} della funzione.
    \item L'insieme \( B \) è detto \textbf{codominio} della funzione.
    \item L'insieme \(\operatorname{Im}(f) = \{ f(x) \mid x \in A \} \subseteq B\) è detto \textbf{immagine} della funzione.
\end{itemize}

\paragraph{Iniettività:}  
Una funzione \( f : A \to B \) è detta \textbf{iniettiva} se
\[
f(x_1) = f(x_2) \Rightarrow x_1 = x_2 \quad \text{per ogni } x_1, x_2 \in A.
\]
In altre parole, valori distinti nel dominio corrispondono a valori distinti nel codominio.

\paragraph{Suriettività:}  
La funzione \( f : A \to B \) è \textbf{suriettiva} se
\[
\forall y \in B,\ \exists x \in A \text{ tale che } f(x) = y,
\]
cioè se l'immagine della funzione coincide con il codominio: \( \operatorname{Im}(f) = B \).

\paragraph{Biettività:}  
La funzione \( f \) è \textbf{biettiva} se è sia iniettiva che suriettiva.  
In tal caso, ogni elemento del codominio è immagine di un unico elemento del dominio.

\paragraph{Funzione inversa:}  
Se \( f : A \to B \) è biettiva, allora esiste una funzione inversa
\[
f^{-1} : B \to A, \quad f^{-1}(y) = x \quad \text{tale che } f(x) = y.
\]
La funzione inversa soddisfa:
\[
f^{-1}(f(x)) = x \quad \text{per ogni } x \in A, \qquad f(f^{-1}(y)) = y \quad \text{per ogni } y \in B.
\]

\paragraph{Composizione di funzioni:}  
Siano \( f : A \to B \) e \( g : B \to C \) due funzioni.  
La \textbf{composizione} di \( g \) con \( f \), denotata con \( g \circ f \), è la funzione
\[
g \circ f : A \to C, \quad (g \circ f)(x) = g(f(x)).
\]
La composizione è \textbf{associativa}:  
\[
h \circ (g \circ f) = (h \circ g) \circ f \quad \text{se le composizioni sono definite.}
\]
\subsection*{Limiti di funzioni}
\addcontentsline{toc}{subsection}{Limiti di funzioni}
Sia \( f : A \to \bar{\mathbb{R}} \), \( x_0 \) punto di accumulazione di \( A \), \( L \in \bar{\mathbb{R}}\)

\paragraph{Definizione classica:}
\[
\lim_{x \to x_0} f(x) = L \quad \Longleftrightarrow \quad \forall \varepsilon > 0, \exists \delta > 0 :
\]
\[
x \in A, 0 < |x - x_0| < \delta \implies |f(x) - L| < \varepsilon.
\]
\[\text{o equivalentemente }\Longleftrightarrow \quad \forall V_L \text{ intorno di } L,
\]
\[
\exists V_{x_0} \text{ intorno di } x_0 : (V_{x_0} \setminus \{x_0\}) \cap A \implies f(x) \in V_L.
\]

\paragraph{Limiti sinistro e destro:}
\[
\lim_{x \to x_0^-} f(x) = L \quad \Longleftrightarrow \quad x \to x_0, x < x_0,
\]
\[
\lim_{x \to x_0^+} f(x) = L \quad \Longleftrightarrow \quad x \to x_0, x > x_0.
\]

\paragraph{Funzioni asintotiche}
Date due funzioni $f,g$ reali di variabile reale ed $x_0$ p.d.a. per l'intersezione dei domini di $f$ e $g$, diciamo che $f$ è asintotica a $g$ per x$\to x_0$ e scriviamo:
\[f\sim g\iff\lim_{x \to x_0} \frac{f(x)}{g(x)} = 1\]
\subsection*{Asintoti}
\addcontentsline{toc}{subsection}{Asintoti}
\begin{itemize}
    \item \textbf{Asintoto verticale} in \( x_0 \) se  
    \[
    \lim_{x \to x_0^{\pm}} f(x) = \pm \infty.
    \]
    \item \textbf{Asintoto orizzontale} in \( y = l \) se  
    \[
    \lim_{x \to \pm \infty} f(x) = l.
    \]
    \item \textbf{Asintoto obliquo} se esistono \( m, q \in \mathbb{R} \) tali che  
    \[
    \lim_{x \to \pm \infty} \bigl[f(x) - (mx + q)\bigr] = 0 \quad(m = \lim_{x \to \pm \infty} \frac{f(x)}{x},q = \lim_{x \to \pm \infty} \bigl[f(x) - mx \bigr])
    \]
\end{itemize}

\subsection*{Continuità, continuità uniforme, Lipschitzianità e derivabilità}
\addcontentsline{toc}{subsection}{Continuità, continuità uniforme, Lipschitzianità e derivabilità}

\begin{itemize}
    \item \( f \) continua in \( x_0 \) se  
    \[
    \lim_{x \to x_0} f(x) = f(x_0).
    \]
    \item \( f \) è \textbf{continua uniformemente} su \( A \) se  
    \[
    \forall \varepsilon > 0, \exists \delta > 0 : \forall x,y \in A,\, |x - y| < \delta \implies |f(x) - f(y)| < \varepsilon.
    \]
    \item \( f \) è \textbf{Lipschitziana} su \( A \) se  
    \[
    \exists L > 0 : \forall x, y \in A,\, |f(x) - f(y)| \leq L |x - y|.
    \]
    \item \( f \) è \textbf{derivabile} in \( x_0 \) se  
    \[
    \lim_{h \to 0} \frac{f(x_0 + h) - f(x_0)}{h}\in\mathbb{R}
    \]
\end{itemize}

\subsection*{Limiti notevoli}
\addcontentsline{toc}{subsection}{Limiti notevoli}

\[
\lim_{x \to 0} \frac{\sin x}{x} = 1, \quad \lim_{x \to 0} \frac{1 - \cos x}{x^2} = \frac{1}{2},\quad
\lim_{x \to 0} (1 + x)^{\frac{1}{x}} = e, \quad \lim_{x \to 0} \frac{e^x - 1}{x} = 1,\quad\dots\]

\text{(Ricavabili facilmente tramite lo sviluppo di Taylor al primo ordine)}

\subsection*{Successioni}
\addcontentsline{toc}{subsection}{Successioni}

\begin{itemize}
    \item Una \textbf{successione} è una funzione \( f : \mathbb{N} \to \mathbb{R} \),  
    indicata con \( (a_n)_{n \in \mathbb{N}} \), cioè  
    \[
    n \mapsto a_n.
    \]

    \item Una \textbf{sottosuccessione} di \( (a_n) \) è una successione \( (\bar{a}_{n}) \) ottenuta applicando  
    una funzione strettamente crescente \( h: \mathbb{N}\to\mathbb{N} \) agli indici n, cioè  
    \[
  \bar{a}_n:=a_{h(n)}
    \]

    \item La successione \( (a_n) \) è \textbf{monotona crescente} se  
    \[
    a_n \leq a_{n+1}, \quad \forall n,
    \]
    ed è \textbf{monotona decrescente} se  
    \[
    a_n \geq a_{n+1}, \quad \forall n.
    \]
\end{itemize}

\subsection*{Teorema di Bolzano-Weierstrass}
\addcontentsline{toc}{subsection}{Teorema di Bolzano-Weierstrass}

\begin{itemize}
    \item Una successione \( (a_n) \subseteq \mathbb{R} \) è \emph{limitata} se  
    \[
    \exists M > 0 : |a_n| \leq M, \quad \forall n.
    \]

    \item \textbf{Teorema di Bolzano-Weierstrass:}  
    \[
    (x_n \text{ limitata}) \implies \exists \text{ una sottosuccessione } (\bar{a}_{n}) : \lim_{n\to \infty} \bar{a}_{n} = L \in \mathbb{R}.
    \]
\end{itemize}

\subsection*{Teorema di Weierstrass}
\addcontentsline{toc}{subsection}{Teorema di Weierstrass}
Se \( f : A\subseteq\mathbb{R} \to \mathbb{R} \) è continua e $\underbrace{\text{limitata in un intervallo }[a,b]}_{\exists M\in\mathbb{R} \text{ t.c. } f(x)\leq M \quad \forall x \in [a,b]}\subseteq A$, allora esistono
\[
x_{\min}, x_{\max} \in [a,b]
\]
tali che
\[
f(x_{\min}) \leq f(x) \leq f(x_{\max}), \quad \forall x \in [a,b].
\]

\subsection*{Teorema dei valori intermedi}
\addcontentsline{toc}{subsection}{Teorema dei valori intermedi}
Se \( f : [a,b] \to \mathbb{R} \) è continua e \( f(a) \neq f(b) \), allora  
\[
\forall y \in [\min(f(a),f(b)), \max(f(a),f(b))], \quad \exists c \in [a,b] : f(c) = y.
\]

\subsection*{Teorema degli zeri}
\addcontentsline{toc}{subsection}{Teorema degli zeri}
Se \( f : [a,b] \to \mathbb{R} \) è continua e
\[
f(a) \cdot f(b) < 0,
\]
allora
\[
\exists c \in (a,b) : f(c) = 0.
\]

\subsection*{Derivate}
\addcontentsline{toc}{subsection}{Derivate}
Data una funzione $f:A\subseteq\mathbb{R}\to\mathbb{R}$ definiamo la funzione derivata $f'$ di $f$ in $x_0$ (se esiste) come il limite del rapporto incrementale di $f$
    \[
    f'(x_0) = \lim_{h \to 0} \frac{f(x_0 + h) - f(x_0)}{h}.
    \]
    Regole principali:  
    \begin{itemize}
    \item$(f + g)' = f' + g', $
    \item$(f \cdot g)'$ = $f' g + f g'$, \item$\left(\frac{f}{g}\right)' = \frac{f' g - f g'}{g^2}, \quad g \neq 0$
    \item$    h = f(g) \Rightarrow h' = (f' (g)) \cdot g'.$

\end{itemize}
\subsection*{Teorema di Lagrange (del valore medio)}
\addcontentsline{toc}{subsection}{Teorema di Lagrange}
Sia \( f : [a,b] \to \mathbb{R} \) continua su \([a,b]\) e derivabile su \((a,b)\).  
Allora esiste \( c \in (a,b) \) tale che
\[
f'(c) = \frac{f(b) - f(a)}{b - a}.
\]

\subsection*{Teorema di de L'Hôpital}
\addcontentsline{toc}{subsection}{Teorema di de L'Hôpital}

Se  
\[
\lim_{x \to x_0} f(x) = \lim_{x \to x_0} g(x) = 0 \quad \text{o entrambi } \pm \infty,
\]
e \( g'(x) \neq 0 \) vicino a \( x_0 \), allora se esiste il limite  
\[
\lim_{x \to x_0} \frac{f'(x)}{g'(x)},
\]
vale  
\[
\lim_{x \to x_0} \frac{f(x)}{g(x)} = \lim_{x \to x_0} \frac{f'(x)}{g'(x)}.
\]

\subsection*{Serie di Taylor}
\addcontentsline{toc}{subsection}{Serie di Taylor}

Sia \( f \in C^\infty((x_0 - R, x_0 + R)) \). Se esistono costanti \( M, L \geq 0 \) tali che
\[
\forall n \in \mathbb{N}, \quad \forall x \in (x_0 - R, x_0 + R): \quad |f^{(n)}(x)| \leq M L^n,
\]
allora la serie di Taylor di \( f \) in \( x_0 \),
\[
\sum_{n=0}^\infty \frac{f^{(n)}(x_0)}{n!}(x - x_0)^n,
\]
converge uniformemente a \( f(x) \) per ogni \( x \in (x_0 - R, x_0 + R) \), cioè
\[
f(x) = \sum_{n=0}^\infty \frac{f^{(n)}(x_0)}{n!}(x - x_0)^n
= \sum_{k=0}^n \frac{f^{(k)}(x_0)}{k!} (x - x_0)^k + R_n(x),
\]
dove \( R_n(x) \) è il \emph{resto}

\paragraph{Definizione di \( o \)-piccolo}
Sia \( g(x) \) funzione definita in un intorno di \( x_0 \).  
Una funzione \( f(x) \) è \( o(g(x)) \) per x \(\to x_0 \) se
\[
\lim_{x \to x_0} \frac{f(x)}{g(x)} = 0.
\]
Si scrive anche
\[
f(x) = o(g(x)) \quad \text{per } x \to x_0.
\]

\subsubsection*{Resto di Peano:}
\addcontentsline{toc}{subsubsection}{Resto di Peano}

\[
R_n(x) = o((x - x_0)^n) \iff \lim_{x \to x_0} \frac{R_n(x)}{(x - x_0)^n} = 0.
\]

\subsubsection*{Resto di Lagrange:}
\addcontentsline{toc}{subsubsection}{Resto di Lagrange}

\[
\exists \xi \in (x_0, x) : R_n(x) = \frac{f^{(n+1)}(\xi)}{(n+1)!} (x - x_0)^{n+1}
\]

\subsection*{Altro}
\addcontentsline{toc}{subsection}{Altro}
\paragraph{Monotonia}
\( f \) è monotona $\underbrace{crescente}_{(decrescente)}$ in \( I \) se
\[
x_1 < x_2 \implies f(x_1) \underbrace{\leq}_{\geq} f(x_2) \quad\forall (x_1,x_2) \in I\times I
\]

\paragraph{Derivate successive}
Si definiscono
\[
f^{(k)}(x) := \frac{d^k}{dx^k} f(x).
\]

\paragraph{Integrale di Riemann}
Sia \( f : [a,b] \to \mathbb{R} \) limitata.  
Il suo integrale, se esiste, è
\[
\int_a^b f(x) dx := \lim_{\|P\| \to 0} \sum_{i=1}^n f(\xi_i)(x_i - x_{i-1}),
\]
dove \( P = \{a = x_0 < \dots < x_n = b\} \) è una partizione e \( \xi_i \in [x_{i-1}, x_i] \).

\paragraph{Teorema fondamentale del calcolo integrale}
Se \( f \) è continua in \([a,b]\) e
\[
F(x) := \int_a^x f(t) dt,
\]
allora
\[
F'(x) = f(x).
\]





\newpage
\section*{Richiami di Analisi 2}
\addcontentsline{toc}{section}{Richiami di Analisi 2}
\subsection*{Funzioni a più variabili (\( \mathbb{R}^n \to \mathbb{R} \))}
\addcontentsline{toc}{subsection}{Funzioni a più variabili (\( \mathbb{R}^n \to \mathbb{R} \))}

Sia \( f : \mathbb{R}^n \to \mathbb{R} \) una funzione scalare di \( n \) variabili reali. Supponiamo che \( f \) sia sufficientemente regolare (cioè \( k \)-volte continuamente derivabile) in un intorno di un punto \( \vec{x}_0 \in \mathbb{R}^n \).

\textbf{Gradiente:}

\[
\nabla f(\vec{x}_0) := \left( \frac{\partial f}{\partial x_1}(\vec{x}_0), \dots, \frac{\partial f}{\partial x_n}(\vec{x}_0) \right) \in \mathbb{R}^n
\]

\textbf{Hessiana (2° derivata):}

\[
\left[\nabla^2 f(\vec{x}_0)\right]_{ij} := \frac{\partial^2 f}{\partial x_i \partial x_j}(\vec{x}_0)
\quad \text{è una matrice simmetrica } \mathbb{R}^{n \times n}
\]

\textbf{Derivata di terzo ordine:} è un tensore trilineare simmetrico:

\[
\nabla^3 f(\vec{x}_0)[\vec{h}, \vec{h}, \vec{h}] := 
\sum_{i,j,k=1}^{n} 
\frac{\partial^3 f}{\partial x_i \partial x_j \partial x_k}(\vec{x}_0) 
\, h_i h_j h_k
\quad \in \mathbb{R}
\]

\textbf{Sviluppo di Taylor fino al terzo ordine:}

\[
f(\vec{x}_0 + \vec{h}) = 
f(\vec{x}_0) 
+ \nabla f(\vec{x}_0) \cdot \vec{h} 
+ \frac{1}{2} \vec{h}^\top \nabla^2 f(\vec{x}_0) \vec{h}
+ \frac{1}{6} \nabla^3 f(\vec{x}_0)[\vec{h}, \vec{h}, \vec{h}]
+ o(\|\vec{h}\|_2^3)
\]

\textbf{Sviluppo di Taylor generale (ordine \( k \)):}

Per \( f \in C^k(\mathbb{R}^n) \), lo sviluppo di Taylor attorno a \( \vec{x}_0 \) è:

\[
f(\vec{x}_0 + \vec{h}) = \sum_{r=0}^{k} \frac{1}{r!} \nabla^r f(\vec{x}_0)[\underbrace{\vec{h}, \dots, \vec{h}}_{r \text{ volte}}] + o(\|\vec{h}\|_2^k)
\]

dove: \begin{itemize}
\item \( \nabla^r f(\vec{x}_0) \) è il tensore simmetrico di ordine \( r \) delle derivate parziali,
\item il simbolo \( \nabla^r f(\vec{x}_0)[\vec{h}, \dots, \vec{h}] \) denota la contrazione multilineare del tensore lungo \( r \) copie del vettore \( \vec{h} \),
\item \( \|\vec{h}\|_2 \) è la norma euclidea.\end{itemize}

\textbf{Nota concettuale:}  
Il termine di ordine \( r \) quantifica la variazione locale di \( f \) lungo \( \vec{h} \), tenendo conto della "curvatura", "torsione", e derivate più alte della funzione. Ogni derivata aggiuntiva descrive comportamenti locali sempre più fini.


\subsection*{Differenziabilità di funzioni scalari}
\addcontentsline{toc}{subsection}{Differenziabilità di funzioni scalari}

Sia \( f : \mathbb{R}^n \to \mathbb{R} \), e sia \( \vec{x}_0 \in \mathbb{R}^n \).

La funzione è \emph{differenziabile} in \( \vec{x}_0 \) se:

\[
\exists \vec{L} \in \mathbb{R}^n \text{ tale che } 
\lim_{\vec{h} \to \vec{0}} \frac{f(\vec{x}_0 + \vec{h}) - f(\vec{x}_0) - \vec{L} \cdot \vec{h}}{\|\vec{h}\|_2} = 0
\]

In tal caso \( \vec{L} = \nabla f(\vec{x}_0) \).

Concettualmente: \( f \) è localmente approssimabile da una funzione lineare.

\textbf{Nota:} questa è l'estensione naturale della derivabilità in una variabile come limite del rapporto incrementale.


\subsection*{Differenziabilità di ordine \( n \)}
\addcontentsline{toc}{subsection}{Differenziabilità di ordine \( n \)}

Sia \( f : \mathbb{R}^m \to \mathbb{R}^p \). La funzione \( f \) è \textbf{differenziabile di ordine \( n \)} in un punto \( \vec{x}_0 \) se esistono i differenziali multilineari simmetrici di ordine \( k \), 
\[ D^k f(\vec{x}_0) : (\mathbb{R}^m)^k \to \mathbb{R}^p \quad \text{per } k = 1, 2, \ldots, n \]
tali che

\[
\lim_{\|\vec{h}\|_2 \to 0} \frac{
\left\| f(\vec{x}_0 + \vec{h}) - f(\vec{x}_0) - \sum_{k=1}^n \frac{1}{k!} D^k f(\vec{x}_0)[\vec{h}, \ldots, \vec{h}] \right\|_2
}{
\|\vec{h}\|_2^n
} = 0
\]

Dove:

\begin{itemize}
    \item \( D^1 f(\vec{x}_0) = J_f(\vec{x}_0) \) è lo Jacobiano (derivata prima, linearizzazione),
    \item \( D^2 f(\vec{x}_0) \) è il tensore del secondo ordine (derivata seconda multilineare simmetrica),
    \item \( \ldots \)
    \item \( D^n f(\vec{x}_0) \) è il differenziale multilineare simmetrico di ordine \( n \).
\end{itemize}

In altre parole, la funzione \( f \) è approssimabile attorno a \( \vec{x}_0 \) da un polinomio multivariato di grado \( n \):

\[
f(\vec{x}_0 + \vec{h}) = f(\vec{x}_0) + \sum_{k=1}^n \frac{1}{k!} D^k f(\vec{x}_0)[\vec{h}, \ldots, \vec{h}] + o(\|\vec{h}\|_2^n)
\]

La richiesta di differenziabilità di ordine \( n \) è necessaria per poter definire e calcolare i differenziali di ordine \( 2, 3, \ldots, n \), che descrivono come la funzione varia localmente con precisione crescente, fino a includere curvature, variazioni di curvature, ecc.

\subsection*{Campi vettoriali (\( \mathbb{R}^n \to \mathbb{R}^m \))}
\addcontentsline{toc}{subsection}{Campi vettoriali (\( \mathbb{R}^n \to \mathbb{R}^m \))}

Sia \( F : \mathbb{R}^n \to \mathbb{R}^m \) un'applicazione vettoriale, detta \textbf{campo vettoriale}.  
Supponiamo che \( F \) sia \( k \)-volte continuamente derivabile in un intorno di un punto \( \vec{x}_0 \in \mathbb{R}^n \). Allora esiste uno sviluppo di Taylor per \( F \) attorno a \( \vec{x}_0 \), in direzione \( \vec{h} \in \mathbb{R}^n \), del tipo:

\[
F(\vec{x}_0 + \vec{h}) = \sum_{r=0}^{k} \frac{1}{r!} D^r F(\vec{x}_0)[\underbrace{\vec{h}, \dots, \vec{h}}_{r \text{ volte}}] + o(\|\vec{h}\|_2^k)
\]

\textbf{Significato dei termini:}
\begin{itemize}
  \item \( D^r F(\vec{x}_0) \) è il \textbf{differenziale di ordine \( r \)} di \( F \) in \( \vec{x}_0 \): un tensore multilineare simmetrico di tipo \( (r,1) \), ovvero:
  \[
  D^r F(\vec{x}_0) \in \mathbb{R}^{m \times n^r}
  \]
  \item Per esempio:
  \begin{itemize}
    \item \( D^1 F(\vec{x}_0) = J_F(\vec{x}_0) \in \mathbb{R}^{m \times n} \): la \textbf{Jacobiana}, matrice delle derivate parziali prime.
    \item \( D^2 F(\vec{x}_0) \in \mathbb{R}^{m \times n \times n} \): contiene le derivate seconde di ogni componente di \( F \).
    \item \( D^3 F(\vec{x}_0) \in \mathbb{R}^{m \times n \times n \times n} \): contiene tutte le derivate terze.
  \end{itemize}
  \item Il termine \( D^r F(\vec{x}_0)[\vec{h}, \dots, \vec{h}] \in \mathbb{R}^m \) è ottenuto contrattando \( D^r F \) lungo \( r \) copie del vettore \( \vec{h} \).
\end{itemize}

\textbf{Esplicitamente:}
\[
D^r F(\vec{x}_0)[\vec{h}, \dots, \vec{h}] = 
\sum_{i_1, \dots, i_r = 1}^n 
\left( \frac{\partial^r F}{\partial x_{i_1} \dots \partial x_{i_r}} (\vec{x}_0) \right) h_{i_1} \dots h_{i_r}
\]

dove la derivata \( \frac{\partial^r F}{\partial x_{i_1} \dots \partial x_{i_r}} \) è intesa componente per componente, cioè per ogni \( \alpha = 1,\dots,m \):

\[
\left[\frac{\partial^r F}{\partial x_{i_1} \dots \partial x_{i_r}}\right]_\alpha = \frac{\partial^r F_\alpha}{\partial x_{i_1} \dots \partial x_{i_r}}
\]

\textbf{Osservazione:}  
Lo sviluppo di Taylor per \( F : \mathbb{R}^n \to \mathbb{R}^m \) è formalmente identico a quello per le funzioni scalari, ma ogni termine è un vettore in \( \mathbb{R}^m \), e le derivate sono matrici o tensori \( m \)-componenti.

\paragraph{Riassumendo:}

\[
F(\vec{x}_0 + \vec{h}) \approx F(\vec{x}_0) 
+ D F(\vec{x}_0)[\vec{h}] 
+ \frac{1}{2} D^2 F(\vec{x}_0)[\vec{h}, \vec{h}] 
+ \frac{1}{6} D^3 F(\vec{x}_0)[\vec{h}, \vec{h}, \vec{h}] 
+ \ldots
\]

\paragraph{Spiegazione:}

Questa formula esprime l'approssimazione di un campo vettoriale \( F \) vicino a un punto \( \vec{x}_0 \) tramite una somma di contributi di ordine crescente:

\begin{itemize}
    \item Il primo termine \( F(\vec{x}_0) \) è il valore del campo nel punto di partenza.
    \item Il termine \( D F(\vec{x}_0)[\vec{h}] \) rappresenta la variazione lineare di \( F \) nella direzione \( \vec{h} \).
    \item Il termine \( \frac{1}{2} D^2 F(\vec{x}_0)[\vec{h}, \vec{h}] \) descrive la curvatura o come la variazione cambia in modo non lineare.
    \item I termini di ordine superiore misurano variazioni sempre più fini, come la variazione della curvatura stessa.
\end{itemize}

Queste informazioni sono utili per prevedere il comportamento del campo vicino a \( \vec{x}_0 \) con precisione crescente. 

\paragraph{Esempio pratico:}

Se \( F \) rappresenta il campo velocità dell'aria intorno a un punto in un fluido, conoscere solo la variazione lineare (Jacobiano) può indicare la direzione e la velocità del flusso. Tuttavia, per prevedere fenomeni più complessi come turbolenze o vortici, è necessario considerare i termini di ordine superiore che descrivono come cambia questa variazione nello spazio.

\subsection*{Derivata Direzionale}
\addcontentsline{toc}{subsection}{Derivata Direzionale}

\textbf{Definizione:}

Sia \( f: \mathbb{R}^n \to \mathbb{R} \), \( \vec{x} \in \mathbb{R}^n \), e \( \vec{v} \in \mathbb{R}^n \) un \textbf{versore} (cioè \( \|\vec{v}\|_2 = 1 \)). La derivata direzionale di \( f \) in \( \vec{x} \) lungo \( \vec{v} \) è:

\[
D_{\vec{v}} f(\vec{x}) := \lim_{h \to 0} \frac{f(\vec{x} + h \vec{v}) - f(\vec{x})}{h}
\]

Se \( f \) è differenziabile:

\[
D_{\vec{v}} f(\vec{x}) = \nabla f(\vec{x}) \cdot \vec{v}
\]

\textbf{Interpretazione:} misura la variazione di \( f \) in direzione \( \vec{v} \), cioè la “pendenza” di \( f \) in quel verso.
\begin{itemize}
\item \( D_{\vec{v}} f(\vec{x}) > 0 \): \( f \) cresce lungo \( \vec{v} \)  
\item \( D_{\vec{v}} f(\vec{x}) < 0 \): \( f \) decresce  
\item \( D_{\vec{v}} f(\vec{x}) = 0 \): \( f \) è localmente costante lungo \( \vec{v} \)
\end{itemize}
\subsection*{Gradiente come direzione di massima salita}
\addcontentsline{toc}{subsection}{Gradiente come direzione di massima salita}
Sia \( f \) una funzione differenziabile in un punto \((a,b)\) e sia \(\vec{v}\) un vettore unitario, cioè \(\|\vec{v}\|_2 = 1\).

La derivata direzionale di \( f \) in \((a,b)\) nella direzione \(\mathbf{v}\) è data dal prodotto scalare tra il gradiente e \(\mathbf{v}\):
\[
D_{\mathbf{v}} f(a,b) = \nabla f(a,b) \cdot \vec{v}
\]

Per trovare la direzione di crescita massima, vogliamo massimizzare \( D_{\mathbf{v}} f(a,b) \) rispetto a \(\vec{v}\). Poiché \(\nabla f(a,b)\) è fisso, la massimizzazione dipende da \(\vec{v}\).

Ricordando che il prodotto scalare si esprime come
\[
\nabla f(a,b) \cdot \vec{v} = \|\nabla f(a,b)\| \, \|\vec{v}\| \cos \theta = \|\nabla f(a,b)\| \cos \theta,
\]
dove \(\theta\) è l'angolo tra \(\nabla f(a,b)\) e \(\vec{v}\), il massimo si ottiene quando \(\cos \theta = 1\), cioè \(\theta = 0\).

Quindi, la direzione di massima crescita è quella in cui \(\vec{v}\) è parallelo e concorde al gradiente:
\[
\vec{v}_{\max} = \frac{\nabla f(a,b)}{\|\nabla f(a,b)\|}.
\]

\subsection*{Antigradiente come direzione di massima discesa}
\addcontentsline{toc}{subsection}{Antigradiente come direzione di massima discesa}

Ora vogliamo trovare una direzione (versore v) che minimizzi la derivata direzionale in un punto (a,b). Cerchiamo $\vec{v}$ tale da minimizzare $\nabla f(a,b)\cdot \vec{v}$

Per trovare la direzione di massima discesa, vogliamo minimizzare \( D_{\mathbf{v}} f(a,b) \) rispetto a \(\vec{v}\). Poiché \(\nabla f(a,b)\) è fisso, la minimizzazione dipende da \(\vec{v}\).

Come avevamo scritto prima
\[
\nabla f(a,b) \cdot \vec{v} = \|\nabla f(a,b)\| \, \|\vec{v}\| \cos \theta = \|\nabla f(a,b)\| \cos \theta,
\]
dove \(\theta\) è l'angolo tra \(\nabla f(a,b)\) e \(\vec{v}\), il minimo si ottiene quando \(\cos \theta = -1\), cioè \(\theta = \frac{\pi}{2}\).

Quindi, la direzione di massima discesa è quella in cui \(\vec{v}\) è opposto al gradiente:
\[
\vec{v}_{\min} = -\frac{\nabla f(a,b)}{\|\nabla f(a,b)\|}.
\]


\section*{Norme vettoriali e matriciali}
\addcontentsline{toc}{section}{Norme vettoriali e matriciali}
\subsection*{Norma su $\mathbb{R}^n$}
\addcontentsline{toc}{subsection}{Norme vettoriali}

Una \textbf{norma} è una funzione $\|\cdot\| : \mathbb{R}^n \to \mathbb{R}$ che soddisfa, per ogni $x,y \in \mathbb{R}^n$ e $\alpha \in \mathbb{R}$:

\begin{itemize}
  \item $\|x\| \ge 0$ (positività)
  \item $\|x\| = 0 \Leftrightarrow x = 0$ (definitività)
  \item $\|\alpha x\| = |\alpha| \cdot \|x\|$ (omogeneità)
  \item $\|x + y\| \le \|x\| + \|y\|$ (disuguaglianza triangolare)
\end{itemize}

\subsection*{Norma $p$ su $\mathbb{R}^n$}

Per $p \in [1, \infty)$ definiamo:
\[
\|x\|_p := \left( \sum_{i=1}^n |x_i|^p \right)^{1/p}
\quad \text{e per } p = \infty: \quad
\|x\|_\infty := \max_{1 \le i \le n} |x_i|
\]

\paragraph{Positività:} ogni $|x_i|^p \ge 0$ quindi $\|x\|_p \ge 0$.

\paragraph{Definitività:} $\|x\|_p = 0 \Leftrightarrow \sum |x_i|^p = 0 \Leftrightarrow x_i = 0 \, \forall i \Rightarrow x = 0$.

\paragraph{Omogeneità:}

$\|\alpha x\|_p = \left( \sum |\alpha x_i|^p \right)^{1/p} = |\alpha| \left( \sum |x_i|^p \right)^{1/p} = |\alpha| \|x\|_p$


\paragraph{Triangolare:} 
$
\|x + y\|_p \le \|x\|_p + \|y\|_p
$ (Disuguaglianza di Minkoswky)

\subsection*{Norma $p$ con $p$ = $\infty$}

Sia $x \in \mathbb{R}^n$, definiamo:
\[
\|x\|_p := \left( \sum_{i=1}^n |x_i|^p \right)^{1/p}, \quad M := \max_{1 \le i \le n} |x_i|
\]

\[
|x_i| \le M\ \forall i \in \{1,\dots,n\}\Rightarrow \sum_{i=1}^n |x_i|^p \le \sum_{i=1}^n M^p = n M^p
\]

\[
\exists j : |x_j| = M \Rightarrow \sum_{i=1}^n |x_i|^p \ge M^p
\]

\[
\Rightarrow M^p \le \sum_{i=1}^n |x_i|^p \le n M^p
\Rightarrow M \le \left( \sum_{i=1}^n |x_i|^p \right)^{1/p} \le M \cdot n^{1/p}
\]

\[
n^{1/p} \to 1 \text{ per } p \to \infty \Rightarrow \|x\|_p \to M
\]

\[
\Rightarrow \|x\|_\infty := \lim_{p \to \infty} \|x\|_p = \max_i |x_i|
\]


\newpage
\subsection*{Norma su $\mathbb{R}^{m \times n}$}
\addcontentsline{toc}{subsection}{Norme matriciali}

Una funzione $\|\cdot\| : \mathbb{R}^{m \times n} \to \mathbb{R}$ è una \textbf{norma} se vale:

\begin{itemize}
  \item $\|A\| \ge 0$
  \item $\|A\| = 0 \Leftrightarrow A = 0$
  \item $\|\alpha A\| = |\alpha| \cdot \|A\|$
  \item $\|A + B\| \le \|A\| + \|B\|$
\end{itemize}

\subsection*{Norma $p$ indotta su matrici}

Data una norma vettoriale $\|\cdot\|_p$, si definisce:
\[
\|A\|_p := \sup_{x \ne 0} \frac{\|Ax\|_p}{\|x\|_p} = \sup_{\|x\|_p = 1} \|Ax\|_p
\]

\paragraph{Positività:} $\|Ax\|_p \ge 0 \Rightarrow \|A\|_p \ge 0$

\paragraph{Definitività:} $A = 0 \Rightarrow Ax = 0 \Rightarrow \|A\|_p = 0$; viceversa $\|A\|_p = 0 \Rightarrow \|Ax\|_p = 0 \Rightarrow Ax = 0 \Rightarrow A = 0$

\paragraph{Omogeneità:} $\|\alpha A\|_p = |\alpha| \cdot \|A\|_p$

\paragraph{Triangolare:} $\|A + B\|_p \le \|A\|_p + \|B\|_p$

\subsection*{Norma $p$ con $p=1$}

Per vettori: $\|x\|_1 = \sum_j |x_j|$.  
Sia $A \in \mathbb{R}^{m \times n}$, $x \in \mathbb{R}^n$, allora:

\[\|A\|_1 = \max_{1 \le j \le n} \sum_{i=1}^m |a_{ij}|\]
\paragraph{Dimostrazione (1): $\|A\|_1 \le \max_{1 \le j \le n} \sum_{i=1}^m |a_{ij}|$}

\[
\|Ax\|_1 = \sum_{i=1}^m \left| \sum_{j=1}^n a_{ij} x_j \right|
\]

\[
\Rightarrow \left| \sum_j a_{ij} x_j \right| \le \sum_j |a_{ij} x_j| = \sum_j |a_{ij}| \cdot |x_j|
\quad \Rightarrow \quad
\|Ax\|_1 \le \sum_{i=1}^m \sum_{j=1}^n |a_{ij}| \cdot |x_j|
\]

\[
\sum_{i=1}^m \sum_{j=1}^n |a_{ij}| \cdot |x_j| = \sum_{j=1}^n |x_j| \cdot \sum_{i=1}^m |a_{ij}|
\]

\[
\|x\|_1 = \sum_{j=1}^n |x_j| = 1
\Rightarrow
\sum_{j=1}^n |x_j| \cdot \sum_{i=1}^m |a_{ij}| \le 
\left( \max_{1 \le j \le n} \sum_{i=1}^m |a_{ij}| \right) \cdot \sum_{j=1}^n |x_j| = \max_{1 \le j \le n} \sum_{i=1}^m |a_{ij}|
\]

\[
\Rightarrow \|Ax\|_1 \le \max_{1 \le j \le n} \sum_{i=1}^m |a_{ij}|
\Rightarrow \|A\|_1 := \sup_{\|x\|_1 = 1} \|Ax\|_1 \le \max_{1 \le j \le n} \sum_{i=1}^m |a_{ij}|
\]

\paragraph{Dimostrazione (2): $\|A\|_1 \ge \max_{1 \le j \le n} \sum_{i=1}^m |a_{ij}|$}

\[\text{Sia }j^* \in \{1, \dots, n\} \text{ tale che: }\sum_{i=1}^m |a_{i j^*}| = \max_{1 \le j \le n} \sum_{i=1}^m |a_{ij}|
\]

Definiamo $x \in \mathbb{R}^n$ come:
\[
x_j = \begin{cases}
\operatorname{sign}(a_{i j^*}) & \text{se } j = j^* \\
0 & \text{altrimenti}
\end{cases}
\Rightarrow \|x\|_1 = |x_{j^*}| = 1
\]

\[
(Ax)_i = \sum_{j=1}^n a_{ij} x_j = a_{i j^*} \cdot \operatorname{sign}(a_{i j^*})
= |a_{i j^*}| \Rightarrow \|Ax\|_1 = \sum_{i=1}^m |a_{i j^*}|
\]

\[
\Rightarrow \sup_{\|x\|_1 = 1} \|Ax\|_1 \ge \|Ax\|_1 = \sum_{i=1}^m |a_{i j^*}|
= \max_{1 \le j \le n} \sum_{i=1}^m |a_{ij}|
\]

\[
\Rightarrow \|A\|_1 \ge \max_{1 \le j \le n} \sum_{i=1}^m |a_{ij}|
\]
\hfill $\square$
\[
A =
\begin{bmatrix}
a_{11} & \cdots & \color{red}{a_{1j^*}} & \cdots & a_{1n} \\
a_{21} & \cdots & \color{red}{a_{2j^*}} & \cdots & a_{2n} \\
\vdots & \ddots & \vdots & \ddots & \vdots \\
a_{m1} & \cdots & \color{red}{a_{mj^*}} & \cdots & a_{mn}
\end{bmatrix}\Rightarrow \|A\|_1 = \max_{1 \le j \le n} \sum_{i=1}^m |a_{ij}|
\]

\subsection*{Norma $\|A\|_\infty$}

Per vettori: $\|x\|_\infty = \max_{1 \le j \le n} |x_j|$  
Sia $A \in \mathbb{R}^{m \times n}$, $x \in \mathbb{R}^n$, allora:

\[
\|A\|_\infty := \sup_{\|x\|_\infty = 1} \|Ax\|_\infty = \sup_{\max_{1 \le j \le n} |x_j| = 1} \max_{1 \le i \le m} \left| \sum_{j=1}^n a_{ij} x_j \right|
\]

\paragraph{Dimostrazione (1): $\|A\|_\infty \le \max_{1 \le i \le m} \sum_{j=1}^n |a_{ij}|$}

\[
\left| \sum_{j=1}^n a_{ij} x_j \right| \le \sum_{j=1}^n |a_{ij}| \cdot |x_j|
\le \left( \max_{1 \le j \le n} |x_j| \right) \cdot \sum_{j=1}^n |a_{ij}| = \sum_{j=1}^n |a_{ij}|
\]

\[
\Rightarrow \|Ax\|_\infty = \max_{1 \le i \le m} \left| \sum_{j=1}^n a_{ij} x_j \right| \le \max_{1 \le i \le m} \sum_{j=1}^n |a_{ij}|
\Rightarrow \|A\|_\infty \le \max_{1 \le i \le m} \sum_{j=1}^n |a_{ij}|
\]

\paragraph{Dimostrazione (2): $\|A\|_\infty \ge \max_{1 \le i \le m} \sum_{j=1}^n |a_{ij}|$}

Sia $i^* \in \{1, \dots, m\}$ tale che:
\[
\sum_{j=1}^n |a_{i^* j}| = \max_{1 \le i \le m} \sum_{j=1}^n |a_{ij}|
\]

Definiamo $x \in \mathbb{R}^n$ come:
\[
x_j = \operatorname{sign}(a_{i^* j}) \Rightarrow |x_j| = 1
\Rightarrow \|x\|_\infty = \max_j |x_j| = 1
\]

\[
(Ax)_{i^*} = \sum_{j=1}^n a_{i^* j} \cdot \operatorname{sign}(a_{i^* j}) = \sum_{j=1}^n |a_{i^* j}|
\Rightarrow \|Ax\|_\infty \ge |(Ax)_{i^*}| = \sum_{j=1}^n |a_{i^* j}|
\]

\[
\Rightarrow \|A\|_\infty \ge \max_{1 \le i \le m} \sum_{j=1}^n |a_{ij}|
\]
\hfill $\square$

\[
A =
\begin{bmatrix}
a_{11} & a_{12} & \cdots & a_{1n} \\
\vdots & \vdots & \ddots & \vdots \\
\color{red}{a_{i^*1}} & \color{red}{a_{i^*2}} & \cdots & \color{red}{a_{i^*n}} \\
\vdots & \vdots & \ddots & \vdots \\
a_{m1} & a_{m2} & \cdots & a_{mn}
\end{bmatrix}
\Rightarrow \|A\|_\infty = \max_{1 \le i \le m} \sum_{j=1}^n |a_{ij}|
\]


\subsection*{Norma $\|A\|_2$}

Per $x \in \mathbb{R}^n$, $\|x\|_2 := \left( \sum_{j=1}^n |x_j|^2 \right)^{1/2}$.\\
Sia $A \in \mathbb{R}^{m \times n}$, allora:

\[
\|A\|_2 := \sup_{\|x\|_2 = 1} \|Ax\|_2 = \sup_{\|x\|_2 = 1} \sqrt{(Ax)^T (Ax)} = \sup_{\|x\|_2 = 1} \sqrt{x^T A^T A x}.
\]

Poiché $A^T A \in \mathbb{R}^{n \times n}$ è simmetrica possiede un sistema completo di autovettori ortonormali $\{v_i\}_{i=1}^n$ tali che $A^T A v_i = \sigma_i^2 v_i$, con $\sigma_i \geq 0$ (valori singolari di $A$).\\
Qualsiasi $x \in \mathbb{R}^n$ con $\|x\|_2 = 1$ si può scrivere come:
\[
x = \sum_{i=1}^n \alpha_i v_i, \quad \text{con } \sum_{i=1}^n \alpha_i^2 = 1.
\]

Allora:
\[
A^T A x = \sum_{i=1}^n \alpha_i A^T A v_i = \sum_{i=1}^n \alpha_i \sigma_i^2 v_i,
\]
\[
x^T A^T A x = \left( \sum_{i=1}^n \alpha_i v_i \right)^T \left( \sum_{i=1}^n \alpha_i \sigma_i^2 v_i \right) = \sum_{i=1}^n \alpha_i^2 \sigma_i^2,
\]
poiché $\{v_i\}$ è ortonormale.

Quindi:
\[
x^T A^T A x \leq \sigma_1^2 \sum_{i=1}^n \alpha_i^2 = \sigma_1^2 \cdot 1 = \sigma_1^2 \quad \text{con }\sigma_1=\max_{1\leq i \leq n}\sigma_i
\]

Il massimo si ottiene per $x = v_1$, cioè:
\[
\|A\|_2 = \sup_{\|x\|_2 = 1} \sqrt{x^T A^T A x} = \sqrt{\sigma_1^2} = \sigma_1.
\]
\hfill $\square$


\subsection*{Norma di Frobenius $\|A\|_F$}

Definita come:
\[
\|A\|_F := \left( \sum_{i=1}^m \sum_{j=1}^n |a_{ij}|^2 \right)^{1/2}
\]

Equivale ad applicare $\| \cdot \|_2$ al vettore ottenuto appiattendo la matrice $A$ in $\mathbb{R}^{mn}$.

\textbf{Proprietà:}
\begin{itemize}
  \item $\|A\|_F \ge 0$
  \item $\|A\|_F = 0 \Leftrightarrow a_{ij} = 0\ \forall i,j$
  \item $\|\alpha A\|_F = |\alpha| \cdot \|A\|_F$
  \item $\|A + B\|_F \le \|A\|_F + \|B\|_F$ (disuguaglianza triangolare)
\end{itemize}
\newpage
\section*{Numero di condizionamento}
\addcontentsline{toc}{section}{Numero di condizionamento}

\begin{flushleft}
$\displaystyle
\kappa_p(A) = \|A\|_p \cdot \|A^{-1}\|_p \quad \text{dove } A \in \mathbb{R}^{n \times n} \text{ non singolare}
$

\vspace{0.5em}
\vspace{0.5em}

Norme matriciali più comuni:

\begin{itemize}
\item Norma 1: $\|A\|_1 = \max\limits_{1 \leq j \leq n} \sum_{i=1}^n |a_{ij}|$
\item Norma 2 (spettrale): $\|A\|_2 = \sqrt{\rho(A^TA)}$
\vspace{0.8em}

\item Norma $\infty$: $\|A\|_\infty = \max\limits_{1 \leq i \leq n} \sum_{j=1}^n |a_{ij}|$
\end{itemize}
\end{flushleft}

\begin{lstlisting}
function k = cond_number(A, p)
    % Calcola il numero di condizionamento in norma p
    % Input:
    %   A - matrice quadrata non singolare
    %   p - tipo di norma (1, 2, inf, 'fro')
    % Output:
    %   k - numero di condizionamento Kp(A)
    
    if nargin < 2
        p = 2; % default norma 2
    end
    
    % Controllo singolarità
    if det(A) == 0
        error('La matrice è singolare');
    end
    
    % Calcolo norma matrice e inversa
    norma_A = norm(A, p);
    norma_A_inv = norm(inv(A), p);
    
    % Numero di condizionamento
    k = norma_A * norma_A_inv;
    
end
\end{lstlisting}


\newpage
\subsection*{Numero di condizionamento di matrici simmetriche}
\addcontentsline{toc}{subsection}{Numero di condizionamento di matrici simmetriche}

\begin{flushleft}
$\displaystyle
\kappa_2(A_{kn}) = \|A_{kn}\|_2 \cdot \|A_{kn}^{-1}\|_2 \quad \text{dove } A_{kn} \in \mathbb{R}^{n \times n}
$ è una matrice definita da
\end{flushleft}

\[
A_{kn} = 
\begin{bmatrix}
k & -1 & 0 & \cdots & 0 \\
-1 & k & -1 & \ddots & \vdots \\
0 & -1 & \ddots & \ddots & 0 \\
\vdots & \ddots & \ddots & k & -1 \\
0 & \cdots & 0 & -1 & k
\end{bmatrix}
\quad 
\]
\vspace{0.5em}

\begin{lstlisting}
% Funzione che costruisce la matrice tridiagonale Akn
function A = Akn(n, k)
    main_diag = k * ones(n, 1);     % Diagonale principale
    off_diag = -1 * ones(n-1, 1);   % Diagonali sopra e sotto
    A = diag(main_diag) + diag(off_diag, 1) + diag(off_diag, -1);  % Somma delle diagonali
end

function cond_num = condizionamento(k, n)
    % Crea la matrice tridiagonale Akn
    A = Akn(n, k);
    
    % Calcola gli autovalori della matrice
    autovalori = eig(A);
    
    % Calcola la norma 2 della matrice (massimo autovalore assoluto)
    norma_A = max(abs(autovalori));
    
    % Calcola la norma 2 dell'inversa della matrice (minimo autovalore assoluto)
    norma_A_inv = 1 / min(abs(autovalori));
    
    % Numero di condizionamento
    cond_num = norma_A * norma_A_inv;
end


n = 100;  % Numero di righe/colonne della matrice
k_values = [2, 4, 1e6];  % Valori di k da testare (1e6 in MatLab significa 1*10^(6), 10 è la base dell'esponenziale)

for k = k_values
    cond_num = condizionamento(k, n);
    
    if k == 1e6
        % Per valori di k molto grandi, stampa in notazione scientifica
        fprintf('Numero di condizionamento per k = %.0e: %.4f\n', k, cond_num);
    else
        % Altrimenti, stampa normalmente
        fprintf('Numero di condizionamento per k = %d: %.4f\n', k, cond_num);
    end
end
\end{lstlisting}
\newpage
\section*{Risoluzione di sistemi triangolari}
\addcontentsline{toc}{section}{Risoluzione di sistemi triangolari}

\subsection*{Caso triangolare inferiore}
\addcontentsline{toc}{subsection}{Caso triangolare inferiore}

\begin{flushleft}
$A x = b \quad \text{dove } A \in \mathbb{R}^{n \times n} \text{ è triangolare inferiore}$

\vspace{0.5em}

$\displaystyle
x_i = \frac{1}{A_{ii}} \left(b_i - \sum_{j=1}^{i-1} A_{ij} x_j \right)
$
\end{flushleft}

\begin{lstlisting}
function x = lowertriangularsystem(A, b)
    n = length(b);
    x = zeros(n,1);
    
    x(1) = b(1) / A(1,1);
    
    for i = 2:n
        somma = A(i,1:i-1) * x(1:i-1);
        x(i) = (b(i) - somma) / A(i,i);
    end
end
\end{lstlisting}


\subsection*{Caso triangolare superiore}
\addcontentsline{toc}{subsection}{Caso triangolare superiore}
\begin{flushleft}
$A x = b \quad \text{dove } A \in \mathbb{R}^{n \times n} \text{ è triangolare superiore}$

\vspace{0.5em}

$\displaystyle
x_i = \frac{1}{A_{ii}} \left(b_i - \sum_{j=i+1}^{n} A_{ij} x_j \right)
$
\end{flushleft}

\begin{lstlisting}
function x = uppertriangularsystem(A, b)
    n = length(b);
    x = zeros(n,1);
    
    x(n) = b(n) / A(n,n);
    
    for i = n-1:-1:1
        somma = A(i,i+1:n) * x(i+1:n);
        x(i) = (b(i) - somma) / A(i,i);
    end
end
\end{lstlisting}

\section*{Fattorizzazione LU}
\addcontentsline{toc}{section}{Fattorizzazione LU}
\begin{flushleft}
\fbox{
\begin{minipage}{\linewidth}
\textsc{Input}: $A \in \mathbb{R}^{n \times n}$

\vspace{0.6em}

\(U, L = A, I_n \quad\) 
\vspace{0.5em}

Per ogni $j$ in $\{1, \dots, n\}$: \\[0.2em]
\hspace*{1.5em} $k = j$ \\[0.2em]
\hspace*{1.5em} Finché \(U_{k,j} = 0\): \\[0.2em]
\hspace*{3em} $k = k+1$ \\[0.2em]
\hspace*{3em} Se \(k > n\)\\[0.2em]
\hspace*{4.5em} $Errore$ \\[0.2em]
\hspace*{3em} Se \(U_{k,j} = 0\)\\[0.2em]
\hspace*{4.5em} Passa all'iterazione successiva \\[0.2em]
\hspace*{3em} Se \(U_{k,j} > 0\) \\[0.2em]
\hspace*{4.5em} Scambia le righe $j$ e $k$ in $U$ ed esci dal ciclo \\[0.2em]
\hspace*{1.5em} Per ogni $i$ in $\{j+1, \dots, n\}$:\\[0.2em]
\hspace*{3em} \(L_{i,j} = \frac{U_{i,j}}{U_{j,j}}\) \quad  \\[0.2em]
\hspace*{3em} Sottrai \(L_{i,j} \cdot U_{j,:}\) all'i-esima riga di \(U_{i,:}\)  \quad

\vspace{0.6em}

\textsc{Output}: $L$ triangolare inferiore, $U$ triangolare superiore $(A = LU)$
\end{minipage}
}
\end{flushleft}

\begin{lstlisting}
function [L, U] = FattorizzazioneLU(A)
    U = A;  % Inizializzo U
    [n, n] = size(U);
    L = eye(n);  % Inizializzo L
    Tol = 10^(-6) * norm(A, inf);  % Tolleranza basata sulla norma infinito di A
    for j = 1:n
        k = j;
        % Eseguiamo il pivoting sulle righe
        while abs(U(k,j)) <= Tol
            k = k + 1;
            if k > n  % Se siamo arrivati alla fine senza trovare un buon pivot
                % La matrice è singolare o mal condizionata
                L = NaN;
                U = NaN;
                return;
            end
            % Qui dentro controlliamo se il pivot è valido
            if abs(U(k,j)) > Tol  % Aggiunto all'interno del while
                % Se il pivot è valido, scambiamo le righe
                U(j,:), U(k,:) = U(k,:), U(j,:);  % Scambia le righe di U
                break;  % Esci dal while una volta trovato un pivot valido
            end
        end
        for i = j+1:n % Calcola i moltiplicatori di Gauss
            m = U(i,j) / U(j,j);  % Moltiplicatore di Gauss
            L(i,j) = m;  % Memorizza il moltiplicatore in L
            U(i,:) = U(i,:) - m * U(j,:);  % Aggiorna la riga i-esima di U
        end
    end
end

\end{lstlisting}
\newpage
\subsection*{Fattorizzazione LU con pivoting: motivazione teorica}
\addcontentsline{toc}{subsection}{Fattorizzazione LU con pivoting: motivazione teorica}

Durante l’eliminazione di Gauss, per evitare divisioni per 0 o instabilità numerica, si eseguono scambi di righe $\Rightarrow$ ogni scambio è rappresentato da una \textbf{matrice di permutazione} $P_{ik} = I_n$ con righe $i,k$ scambiate.\\
\[
P_{ik}A = \text{righe $i,k$ scambiate},\quad AP_{ik} = \text{colonne $i,k$ scambiate}
\]
\[
P_{ik}^2 = I_n \Rightarrow P_{ik}^{-1} = P_{ik}, \quad \det(P_{ik}) = -1
\]
Sia $P_k = \begin{cases} P_{ik} & \text{se righe } i,k \text{ scambiate} \\ I_n & \text{altrimenti} \end{cases}$

Sequenza eliminazione con pivoting:
\[
M_{n-1}P_{n-1} \cdots M_1P_1 A = U \Rightarrow A = P_1^{-1}M_1^{-1} \cdots P_{n-1}^{-1}M_{n-1}^{-1}U
\]
Poiché $P_k^{-1} = P_k$, possiamo scrivere:
\[
A = LU,\quad \text{dove } L = P_1M_1^{-1}P_2M_2^{-1} \cdots P_{n-1}M_{n-1}^{-1}
\]
Definendo $P = P_{n-1} \cdots P_1$ si ottiene la fattorizzazione:
$PA = LU$
(con pivoting parziale: solo righe, non colonne).

\vspace{1em}

\subsection*{Fattorizzazione LU senza pivoting}
\addcontentsline{toc}{subsection}{Fattorizzazione LU senza pivoting}

\begin{flushleft}
\fbox{
\begin{minipage}{\linewidth}
\textsc{Input}: $A \in \mathbb{R}^{n \times n}$, \quad $U := A,\ L := I_n$\\
Per $j = 1 \dots n$:\\
\hspace*{1em} Se $U_{j,j} = 0$: errore (pivot nullo)\\
\hspace*{1em} Per $i = j+1 \dots n$: \\
\hspace*{2em} $L_{i,j} := U_{i,j}/U_{j,j}$\\ \hspace*{2em}$U_{i,:} := U_{i,:} - L_{i,j} \cdot U_{j,:}$\\
\textsc{Output}: $L$ triangolare inf., $U$ triangolare sup., $A = LU$
\end{minipage}
}
\end{flushleft}

\subsection*{Applicazioni della fattorizzazione LU}
\addcontentsline{toc}{subsection}{Applicazioni della fattorizzazione LU}

\begin{itemize}
\item \textbf{Determinante}: $A = LU \Rightarrow\det(A) = \det(L)\det(U)$ $\Rightarrow \det(A) = (-1)^S\prod_{i=1}^n U_{ii}$ con S = numero di scambi di riga effettuati
\item \textbf{Rango}: $\mathrm{rank}(A) = |\{i : U_{ii} \neq 0\}|$
\item \textbf{Sistemi $Ax = b$}: risolvi $Ly = b$ (sost. in avanti), poi $Ux = y$ (sost. all'indietro)
\item \textbf{Inversa}: $\forall i \in \{1,\dots\,n\}:$ risolvi $Ax_i = e_i \Rightarrow A^{-1} = [x_1 \dots x_n]$
\item \textbf{Numero di condizionamento}: Possiamo sfruttare la fattorizzazione $LU$ per calcolare l'inversa e di conseguenza il numero di condizionamento
\end{itemize}


\newpage
\textbf{Esempio: Calcolare la fattorizzazione LU di $A = \begin{bmatrix} 2 & 3 \\ 4 & 7 \end{bmatrix}
$}
\addcontentsline{toc}{subsection}{Esempio}

\[
A = \begin{bmatrix} 2 & 3 \\ 4 & 7 \end{bmatrix}
\]

\[
A = LU,\quad
L = \begin{bmatrix} 1 & 0 \\ l_{21} & 1 \end{bmatrix},\quad
U = \begin{bmatrix} u_{11} & u_{12} \\ 0 & u_{22} \end{bmatrix}
\]

\[
LU =
\begin{bmatrix}
1 & 0 \\
l_{21} & 1
\end{bmatrix}
\begin{bmatrix}
u_{11} & u_{12} \\
0 & u_{22}
\end{bmatrix}
=
\begin{bmatrix}
u_{11} & u_{12} \\
l_{21} u_{11} & l_{21} u_{12} + u_{22}
\end{bmatrix}
\]

\[
\begin{cases}
u_{11} = 2 \\
u_{12} = 3 \\
l_{21} \cdot 2 = 4 \Rightarrow l_{21} = 2 = \frac{a_{21}}{a_{11}} \\
2 \cdot 3 + u_{22} = 7 \Rightarrow u_{22} = 1 = a_{22}-l_{21}u_{12}
\end{cases}
\]

\[
L = \begin{bmatrix} 1 & 0 \\ 2 & 1 \end{bmatrix},\quad
U = \begin{bmatrix} 2 & 3 \\ 0 & 1 \end{bmatrix},\quad
A = LU
\]

\textbf{Esempio: Calcolare la fattorizzazione LU di }  
$A = \begin{bmatrix} \varepsilon & 2 & 1 \\ 1 & -2 & 0.5 \\ 2 & 3 & 1 \end{bmatrix}$  
\addcontentsline{toc}{subsection}{Esempio}

\[
A = \begin{bmatrix} \varepsilon & 2 & 1 \\ 1 & -2 & 0.5 \\ 2 & 3 & 1 \end{bmatrix}
\]

\[
P A = LU,\quad
P = \text{matrice di permutazione}
\]

Scambio righe: \( r_1 = r_3 \Rightarrow P = I_\text{ con righe 1 e 3 scambiate} = \begin{bmatrix} 0 & 0 & 1 \\ 0 & 1 & 0 \\ 1 & 0 & 0 \end{bmatrix} \)

\[
PA = \begin{bmatrix} 2 & 3 & 1 \\ 1 & -2 & 0.5 \\ \varepsilon & 2 & 1 \end{bmatrix}
\]

\[
L = \begin{bmatrix}
1 & 0 & 0 \\
l_{21} & 1 & 0 \\
l_{31} & l_{32} & 1
\end{bmatrix},\quad
U = \begin{bmatrix}
\textcolor{red}{u_{11}} & \textcolor{red}{u_{12}} & \textcolor{red}{u_{13}} \\
0 & u_{22} & u_{23} \\
0 & 0 & u_{33}
\end{bmatrix}
\]

\[
\text{Fase 1:}
\]

\[
\textcolor{red}{u_{11}} = 2,\quad \textcolor{red}{u_{12}} = 3,\quad \textcolor{red}{u_{13}} = 1
\]

\[
\textcolor{blue}{l_{21}} = \dfrac{(PA)_{21}}{\textcolor{red}{u_{11}}} = \dfrac{1}{2},\quad \textcolor{blue}{l_{31}} = \dfrac{(PA)_{31}}{\textcolor{red}{u_{11}}} = \dfrac{\varepsilon}{2}
\]

\[
A_{2,:} = A_{2,:} - \textcolor{blue}{l_{21}} \cdot A_{1,:} =
\begin{bmatrix}
1 & -2 & 0.5
\end{bmatrix}
-
\dfrac{1}{2}
\begin{bmatrix}
\textcolor{red}{2} & \textcolor{red}{3} & \textcolor{red}{1}
\end{bmatrix}
=
\begin{bmatrix}
0 & -3.5 & 0
\end{bmatrix}
\]

\[
A_{3,:} = A_{3,:} - \textcolor{blue}{l_{31}} \cdot A_{1,:} =
\begin{bmatrix}
\varepsilon & 2 & 1
\end{bmatrix}
-
\dfrac{\varepsilon}{2}
\begin{bmatrix}
\textcolor{red}{2} & \textcolor{red}{3} & \textcolor{red}{1}
\end{bmatrix}
=
\begin{bmatrix}
0 & 2 - \dfrac{3\varepsilon}{2} & 1 - \dfrac{\varepsilon}{2}
\end{bmatrix}
\]

\[
U =
\begin{bmatrix}
\textcolor{red}{2} & \textcolor{red}{3} & \textcolor{red}{1} \\
0 & u_{22} & u_{23} \\
0 & 0 & u_{33}
\end{bmatrix}
\]

\[
\text{Fase 2:}
\]

\[
u_{22} = -3.5,\quad u_{23} = 0,\quad \textcolor{brown}{l_{32}} = \dfrac{(PA)_{32}}{u_{22}} = \dfrac{2 - \dfrac{3\varepsilon}{2}}{-3.5}
\]

\[
A_{3,:} = A_{3,:} - \textcolor{brown}{l_{32}} \cdot A_{2,:} =
\begin{bmatrix}
0 & 2 - \dfrac{3\varepsilon}{2} & 1 - \dfrac{\varepsilon}{2}
\end{bmatrix}
-
\textcolor{brown}{l_{32}} \cdot
\begin{bmatrix}
0 & -3.5 & 0
\end{bmatrix}
=
\begin{bmatrix}
0 & 0 & \textcolor{green}{1 - \dfrac{\varepsilon}{2}}
\end{bmatrix}
\]

\[
\textcolor{green}{u_{33}} = 1 - \dfrac{\varepsilon}{2}
\]

\[
L =
\begin{bmatrix}
1 & 0 & 0 \\
\textcolor{blue}{\dfrac{1}{2}} & 1 & 0 \\
\textcolor{blue}{\dfrac{\varepsilon}{2}} & \textcolor{brown}{\dfrac{2 - \dfrac{3\varepsilon}{2}}{-3.5}} & 1
\end{bmatrix},\quad
U =
\begin{bmatrix}
\textcolor{red}{2} & \textcolor{red}{3} & \textcolor{red}{1} \\
0 & -3.5 & 0 \\
0 & 0 & \textcolor{green}{1 - \dfrac{\varepsilon}{2}}
\end{bmatrix},\quad
P =
\begin{bmatrix}
0 & 0 & 1 \\
0 & 1 & 0 \\
1 & 0 & 0
\end{bmatrix}
\]

\[
PA = LU
\]

\newpage

\subsection*{Analisi della fattorizzazione LU delle matrici di Hilbert}
\addcontentsline{toc}{subsection}{Analisi della fattorizzazione LU delle matrici di Hilbert}

La \textbf{matrice di Hilbert} è una matrice quadrata \( H \in \mathbb{R}^{n \times n} \) definita da:
\[
H_{i,j} = \frac{1}{i + j - 1}, \quad \text{per } 1 \leq i,j \leq n.
\]
Essa è simmetrica, definita positiva e mal condizionata (il suo numero di condizionamento è alto) per valori crescenti di \( n \), il che la rende un banco di prova classico per lo studio della stabilità numerica degli algoritmi di fattorizzazione.

Il seguente script MATLAB analizza il comportamento del residuo 
\(
\| H - LU \|
\)
(nella norma di Frobenius) quando si applica la fattorizzazione LU alle matrici di Hilbert, al variare della dimensione \( n \).

\begin{lstlisting}[language=Matlab]
% Analisi della fattorizzazione LU delle matrici di Hilbert
clear all        % Pulisce tutte le variabili presenti nello workspace.
close all        % Chiude tutte le finestre delle figure aperte.

% Calcola il residuo per diverse dimensioni
max_n = 20;                  % Imposta la dimensione massima della matrice di Hilbert
residui = zeros(max_n, 1);   % Vettore per memorizzare i residui

for n = 2:max_n              % Ciclo su dimensioni crescenti
    % Costruzione della matrice di Hilbert H di dimensione n x n
    H = zeros(n, n);
    for i = 1:n
        for j = 1:n
            H(i, j) = 1 / (i + j - 1);   % Definizione della matrice di Hilbert
        end
    end

    % Fattorizzazione LU: H = L * U
    [L, U] = FattorizzazioneLU(H);

    % Calcolo del residuo: ||H - LU||_F
    residuo = sqrt(sum(sum((H - L*U).^2)));
    residui(n) = residuo;

    fprintf('n = %d, Residuo = %e\n', n, residuo);
end

% Visualizzazione del residuo in scala logaritmica
figure;
semilogy(2:max_n, residui(2:max_n), 'o-');
xlabel('Dimensione n');
ylabel('Residuo ||H^n - LU||');
title('Residuo della fattorizzazione LU per le matrici di Hilbert');
grid on;
\end{lstlisting}


\newpage
\subsection*{Fattorizzazione LU di matrici tridiagonali}
\addcontentsline{toc}{subsection}{Fattorizzazione LU di matrici tridiagonali}

Il codice MATLAB riportato calcola la fattorizzazione LU ottimizzata per matrici tridiagonali, sfruttando la loro struttura speciale per ridurre i calcoli rispetto a una fattorizzazione generale. Inoltre, uno script confronta i tempi di esecuzione tra questa versione specifica e una fattorizzazione LU generale.

\begin{lstlisting}[language=Matlab]
% Funzione per calcolare la fattorizzazione LU di una matrice tridiagonale
function [L,U]= tridiag_LU(Ah)
    d0=diag(Ah);              % Estrae la diagonale principale
    d1=diag(Ah,1);            % Estrae la diagonale superiore (+1)
    d2=diag(Ah,-1);           % Estrae la diagonale inferiore (-1)
    n=(size(Ah,1))-1;         % Dimensione matrice meno 1
    
    % Algoritmo ottimizzato per fattorizzazione LU
    for i = 1:n
        d2(i)=d2(i)/d0(i);              % Moltiplicatori L
        d0(i+1)=d0(i+1)-d2(i)*d1(i);   % Aggiorna diagonale U
    end
    
    U = diag(d0)+diag(d1,1);            % Costruisce U
    L = eye(n+1)+ diag(d2,-1);          % Costruisce L
end

% Script principale: confronto tempi LU tridiagonale vs generale
tdiag=zeros(199,1); t=zeros(199,1);
for n=2:200
    Ah=diag(2*ones(n,1));
    Ah=Ah-diag(ones(n-1,1),1);
    Ah=Ah-diag(ones(n-1,1),-1);
    Ah=Ah*((n+1)^2);  % Scala per discretizzazione
    
    tic; [L,U]= tridiag_LU(Ah); tdiag(n-1)=toc;
    tic; [L,U]=FattorizzazioneLU(Ah); t(n-1)=toc;
end

plot(2:200,tdiag)
hold on
plot(2:200,t)
hold off
\end{lstlisting}


\newpage
\subsection*{Fattorizzazione LU di matrici di Poisson}
\addcontentsline{toc}{subsection}{Fattorizzazione LU di matrici di Poisson}

La matrice di Poisson è una matrice sparsa tipica delle discretizzazioni a differenze finite del problema di Poisson (equazione differenziale parziale del tipo \(-\Delta u = f\)). Ha una struttura a blocchi tridiagonale e, per la dimensione \(n^2\), si presenta come

\[
A = \begin{bmatrix}
T & -I & 0 & \cdots & 0 \\
-I & T & -I & \cdots & 0 \\
0 & -I & T & \cdots & 0 \\
\vdots & \vdots & \vdots & \ddots & \vdots \\
0 & 0 & 0 & \cdots & T
\end{bmatrix}
\quad \text{dove} \quad
T = \begin{bmatrix}
4 & -1 & 0 & \cdots & 0 \\
-1 & 4 & -1 & \cdots & 0 \\
0 & -1 & 4 & \cdots & 0 \\
\vdots & \vdots & \vdots & \ddots & -1 \\
0 & 0 & 0 & -1 & 4
\end{bmatrix}
\]

Questa matrice rappresenta il sistema lineare derivante dalla discretizzazione del problema di Poisson su una griglia bidimensionale. Il nome “matrice di Poisson” deriva dal fatto che essa nasce direttamente dall’equazione di Poisson, molto usata in fisica e ingegneria per modellare fenomeni quali la diffusione del calore, l’elettrostatica e altri processi.

Il codice Matlab seguente calcola, per diverse dimensioni \(n\), il numero di elementi nulli nella matrice di Poisson \(A\) e nelle matrici triangolari \(L\) e \(U\) ottenute dalla fattorizzazione LU, per studiarne la sparsità.

\begin{lstlisting}[language=Matlab]
k = 30;  % Numero massimo di dimensioni da considerare
LL = zeros(k-1, 1);  % Zeri in L per ogni n
UU = zeros(k-1, 1);  % Zeri in U per ogni n
values = zeros(k-1, 1);  % Zeri in A per ogni n

for n = 2:k
    A = gallery("poisson", n);  % Matrice di Poisson di dimensione n^2 x n^2
    
    values(n-1) = sum(sum(A == 0));  % Conta zeri in A
    
    [L, U] = lu(A);  % Fattorizzazione LU
    
    LL(n-1) = sum(sum(L == 0));  % Conta zeri in L
    UU(n-1) = sum(sum(U == 0));  % Conta zeri in U
end

plot(2:k, LL)  % Grafico zeri in L
hold on
plot(2:k, values)  % Grafico zeri in A
hold off
\end{lstlisting}

L’operatore \emph{laplaciano} in due dimensioni è definito come

\[
\Delta u = \frac{\partial^2 u}{\partial x^2} + \frac{\partial^2 u}{\partial y^2}.
\]

Nella risoluzione numerica di problemi su domini continui, come l’equazione di Poisson

\[
-\Delta u = f,
\]

si usa una \emph{discretizzazione} tramite differenze finite su una griglia regolare.

Consideriamo una griglia 2D con passo \(h\) in entrambe le direzioni; le derivate seconde si approssimano come

\[
\frac{\partial^2 u}{\partial x^2}(x_i,y_j) \approx \frac{u_{i+1,j} - 2 u_{i,j} + u_{i-1,j}}{h^2},
\]
\[
\frac{\partial^2 u}{\partial y^2}(x_i,y_j) \approx \frac{u_{i,j+1} - 2 u_{i,j} + u_{i,j-1}}{h^2}.
\]

Sommando otteniamo

\[
-\Delta u(x_i,y_j) \approx \frac{-u_{i+1,j} - u_{i-1,j} - u_{i,j+1} - u_{i,j-1} + 4 u_{i,j}}{h^2}.
\]

La matrice \(A\) associata al sistema lineare ha quindi le seguenti caratteristiche:

\begin{itemize}
    \item sulla diagonale principale ci sono valori pari a \(4\), perché il punto \((i,j)\) dipende da se stesso con coefficiente \(4\),
    \item nelle posizioni corrispondenti ai punti adiacenti (destra, sinistra, sopra, sotto) ci sono valori \(-1\),
    \item tutti gli altri elementi sono zero.
\end{itemize}

La matrice \(T\) è la parte principale di questa matrice \(A\), che rappresenta la discretizzazione lungo una dimensione, mentre la struttura a blocchi con \(T\) e \(-I\) tiene conto della seconda dimensione.

Questa struttura è tipica nella discretizzazione di problemi ellittici come l’equazione di Poisson, e garantisce che la matrice sia:

\begin{itemize}
    \item sparsa,
    \item simmetrica,
    \item definita positiva,
\end{itemize}

proprietà utili per sviluppare metodi numerici efficienti di risoluzione.

\newpage
\section*{Sistemi lineari}
\addcontentsline{toc}{section}{Sistemi lineari}

\begin{flushleft}
\begin{minipage}{\linewidth}
Un sistema lineare in forma matriciale è:
\[
Ax = b
\]
dove:
\begin{itemize}
  \item \( A \in \mathbb{R}^{n \times n} \) è la matrice dei coefficienti,
  \item \( x \in \mathbb{R}^n \) è il vettore incognito,
  \item \( b \in \mathbb{R}^n \) è il termine noto.
\end{itemize}

\textbf{Obiettivo:} Trovare \( x \) tale che \( Ax = b \).\\

\textbf{Condizione di compatibilità:} il sistema ammette soluzione se e solo se \( b \in \text{Im}(A) \), cioè se \( b \) è combinazione lineare delle colonne di \( A \) e in quel caso $A$ è invertibile.

Quindi se \( A \) è invertibile (cioè \( \det(A) \neq 0 \)), allora:
$x = A^{-1}b$
\end{minipage}
\end{flushleft}

\section*{Metodo di Jacobi}
\addcontentsline{toc}{section}{Metodo di Jacobi}

\begin{flushleft}
\fbox{
\begin{minipage}{\linewidth}
\textsc{Input}: $A \in \mathbb{R}^{n \times n}$, $b \in \mathbb{R}^n$, $x^{(0)} \in \mathbb{R}^n$, $\varepsilon > 0$

\vspace{0.6em}

$x = x^{(0)}$ \quad
\vspace{0.5em}

Finché $\|Ax - b\|_{\infty} > \varepsilon$: \\[0.2em]
\hspace*{1.5em} Per ogni $i$ in $\{1, \dots, n\}$: \\[0.2em]
\hspace*{3em} $x_i^{(k+1)} = \frac{1}{A_{i,i}} \left( b_i - \sum_{j \neq i} A_{i,j} \, x_j^{(k)} \right)$

\vspace{0.6em}

\textsc{Output}: $x$ (approssimazione della soluzione del sistema $Ax = b$)
\end{minipage}
}
\end{flushleft}

\begin{lstlisting}
function x = Jacobi(A, b, x0, tol)
    % Metodo iterativo di Jacobi per risolvere Ax = b

    x = x0;                   % Inizializzo x con x^(0)
    n = length(b);            % Dimensione del sistema
    maxIter = 1000;           % Numero massimo di iterazioni
    for k = 1:maxIter
        x_new = zeros(n,1);
        for i = 1:n
            s = A(i,1:i-1) * x(1:i-1) + A(i,i+1:n) * x(i+1:n);  % Somma pesata esclusa la diagonale
            x_new(i) = (b(i) - s) / A(i,i);                     % Formula di Jacobi
        end
        if norm(A*x_new - b, inf) < tol
            x = x_new;
            return;
        end
        x = x_new;
    end
end
\end{lstlisting}
\newpage
\section*{Metodo di Gauss-Seidel}
\addcontentsline{toc}{section}{Metodo di Gauss-Seidel}

\begin{flushleft}
\fbox{
\begin{minipage}{\linewidth}
\textsc{Input}: $A \in \mathbb{R}^{n \times n}, \, b \in \mathbb{R}^n, \, x^{(0)} \in \mathbb{R}^n$, $\varepsilon > 0$

\vspace{0.6em}

Scriviamo $A = D + L + U$ con:
\begin{itemize}
  \item[] $D =$ matrice diagonale $\Rightarrow D_{i,j} = \delta_{i,j} \cdot A_{i,j}$
  \item[] $L =$ parte strettamente inferiore ($i > j$)
  \item[] $U =$ parte strettamente superiore ($i < j$)
\end{itemize}

\vspace{0.6em}

$Ax = b \iff (D + L + U)x = b$

\vspace{0.3em}

Ora porto $Ux$ dall'altra parte, quindi risolvo iterativamente 

$$(D + L)x^{(k+1)} = b - Ux^{(k)}$$

\noindent
Che equivale a:
$$x^{(k+1)} = (D + L)^{-1}(b - Ux^{(k)})$$

\noindent
Componente per componente ($i$-esima equazione di $Ax = b$):
$$x_i^{(k+1)} = \frac{1}{A_{i,i}} \left( b_i - \sum_{j=1}^{i-1} A_{i,j} x_j^{(k+1)} - \sum_{j=i+1}^{n} A_{i,j} x_j^{(k)} \right)$$

\vspace{0.6em}

\textsc{Output}: $x$ (approssimazione della soluzione del sistema $Ax = b$)

\end{minipage}
}
\end{flushleft}

\vspace{1em}
\begin{lstlisting}
function x = GaussSeidel(A, b, x0, tol)
    % Metodo iterativo di Gauss-Seidel per risolvere Ax = b

    x = x0;                   % Inizializzo con la stima iniziale x^(0)
    n = length(b);            % Dimensione del sistema
    maxIter = 1000;           % Numero massimo di iterazioni
    
    for k = 1:maxIter
        x_old = x;            % Salvo la soluzione precedente
        for i = 1:n
            % Calcolo:
            % x_i^{(k+1)} = (1/A_ii) * (b_i - somma_{j<i} A_ij x_j^(k+1) - somma_{j>i} A_ij x_j^(k))
            s1 = A(i,1:i-1) * x(1:i-1);       % somma A_ij * x_j^(k+1) per j < i
            s2 = A(i,i+1:n) * x_old(i+1:n);   % somma A_ij * x_j^(k) per j > i
            x(i) = (b(i) - s1 - s2) / A(i,i); % nuova x_i
        end
        
        % Verifica se il residuo è sufficientemente piccolo
        if norm(A*x - b, inf) < tol
            return;
        end
    end
end
\end{lstlisting}
\newpage
\subsection*{Metodo di Gauss-Seidel Simmetrico}
\addcontentsline{toc}{subsection}{Metodo di Gauss-Seidel Simmetrico}

\begin{flushleft}
\fbox{
\begin{minipage}{\linewidth}
\textsc{Input}: $A \in \mathbb{R}^{n \times n}, \, b \in \mathbb{R}^n, \, x^{(0)} \in \mathbb{R}^n$, $\varepsilon > 0$

\vspace{0.6em}

Il metodo di Gauss-Seidel simmetrico è un metodo iterativo a due fasi, utile per sistemi $Ax = b$ con $A$ simmetrica definita positiva.

\textbf{Fase 1 (down sweep)}: si scrive
\[
A = E_1 + F_1,\quad \text{con } E_1 = \{A_{ij} \mid i \ge j\},\quad F_1 = \{A_{ij} \mid i < j\}
\]
e si risolve:
\[
E_1 x^{(k+1/2)} = b - F_1 x^{(k)} \text{ per }x^{(k+1/2)}.
\]

\textbf{Fase 2 (up sweep)}: si scrive
\[
A = E_2 + F_2,\quad \text{con } E_2 = \{A_{ij} \mid i > j\},\quad F_2 = \{A_{ij} \mid i \le j\}
\]
e si risolve:
\[
F_2 x^{(k+1)} = b - E_2 x^{(k+1/2)}\text{ per }x^{(k+1)}.
\]

\textbf{Condizione di arresto}:
\[
\|x^{(k+1)} - x^{(k)}\| < \varepsilon
\]

\textsc{Output}: $x$ (approssimazione della soluzione)

\end{minipage}
}
\end{flushleft}

\vspace{0.6em}

\textbf{Differenze rispetto al Gauss-Seidel classico}:
\begin{itemize}
    \item Il metodo classico aggiorna $x^{(k+1)}$ con una sola fase: $(L + D)x^{(k+1)} = b - Ux^{(k)}$
    \item Il metodo simmetrico esegue due risoluzioni per iterazione, rispettivamente con $E_1$ e $F_2$
    \item È simmetrico: usa sia parte inferiore che superiore della matrice
    \item Più adatto a matrici simmetriche definite positive
\end{itemize}

\newpage
\vspace{1em}
\begin{lstlisting}
function [x, k] = GaussSeidelSimmetrico(A, b, x0, tol)
    % Metodo iterativo Gauss-Seidel simmetrico per risolvere Ax = b

    x = x0;
    n = length(b);
    maxIter = 1000;

    % Decomposizione E1 (i >= j), F1 (i < j)
    E1 = tril(A);
    F1 = A - E1;

    % Decomposizione E2 (i > j), F2 (i <= j)
    E2 = tril(A, -1);
    F2 = A - E2;

    for k = 1:maxIter
        % Fase 1
        y = b - F1 * x;
        x_half = E1 \ y;

        % Fase 2
        z = b - E2 * x_half;
        x_new = F2 \ z;

        if norm(x_new - x, inf) < tol
            x = x_new;
            return;
        end

        x = x_new;
    end
end
\end{lstlisting}

\newpage
\section*{Metodo di Richardson}
\addcontentsline{toc}{section}{Metodo di Richardson}

\begin{flushleft}
\fbox{
\begin{minipage}{\linewidth}
\textsc{Input}: $A \in \mathbb{R}^{n \times n}$, $b \in \mathbb{R}^n$, $x_0 \in \mathbb{R}^n$, $\varepsilon > 0$, $\alpha > 0$

\item $r = b - A x_0$
\item $\text{residui} = [\|r\|_2]$
\item Finch\'e $\|r\|_2 > \varepsilon$:
\begin{itemize}
\item $x = x + \alpha r$
\item $r = b - A x$
\item $\text{residui} = [\text{residui}, \|r\|_2]$
\end{itemize}

\textsc{Output}: $x$ (approssimazione della soluzione del sistema $Ax = b$)
\end{minipage}
}
\end{flushleft}

\vspace{1em}
\begin{lstlisting}
function [x, residuals] = richardson(A, b, x0, epsilon, alpha)
    % Metodo iterativo di Richardson per risolvere Ax = b

    x = x0;
    r = b - A*x;
    residuals = [norm(r)];
    max_iter = 1000; % Numero massimo iterazioni (come default)
    
    while norm(r) > epsilon && length(residuals) <= max_iter
        x = x + alpha*r;
        r = b - A*x;
        residuals = [residuals; norm(r)];
        % Aggiunge la norma del nuovo residuo in fondo al vettore
    end
end
\end{lstlisting}

\begin{lstlisting}
% Esempio di utilizzo
n = 6;
A = full(gallery('poisson', n));
b = ones(n, 1);
x0 = zeros(n, 1);
epsilon = 1e-6;
alpha = 0.5; % Parametro da scegliere opportunamente
[x, res] = richardson(A, b, x0, epsilon, alpha);
figure;
semilogy(0:length(res)-1, res, '-o', 'LineWidth', 1.5);
xlabel('Iterazione');
ylabel('||r||_2');
title('Convergenza del Metodo di Richardson');
grid on;
\end{lstlisting}


\newpage
\section*{Metodi di tipo gradiente}
\addcontentsline{toc}{section}{Metodi di tipo gradiente}

Vogliamo risolvere il sistema lineare \( Ax = b \). Una strategia consiste nel minimizzare la funzione
$\varphi(x) = \frac{1}{2} x^\top A x - b^\top x$
che ha minimo unico in \( x^* = A^{-1}b \) quando \( A = A^\top \succ 0 \).  
Il gradiente è \( \nabla \varphi(x_k) = A x_k - b = -r_k \), quindi il residuo \( r_k = b - A x_k \) è la direzione di massima discesa per \( \varphi(x) \).

I metodi iterativi aggiornano la soluzione con
\[
x_{k+1} = x_k + \alpha_k r_k \quad (p_k \text{ per il gradiente coniugato}),
\]
dove \( \alpha_k \) è scelto in modi diversi.

\paragraph{Ortomin}
 valido per \( A \) non indefinita ($x^TAx\ne0\quad\forall x \in \mathbb{R}$). Minimizza il residuo \( \|r_{k+1}\|_2^2 \) scegliendo un \( \alpha_k \) tale che il nuovo residuo sia ortogonale a \( A r_k \).
\begin{itemize}
    \item \textbf{Aggiornamento}: \( x_{k+1} = x_k + \alpha_k r_k \)
    \item \textbf{Condizione di ortogonalità}: \( r_{k+1} \perp A r_k \)
    \item \textbf{Funzionale minimizzato}: \( \|r_{k+1}\|_2^2 \)
    \item \textbf{Spazio di ricerca}: \( r_k + \text{span}(A r_k) \)
    \item \textbf{Passo ottimale}: 
    \[
    \alpha_k = \frac{r_k^\top A r_k}{(A r_k)^\top A r_k}
    \]
\end{itemize}

\paragraph{Metodo del gradiente}
 valido solo per \( A = A^\top \succ 0 \). Minimizza l’errore in norma dell’energia: \( \|e_{k+1}\|_A^2 \), con \( e_k = x^* - x_k \), scegliendo \( \alpha_k \) tale che \( e_{k+1} \perp_A r_k \).
\begin{itemize}
    \item \textbf{Aggiornamento}: \( x_{k+1} = x_k + \alpha_k r_k \)
    \item \textbf{Condizione di ortogonalità}: \( e_{k+1} \perp_A r_k \)
    \item \textbf{Funzionale minimizzato}: \( \|e_{k+1}\|_A^2 \)
    \item \textbf{Spazio di ricerca}: \( e_k + \text{span}(r_k) \)
    \item \textbf{Passo ottimale}:
    \[
    \alpha_k = \frac{r_k^\top r_k}{r_k^\top A r_k}
    \]
\end{itemize}

\paragraph{Gradiente coniugato}
 come il gradiente, ma le direzioni sono rese tra loro \( A \)-ortogonali (coniugate), per accelerare la convergenza.
\begin{itemize}
    \item \textbf{Aggiornamento}: \( x_{k+1} = x_k + \alpha_k p_k \), con \( p_k = r_k + \beta_k p_{k-1} \text{ (alla $1^a$ iterazione }p_k=r_k)\)
    \item \textbf{Condizione di ortogonalità}: \( e_{k+1} \perp_A p_k \), con \( p_i^\top A p_j = 0 \) per \( i \ne j \)
    \item \textbf{Funzionale minimizzato}: \( \|e_{k+1}\|_A^2 \)
    \item \textbf{Spazio di ricerca}: \( e_k + \text{span}(p_k, p_{k-1}) \)
    \item \textbf{Passo ottimale}:
    \[
    \alpha_k = \frac{r_k^\top r_k}{p_k^\top A p_k}, \quad \beta_k = \frac{r_k^\top r_k}{r_{k-1}^\top r_{k-1}}
    \]
\end{itemize}


\newpage
\subsection*{Metodo Ortomin}
\addcontentsline{toc}{subsection}{Metodo Ortomin}

\begin{flushleft}
\fbox{
\begin{minipage}{\linewidth}
\textsc{Input}: $A \in \mathbb{R}^{n \times n}$ (simmetrica def. pos.), $b \in \mathbb{R}^n$, $x_0 = \mathbf{1}$, tol = $10^{-10}\|b\|_2$, $it_{\text{max}} = 1000$

\item Calcola $r = b - A x_0$, $\text{residui} = [\|r\|_2]$
\item Per $k = 1$ a $it_{\text{max}}$:
\begin{itemize}
\item $Ar = A \cdot r$
\item $\alpha = \frac{r^T Ar}{(Ar)^T Ar}$ \quad (Passo ottimale)
\item $x = x + \alpha r$
\item $r = r - \alpha Ar$
\item $\text{residui} = [\text{residui}, \|r\|_2]$
\item Se $\|r\|_2 \leq $ tol: termina
\end{itemize}

\textsc{Output}: $x$, $\text{residui}$
\end{minipage}
}
\end{flushleft}

\vspace{1em}
\begin{lstlisting}[language=Matlab]
function [x, residuals] = ortomin(A, b, x0, tol, itmax)

x = x0;
r = b - A*x;
residuals = zeros(itmax+1, 1);
residuals(1) = norm(r);
it = 0;

while residuals(it+1) > tol && it < itmax
    Ar = A * r;
    alpha = (r' * Ar) / (Ar' * Ar);
    x = x + alpha * r;
    r_new = r - alpha * Ar;

    it = it + 1;
    residuals(it+1) = norm(r_new);
    r = r_new;
end

residuals = residuals(1:it+1);

% Plot risultati
figure;
plot(0:it, residuals, '-o', 'LineWidth', 1.5);
xlabel('Iterazione');
ylabel('||A x - b||_2');
title('Convergenza del Metodo ORTOMIN');
grid on;
end

\end{lstlisting}


\newpage
\subsection*{Metodo del Gradiente}
\addcontentsline{toc}{subsection}{Metodo del Gradiente}

\begin{flushleft}
\fbox{
\begin{minipage}{\linewidth}
\textsc{Input}: $A \in \mathbb{R}^{n \times n}$ (simmetrica def. pos.), $b \in \mathbb{R}^n$, $x_0 = \mathbf{1}$, tol = $10^{-6}\|b\|_2$, $it_{\text{max}} = 1000$

\item Calcola $r = b - A x_0$, $\text{residui} = [\|r\|_2]$
\item Per $k = 1$ a $it_{\text{max}}$:
\begin{itemize}
\item $Ar = A \cdot r$
\item $\alpha = \frac{r^T r}{r^T A r}$ \quad (Passo ottimale)
\item $x = x + \alpha r$
\item $r = b - A x$
\item $\text{residui} = [\text{residui}, \|r\|_2]$
\item Se $\|r\|_2 \leq $ tol: termina
\end{itemize}

\textsc{Output}: $x$, $\text{residui}$
\end{minipage}
}
\end{flushleft}

\vspace{1em}
\begin{lstlisting}[language=Matlab]
function [x, residuals] = gradiente(A, b, x0, tol, itmax)

x = x0;
r = b - A*x;
residuals = zeros(itmax+1, 1);
residuals(1) = norm(A*x - b);
it = 0;

while residuals(it+1) > tol && it < itmax
    Ar = A * r;
    alpha_opt = (r' * r) / (r' * Ar);
    x = x + alpha_opt * r;
    r = b - A * x;
    it = it + 1;
    residuals(it+1) = norm(r);
end

residuals = residuals(1:it+1);

% Plot risultati
figure;
plot(0:it, residuals, '-o', 'LineWidth', 1.5);
xlabel('Iterazione');
ylabel('||A x - b||_2');
title('Convergenza del Metodo del Gradiente');
grid on;
end

\end{lstlisting}

\newpage
\subsection*{Metodo del Gradiente Coniugato}
\addcontentsline{toc}{subsection}{Metodo del Gradiente Coniugato}

\begin{flushleft}
\fbox{
\begin{minipage}{\linewidth}
\textsc{Input}: $A \in \mathbb{R}^{n \times n}$ (simmetrica def. pos.), $b \in \mathbb{R}^n$, $x_0 = \mathbf{1}$, tol = $10^{-10}\|b\|_2$, $it_{\text{max}} = 1000$

\item Calcola $r = b - A x_0$, $p = r$, $\text{residui} = [\|r\|_2]$
\item Per $k = 1$ a $it_{\text{max}}$:
\begin{itemize}
\item $Ap = A \cdot p$
\item $\alpha = \frac{r^T p}{p^T A p}$ \quad (Passo ottimale)
\item $x = x + \alpha p$
\item $r' = r - \alpha Ap$
\item $\beta = \frac{r'^T r'}{r^T r}$
\item $p = r' + \beta p$
\item $r=r'$
\item $\text{residui} = [\text{residui}, \|r\|_2]$
\item Se $\|r\|_2 \leq $ tol: termina
\end{itemize}

\textsc{Output}: $x$, $\text{residui}$
\end{minipage}
}
\end{flushleft}

\vspace{1em}
\begin{lstlisting}
function [x, residuals] = gradiente_coniugato(A, b, x0, tol, itmax)

x = x0;
r = b - A*x;
p = r;
residuals = zeros(itmax+1, 1);
residuals(1) = norm(r);
it = 0;
while residuals(it+1) > tol && it < itmax
    Ap = A * p;
    alpha = (r' * p) / (p' * Ap);
    x = x + alpha * p;
    r_new = r - alpha * Ap;
    beta = (r_new' * r_new) / (r' * r);
    p = r_new + beta * p;
    it = it + 1;
    residuals(it+1) = norm(r_new);
    r = r_new;
end

residuals = residuals(1:it+1);

% Plot risultati
figure;
plot(0:it, residuals, '-o', 'LineWidth', 1.5);
xlabel('Iterazione');
ylabel('||A x - b||_2');
title('Convergenza del Metodo del Gradiente Coniugato');
grid on;
end


\end{lstlisting}
\section*{Equazioni e sistemi non lineari}

Dato \( f : \mathbb{R}^n \to \mathbb{R}^n \), vogliamo trovare \( x \in \mathbb{R}^n \) tale che:
\[
f(x) = 0
\]
  
\[
f(x) = \begin{pmatrix}
f_1(x_1, \dots, x_n) \\
\vdots \\
f_n(x_1, \dots, x_n)
\end{pmatrix}, \quad f(x) = 0 \iff
\begin{cases}
f_1(x) = 0 \\
\vdots \\
f_n(x) = 0
\end{cases}
\]

\textbf{Obiettivo:} trovare \( x^* \in \mathbb{R}^n \) tale che \( f(x^*) = 0 \).

\textbf{Idea:}  
Costruire una successione \( (x_k) \) tale che:
\[
x_{k+1} = g(x_k), \quad \text{con } \lim_{k \to \infty} x_k = x^*, \text{ e } f(x^*) = 0
\]

\textbf{Metodi principali:} Bisezione, Secanti (\( n = 1 \)), Newton, Punto fisso (\( n \in \mathbb{N}^*=\mathbb{N}\setminus\{0\}\))

\section*{Metodo di Bisezione}
\addcontentsline{toc}{section}{Metodo di Bisezione}

\begin{flushleft}
\fbox{
\begin{minipage}{\linewidth}
\textsc{Input}: funzione $f$, intervallo $[a,b]$ tale che $f(a) \cdot f(b) < 0$, tolleranza $\text{tol}$, iterazioni massime $it_{\max}$

\item Per $k = 1$ a $it_{\max}$:
\begin{itemize}
\item $c = \frac{a + b}{2}$ \quad (punto medio)
\item Se $f(c) = 0$ oppure $\frac{b - a}{2} < \text{tol}$: termina e restituisce $c$
\item Se $f(a) \cdot f(c) < 0$ allora $b = c$, altrimenti $a = c$
\end{itemize}

\textsc{Output}: approssimazione della radice $c$
\end{minipage}
}
\end{flushleft}

\vspace{1em}
\begin{lstlisting}[language=Matlab]
function root = bisezione(f, a, b, tol, itmax)
    if f(a)*f(b) >= 0
        error('La funzione deve avere segni opposti in a e b')
    end
    
    for k = 1:itmax
        c = (a + b) / 2; % punto medio
        if f(c) == 0 || (b - a)/2 < tol
            root = c;
            return
        end
        if f(a)*f(c) < 0
            b = c;
        else
            a = c;
        end
    end
    root = (a + b) / 2;
end
\end{lstlisting}

\newpage
\subsection*{Stima dell'errore nel metodo di Bisezione}
\addcontentsline{toc}{subsection}{Stima dell'errore}
Sia \(f:[a,b] \to \mathbb{R}\) continua, con \(f(a)f(b) < 0\).  
Allora esiste \(x^* \in [a,b]\) tale che \(f(x^*) = 0\).

Il metodo di bisezione costruisce una successione \(\{x_n\}\) di approssimazioni di \(x^*\), con:
\[
x_n = \frac{a_n + b_n}{2},\quad \text{e } x^* \in [a_n, b_n]
\]

Ad ogni iterazione:
\[
[a_{n+1}, b_{n+1}] \subset [a_n, b_n],\quad \text{con } b_{n+1} - a_{n+1} = \frac{1}{2}(b_n - a_n)
\]

\textbf{Errore assoluto all’iterazione \(n\):}
\[
|x_n - x^*| \leq \frac{b_n - a_n}{2} = \frac{b - a}{2^{n+1}}
\]
\[
|x_{n+1} - x^*| \leq \frac{1}{2} |x_n - x^*|
\]

\textbf{Quindi:} l'errore si dimezza ad ogni iterazione  
\(\Rightarrow\) \textbf{convergenza lineare}
\[
|e_{n+1}| \leq \rho |e_n|,\quad \text{con } \rho = \frac{1}{2} < 1
\]


\newpage
\section*{Metodo di Newton}
\addcontentsline{toc}{section}{Metodo di Newton}

\begin{flushleft}
\fbox{
\begin{minipage}{\linewidth}
\textsc{Input}: funzione $f$, derivata $f'$, punto iniziale $x_0$, tolleranza $\text{tol}$, iterazioni massime $it_{\max}$

\item Per $k = 1$ a $it_{\max}$:
\begin{itemize}
\item $x_{k+1} = x_k - \frac{f(x_k)}{f'(x_k)}$
\item Se $|x_{k+1} - x_k| < \text{tol}$: termina
\end{itemize}

\textsc{Output}: approssimazione della radice $x_{k+1}$
\end{minipage}
}
\end{flushleft}

\vspace{1em}
\begin{lstlisting}[language=Matlab]
function root = newton(f, df, x0, tol, itmax)
    for k = 1:itmax
        fx = f(x0);
        dfx = df(x0);
        
        if dfx == 0
            error('Derivata nulla: il metodo fallisce')
        end
        
        x1 = x0 - fx / dfx;
        
        if abs(x1 - x0) < tol
            root = x1;
            return
        end
        
        x0 = x1;
    end
    root = x0
end
\end{lstlisting}

\newpage
\subsection*{Metodo di Newton in \(\mathbb{R}^n\)}
\addcontentsline{toc}{subsection}{Metodo di Newton in \(\mathbb{R}^n\)}

\begin{flushleft}
\fbox{
\begin{minipage}{\linewidth}
\textsc{Input}: funzione vettoriale \(F: \mathbb{R}^n \to \mathbb{R}^n\), Jacobiano \(J_F\), punto iniziale \(x_0 \in \mathbb{R}^n\), tolleranza \(\text{tol}\), iterazioni massime \(it_{\max}\)

\begin{itemize}
\item Per \(k = 0, \dots, it_{\max}-1\):
\begin{itemize}
\item risolvi \(J_F(x_k) h = -F(x_k)\)
\item aggiorna \(x_{k+1} = x_k + h\)
\item se \(\|h\| < \text{tol}\), termina
\end{itemize}
\end{itemize}

\textsc{Output}: approssimazione della radice \(x_{k+1}\)
\end{minipage}
}
\end{flushleft}

\vspace{1em}
\textbf{Spiegazione:}  
Lo sviluppo di Taylor per \(F: \mathbb{R}^n \to \mathbb{R}^n\) attorno a \(x_k\) è:
\[
F(x) = \sum_{j=0}^{\infty} \frac{1}{j!} D^jF(x_k)[(x - x_k)^{\otimes j}]
\]
dove:
- \( D^jF(x_k) \) è il tensore delle derivate parziali di ordine \(j\),
- \( (x - x_k)^{\otimes j} \) è il prodotto tensoriale di ordine \(j\) di \( (x - x_k) \) con sé stesso.

Troncando al primo ordine:
\[
F(x) = F(x_k) + J_F(x_k)(x - x_k) + o(\|x - x_k\|)
\]
\[
F(x) = 0 \Rightarrow F(x_k) + J_F(x_k)(x - x_k) + \underbrace{o(\|x - x_k\|)}_{\rightarrow 0} = 0
\]
\[
J_F(x_k)(x - x_k) = -F(x_k)
\]
\[
h := x - x_k \Rightarrow J_F(x_k)h = -F(x_k)
\]

\( n = 1 \Rightarrow J_F(x_k) = f'(x_k) \Rightarrow h = -\frac{f(x_k)}{f'(x_k)}  \Rightarrow x_{k+1} = x_k - \frac{f(x_k)}{f'(x_k)} \) \{Newton in \(\mathbb{R}\)\}

\vspace{1em}


\begin{lstlisting}[language=Matlab]
function root = newtonRn(F, JF, x0, tol, itmax)
    x = x0;
    for k = 1:itmax
        Fx = F(x);
        Jx = JF(x);
        h = -Jx \ Fx;
        x_new = x + h;
        if norm(h) < tol
            root = x_new;
            return
        end
        x = x_new;
    end
    root = x;
end
\end{lstlisting}

\newpage
\subsection*{Stima dell'errore nel metodo di Newton}
\addcontentsline{toc}{subsection}{Stima dell'errore}
\textbf{Definizione:}
Sia \( f : \mathbb{R} \to \mathbb{R} \) una funzione derivabile. Si dice che \( x^* \in \mathbb{R} \) è una \emph{radice semplice} di \( f \) se:
\[
f(x^*) = 0 \quad \text{e} \quad f'(x^*) \ne 0
\]

Sia \( f:\mathbb{R} \to \mathbb{R} \) una funzione due volte derivabile, e \( x^* \) una radice semplice semplice,

Il metodo di Newton genera la successione:
\[
x_{n+1} = x_n - \frac{f(x_n)}{f'(x_n)}
\]

Supponiamo che \( x_n \) sia vicino a \( x^* \), definiamo l’errore \( e_n = x_n - x^* \).  
Sviluppiamo \( f(x_n) \) con Taylor attorno a \( x^* \):

\[
f(x_n) = f(x^* + e_n) = f(x^*) + f'(x^*)e_n + \frac{f''(\xi_n)}{2} e_n^2 = f'(x^*)e_n + \frac{f''(\xi_n)}{2} e_n^2\text{ con }\xi_n \in (x^*, x_n)\]

\[
x_{n+1} = x_n - \frac{f'(x^*)e_n + \frac{f''(\xi_n)}{2} e_n^2}{f'(x_n)} \approx x_n - \frac{f'(x^*)e_n}{f'(x^*)} - \frac{f''(\xi_n)}{2f'(x^*)} e_n^2
\]
\[
\Rightarrow x_{n+1} - x^* = e_{n+1} \approx \cancel{e_n} - \cancel{e_n} - \frac{f''(\xi_n)}{2f'(x^*)} e_n^2 = -\frac{f''(\xi_n)}{2f'(x^*)} e_n^2 \Rightarrow
|e_{n+1}| \approx \left|\frac{f''(\xi_n)}{2f'(x^*)}\right| \cdot e_n^2
\]

Quindi Newton quando converge lo fa quadraticamente:
\[
e_{n+1} \approx C e_n^2 \quad \text{con} \quad C = \left|\frac{f''(\xi_n)}{2f'(x^*)}\right|
\]

\newpage

\section*{Metodo delle Secanti}


\addcontentsline{toc}{section}{Metodo delle Secanti}

\begin{flushleft}
\fbox{
\begin{minipage}{\linewidth}
\textsc{Input}: funzione $f$, due punti iniziali $x_{-1}, x_0$, tolleranza $\text{tol}$, iterazioni massime $it_{\max}$

\item Per $k = 1$ a $it_{\max}$:
\begin{itemize}
\item $x_{k+1} = x_k - f(x_k) \cdot \frac{x_k - x_{k-1}}{f(x_k) - f(x_{k-1})}$
\item Se $|x_{k+1} - x_k| < \text{tol}$: termina
\end{itemize}

\textsc{Output}: approssimazione della radice $x_{k+1}$
\end{minipage}
}
\end{flushleft}

\vspace{1em}
\begin{lstlisting}[language=Matlab]
function root = secanti(f, x0, x1, tol, itmax)
    for k = 1:itmax
        fx0 = f(x0);
        fx1 = f(x1);
        
        if fx1 == fx0
            error('Divisione per zero: metodo fallisce')
        end
        
        x2 = x1 - fx1 * (x1 - x0) / (fx1 - fx0);
        
        if abs(x2 - x1) < tol
            root = x2;
            return
        end
        
        x0 = x1;
        x1 = x2;
    end
    root = x1;
end
\end{lstlisting}

\newpage
\newpage
\subsection*{Stima dell'errore nel metodo delle Secanti}
\addcontentsline{toc}{subsection}{Stima dell'errore}
Sia \(x^*\) radice semplice di \(f\), con \(f(x^*)=0\), \(f'(x^*) \neq 0\), l'errore all'n-esima iterazione è \(e_n = x_n - x^*\) $\Rightarrow$ \(x_n=e_n+x^*\)

\[
x_{n+1} = x_n - f(x_n) \frac{x_n - x_{n-1}}{f(x_n) - f(x_{n-1})}
\]

Sviluppando \(f(x_n), f(x_{n-1})\) con Taylor, intorno a \(x^*\):
\[
f(x_n) = \underbrace{f(x^*)}_{0}+f'(x^*) e_n + \frac{f''(x^*)}{2} e_n^2 + o(e_n^2),\quad f(x_{n-1}) = \underbrace{f(x^*)}_{0} + f'(x^*) e_{n-1} + \frac{f''(x^*)}{2} e_{n-1}^2+ o(e_{n-1}^2)
\]

\[
x_{n+1} = \underbrace{x_n}_{x^*+e_n} - \underbrace{f(x_n)}_{f'(x^*) e_n + \frac{f''(x^*)}{2} e_n^2 + o(e_n)} \frac{\overbrace{x_n - x_{n-1}}^{e_n+x^*-(e_{n-1}+x^*)=e_n-e_{n-1}}}{f'(x^*)(e_n-e_{n-1})+\frac{1}{2}f''(x^*)(e_n^2-e_{n-1}^2)+o(e_n^2)-o(e_{n-1}^2)}
\]

\[x_{n+1} \approx x^* + e_n - \left( f'(x^*) e_n + \frac{f''(x^*)}{2} e_n^2 \right) \cdot \frac{\cancel{e_n - e_{n-1}}}{\cancel{(e_n - e_{n-1})}[f'(x^*) + \frac{f''(x^*)}{2}(e_n + e_{n-1})]}\]
\[= x^* + e_n - \left( f'(x^*) e_n + \frac{f''(x^*)}{2} e_n^2 \right) \cdot \frac{1}{1 + \left( \frac{\frac{f''(x^*)}{2}(e_n + e_{n-1})}{f'(x^*)} \right)}\cdot\frac{1}{f'(x^*)}
\]
\[= x^* + e_n - \frac{ f'(x^*) e_n + \frac{f''(x^*)}{2} e_n^2}{f'(x^*)} \cdot \frac{1}{1 + \left( \frac{\frac{f''(x^*)}{2}(e_n + e_{n-1})}{f'(x^*)} \right)}\]
\[\underbrace{\approx}_{\frac{1}{1+t}\underbrace{\sim}_{t\to0}{1-t}} x^* + e_n - \frac{ f'(x^*) e_n + \frac{f''(x^*)}{2} e_n^2}{f'(x^*)} \cdot \left[
1 - \left( \frac{\frac{f''(x^*)}{2}(e_n + e_{n-1})}{f'(x^*)} \right)\right]\]
\[= x^*+e_n-\left[\frac{\cancel{f'(x^*)}e_n}{\cancel{f'(x^*)}}+\frac{1}{2}\frac{f''(x^*)e_n^2}{f'(x^*)}-\left(\frac{\cancel{f'(x^*)}e_n}{\cancel{f'(x^*)}}+\frac{1}{2}\frac{f''(x^*)e_n^2}{f'(x^*)}\right)\left( \frac{\frac{f''(x^*)}{2}(e_n + e_{n-1})}{f'(x^*)} \right)\right]\]
\[= x^*+e_n-\left[e_n+\frac{1}{2}\frac{f''(x^*)e_n^2}{f'(x^*)}-\left(e_n(e_n + e_{n-1})\left( \frac{f''(x^*)}{2f'(x^*)} \right)+\underbrace{e_n^2(e_n + e_{n-1})\left(\frac{1}{2}\frac{f''(x^*)}{f'(x^*)}\right)\left( \frac{f''(x^*)}{2f'(x^*)}\right)}_{\textit{Questi termini sono trascurabili }\approx e_n^3} \right)\right]\]
\[= x^*+\cancel{e_n}-\left[\cancel{e_n}+\frac{1}{2}\frac{f''(x^*)e_n^2}{f'(x^*)}-\left(e_n(e_n + e_{n-1})\left( \frac{f''(x^*)}{2f'(x^*)} \right) \right)\right]\]

\[\Rightarrow e_{n+1}=x_{n+1}-x^*= -e_n^2\frac{1}{2}\frac{f''(x^*)}{f'(x^*)}+e_n(e_n + e_{n-1})\frac{1}{2}\frac{f''(x^*)}{f'(x^*)}=\frac{1}{2}\frac{f''(x^*)}{f'(x^*)}e_n(-e_n+e_n+e_{n-1})\]
$$=\underbrace{\frac{1}{2}\frac{f''(x^*)}{f'(x^*)}}_{C}e_ne_{n-1}$$

Ora $\begin{cases}e_{n+1}=Ce_n^{\alpha}\Rightarrow Ce_n=e_{n+1}^{\frac{1}{\alpha}}\Rightarrow e_{n-1}=e_n^{\frac{1}{\alpha}}\\e_{n+1}=Ce_ne_{n-1}=Ce_ne_n^{\frac{1}{\alpha}}=Ce_n^\alpha\end{cases}\Rightarrow \frac{1}{\alpha}+1=\alpha\Rightarrow\alpha=\frac{1+\sqrt{5}}{2}=\phi$
Quindi il metodo delle secanti converge con ordine \(\varphi\), più lentamente di Newton ma funziona senza bisogno di calcolare la derivata di f

\newpage
\section*{Metodo del Punto Fisso}
\addcontentsline{toc}{section}{Metodo del Punto Fisso}

\begin{flushleft}
\fbox{
\begin{minipage}{\linewidth}
\textsc{Input}: funzione $g$, punto iniziale $x_0$, tolleranza $\text{tol}$, iterazioni massime $it_{\max}$

\item Per $k = 1$ a $it_{\max}$:
\begin{itemize}
\item $x_{k+1} = g(x_k)$
\item Se $|x_{k+1} - x_k| < \text{tol}$: termina
\end{itemize}

\textsc{Output}: approssimazione del punto fisso $x_{k+1}$
\end{minipage}
}
\end{flushleft}

\vspace{1em}
\begin{lstlisting}[language=Matlab]
function root = punto_fisso(g, x0, tol, itmax)
    for k = 1:itmax
        x1 = g(x0);
        
        if abs(x1 - x0) < tol
            root = x1;
            return
        end
        
        x0 = x1;
    end
    root = x0;
end
\end{lstlisting}
\newpage
\subsection*{Stima dell'errore nel metodo del punto fisso}
\addcontentsline{toc}{subsection}{Stima dell'errore}
Sia \(x^*\) punto fisso di \(g\), cioè \(g(x^*) = x^*\). Definiamo l'errore \(e_k = |x_k - x^*|\).

Il metodo è:
\[
x_{k+1} = g(x_k)
\]

Quindi:
\[
e_{k+1} = |x_{k+1} - x^*| = |g(x_k) - g(x^*)| =g(
\]

\[
e_{k+1} \le M e_k
\]

Perciò:
\[
e_k \le M^k e_0
\]

Il fatto che \(M < 1\) garantisce che \(M^k \to 0\) per \(k \to \infty\), cioè l'errore tende a zero e il metodo converge linearmente.

\textbf{Motivazione:} ad ogni iterazione l'errore si riduce almeno di un fattore \(M < 1\), quindi l’intervallo degli \(x_k\) si restringe e \(x_k \to x^*\).

\newpage

\section*{Matrici ortogonali}
\addcontentsline{toc}{subsection}{Matrici ortogonali}
Una matrice \( Q \in \mathbb{R}^{n \times n} \) si dice ortogonale se:
\[
Q^\top Q = I \quad \Longleftrightarrow \quad Q^\top = Q^{-1}
\]
\[
\quad \|Qx\|^2 = (Qx)^\top (Qx) = x^\top Q^\top Q x = x^\top x = \|x\|^2 \text{ 
 }\forall x \in \mathbb{R}^n
\]

\(\Rightarrow\) \( Q \) conserva il prodotto scalare e le norme: \textbf{è un'isometria lineare}.

In \( \mathbb{R}^2 \), ogni matrice ortogonale \( Q \) con \( \det(Q) = 1 \) rappresenta una rotazione di un angolo \( \theta \) attorno all'origine.

Base canonica:
\[
e_1 = \begin{pmatrix}1 \\ 0\end{pmatrix}, \quad e_2 = \begin{pmatrix}0 \\ 1\end{pmatrix}
\]

Effetto di una rotazione:
\[
Q e_1 = \begin{pmatrix} \cos(\theta) \\ \sin(\theta) \end{pmatrix}, \quad 
Q e_2 = \begin{pmatrix} \cos(\theta + \pi) = -\sin(\theta) \\ \sin(\theta + \pi)= \cos(\theta) \end{pmatrix}
\]

Quindi:
\[
Q = [Q e_1 \mid Q e_2] = 
\begin{pmatrix}
\cos(\theta) & -\sin(\theta) \\
\sin(\theta) & \cos(\theta)
\end{pmatrix}
\]

E si verifica che:
\[
Q^\top Q = 
\begin{pmatrix}
\cos(\theta) & \sin(\theta) \\
-\sin(\theta) & \cos(\theta)
\end{pmatrix}
\begin{pmatrix}
\cos(\theta) & -\sin(\theta) \\
\sin(\theta) & \cos(\theta)
\end{pmatrix}
=
\begin{pmatrix}
\cos^2(\theta) + \sin^2(\theta) & 0 \\
0 & \cos^2(\theta) + \sin^2(\theta)
\end{pmatrix}
= I
\]

\section*{Fattorizzazione QR con Gram-Schmidt}
\addcontentsline{toc}{section}{Fattorizzazione QR}

\begin{flushleft}
\fbox{
\begin{minipage}{\linewidth}
\textsc{Input}: $A \in \mathbb{R}^{m \times n}$ \quad ($m \geq n$)

\vspace{0.6em}

Per ogni $j$ in $\{1, \dots, n\}$:

\vspace{0.4em}

\hspace*{1.5em} $q_j = a_j - \sum_{k=1}^{j-1} (q_k, a_j) \, q_k$ \\[0.2em]
\hspace*{1.5em} $r_{k,j} = (q_k, a_j)$ \quad per $k = 1, \dots, j-1$ \\[0.2em]
\hspace*{1.5em} $r_{j,j} = \|q_j\|_2$ \\[0.2em]
\hspace*{1.5em} Se $r_{j,j} \neq 0$, allora $q_j = q_j / r_{j,j}$

\vspace{0.6em}

Se $m > \text{rank}(Q(:,1:n))$, completa la base ortonormale:

\vspace{0.4em}

\hspace*{1.5em} Per ogni $j$ in $\{\text{rank}(Q)+1, \dots, m\}$:\\[0.2em]
\hspace*{3em} $v = \text{random}(m)$ \\[0.2em]
\hspace*{3em} $v = v - \sum_{k=1}^{j-1} (q_k, v) \, q_k$ \\[0.2em]
\hspace*{3em} $q_j = v / \|v\|_2$

\vspace{0.6em}

\textsc{Output}: $Q$ ortonormale, $R$ triangolare superiore \quad ($A = QR$)
\end{minipage}
}
\end{flushleft}



\vspace{1em}
\begin{lstlisting}
function [Q, R] = qr_completa(A)
    [m, n] = size(A);                  % Dimensioni di A (m righe, n colonne)
    Q = zeros(m, m);                   % Inizializzo Q come matrice m * m (ortonormale)
    R = zeros(m, n);                   % Inizializzo R come matrice m * n (triangolare superiore)
    for j = 1:n
        Q(:,j) = A(:,j) - Q(:,1:j-1) * (Q(:,1:j-1)' * A(:,j));
        
        % Rende aj ortogonale rispetto alle altre j-1 colonne di Q
        % per poi salvarla come qj
        % (Q(:,1:j-1)' * A(:,j)) e' un vettore con (qi,aj) come i-esima componente
        % che viene moltiplicato per qi ovvero Q(:,1:j-1) e infine sottratto ad aj
        R(1:j-1, j) = Q(:,1:j-1)' * A(:,j);
        % Salva i valori (aj, qi) per i<j sopra la diagonale di R in modo che sia triangolare superiore
        R(j,j) = norm(Q(:,j));          
        % Valore diagonale: norma del vettore ortogonalizzato
        if R(j,j) ~= 0
            Q(:,j) = Q(:,j) / R(j,j);  
            % Normalizzazione di qj
        end
    end
    r = rank(Q(:,1:n)); % Calcolo il rango delle n colonne di Q per sapere se e da dove iniziare il completamento ortogonale (se le colonne di A fossero linearmente indipendenti o avessimo una matrice A rettangolare, con m>n)
    for j = r+1:m
        v = randn(m,1);                           % Vettore casuale
        v = v - Q(:,1:j-1) * (Q(:,1:j-1)' * v);   % Ortogonalizzazione rispetto alle j-1 colonne di Q
        Q(:,j) = v / norm(v);                     % Normalizzo la nuova colonna
    end
end
\end{lstlisting}

\newpage
\subsection*{Applicazioni della fattorizzazione QR}

Sia \( A \in \mathbb{R}^{m \times n} \) con \( m \geq n \). La fattorizzazione QR consiste nello scrivere
\[
A = QR \quad \text{con } Q \in \mathbb{R}^{m \times m} \text{ ortogonale } (Q^T Q = I_m),\ R \in \mathbb{R}^{m \times n} \text{ triangolare superiore}.
\]
Se \( A \) ha rango \( n \), allora si può anche usare la forma ridotta:
\[
A = \tilde{Q} R \quad \text{con } \tilde{Q} \in \mathbb{R}^{m \times n} \text{ con } \tilde{Q}^T \tilde{Q} = I_n,\ R \in \mathbb{R}^{n \times n} \text{ triangolare superiore}.
\]

\paragraph{1. Risoluzione di sistemi sovradeterminati \( Ax = b \) con \( m > n \):}

Vogliamo minimizzare \( \|Ax - b\|_2 \). Sostituendo \( A = QR \), si ha:
\[
\|Ax - b\|_2 = \|QRx - b\|_2 = \|Rx - Q^T b\|_2 \quad (\text{poiché } Q \text{ è ortogonale}).
\]
Sia \( Q^T b = \begin{bmatrix} c \\ d \end{bmatrix} \), con \( c \in \mathbb{R}^n \), \( d \in \mathbb{R}^{m-n} \), e sia \( R = \begin{bmatrix} R_1 \\ 0 \end{bmatrix} \) con \( R_1 \in \mathbb{R}^{n \times n} \). Allora
\[
\|Rx - Q^T b\|_2^2 = \|R_1 x - c\|_2^2 + \|d\|_2^2,
\]
minimo per \( x \) tale che \( R_1 x = c \). Essendo \( R_1 \) triangolare superiore, il sistema si risolve per sostituzione all'indietro.

\paragraph{2. Ortonormalizzazione (Gram-Schmidt numerico):}

Le colonne di \( Q \) formano una base ortonormale per \( span\{a_1, \dots, a_n\} \), dove \( a_i \) sono le colonne di \( A \). Si ottiene con processo di Gram-Schmidt:
\[
q_1 = \frac{a_1}{\|a_1\|},\quad
q_k = \frac{a_k - \sum_{j=1}^{k-1} (q_j^T a_k) q_j}{\left\| a_k - \sum_{j=1}^{k-1} (q_j^T a_k) q_j \right\|},\quad k = 2, \dots, n.
\]
I coefficienti \( r_{ij} = q_i^T a_j \) costituiscono \( R \), quindi si ottiene \( A = QR \).

\paragraph{3. Calcolo degli autovalori (metodo QR iterativo):}

Sia \( A \in \mathbb{R}^{n \times n} \), si costruisce la sequenza:
\[
A^{(0)} := A,\quad A^{(k)} = Q^{(k)} R^{(k)},\quad A^{(k+1)} := R^{(k)} Q^{(k)}.
\]
Ogni \( A^{(k+1)} = (Q^{(k)})^{-1} A^{(k)} Q^{(k)} \Rightarrow \) tutti gli \( A^{(k)} \) sono simili ad \( A \), quindi hanno gli stessi autovalori. Se \( A \) è simmetrica o diagonalizzabile con autovalori distinti, allora:
\[
\lim_{k \to \infty} A^{(k)} = T \quad \text{triangolare superiore con } T_{ii} = \lambda_i(A).
\]
Quindi si possono leggere gli autovalori sulla diagonale.

\textbf{Algoritmo con shift (per migliorare la convergenza):} dato uno shift \( \mu_k \in \mathbb{R} \), si applica:
\[
A^{(k)} - \mu_k I = Q^{(k)} R^{(k)} \Rightarrow A^{(k+1)} = R^{(k)} Q^{(k)} + \mu_k I.
\]
Scelta tipica: \( \mu_k = A^{(k)}_{nn} \). Converge più velocemente agli autovalori.

\paragraph{4. Risoluzione di sistemi lineari \( Ax = b \) con \( A \) quadrata e invertibile:}

Dato \( A = QR \), si ha:
\[
Ax = b \Rightarrow QRx = b \Rightarrow Rx = Q^T b.
\]
Essendo \( R \) triangolare superiore, si risolve facilmente per sostituzione all'indietro.

\paragraph{5. Calcolo del determinante (solo se \( A \in \mathbb{R}^{n \times n} \)):}

Se \( A = QR \), e \( Q \) ortogonale (quindi \( |\det(Q)| = 1 \)), allora:
\[
\det(A) = \det(Q)\det(R) = \pm \prod_{i=1}^n r_{ii}.
\]

\paragraph{6. Calcolo numerico della pseudoinversa \( A^\dagger \):}

Se \( A \in \mathbb{R}^{m \times n} \), rango \( n \), \( A = QR \) con \( Q \in \mathbb{R}^{m \times n} \), \( Q^T Q = I \), \( R \in \mathbb{R}^{n \times n} \), allora:
\[
A^\dagger = R^{-1} Q^T.
\]

\paragraph{7. Metodo di iterazione di potenza inversa con shift (per autovalori):}

Se si vuole l’autovalore \( \lambda \) più vicino a \( \mu \), si applica iterativamente:
\[
(A - \mu I) z^{(k)} = y^{(k)} \Rightarrow y^{(k+1)} = \frac{z^{(k)}}{\|z^{(k)}\|},
\]
risolvendo con fattorizzazione QR di \( A - \mu I \Rightarrow \) stabile e efficiente.

\paragraph{8. Compressione e SVD approssimata (quando SVD è troppo costosa):}

Per \( A \in \mathbb{R}^{m \times n} \) con \( m \gg n \), QR è più economica di SVD. Si può usare:
\[
A \approx QR \Rightarrow \text{usato in algoritmi di compressione, proiezione randomizzata, riduzione dimensionale}.
\]

\paragraph{9. Stabilità numerica:}

La QR è preferita alla decomposizione di Cholesky per problemi numerici malcondizionati, perché evita il calcolo di \( A^T A \), che può aumentare la condizione:
\[
cond(A^T A) = cond(A)^2 \gg cond(A).
\]

\paragraph{10. Algoritmi PCA (Principal Component Analysis):}

Dato \( X \in \mathbb{R}^{m \times n} \), centrare i dati e applicare QR:
\[
X = QR \Rightarrow \text{colonne di } Q \text{ sono componenti principali (non ordinate)}.
\]
Oppure si applica QR a \( X^T \) per trovare i vettori ortonormali su cui proiettare i dati.




\newpage
\textbf{Esempio: Calcolo la fattorizzazione QR di $A = \begin{pmatrix}4 & 1 \\ 3 & 2\end{pmatrix}$}
\addcontentsline{toc}{subsection}{Esempio}

\[
 \quad x = A(:,1) = \begin{pmatrix}4 \\ 3\end{pmatrix}
\]

\[
\|x\|_2 = \sqrt{4^2 + 3^2} = 5
\]

\[
v = x + \|x\|_2 \cdot e_1 = \begin{pmatrix}4 \\ 3\end{pmatrix} + 5 \begin{pmatrix}1 \\ 0\end{pmatrix} = \begin{pmatrix}9 \\ 3\end{pmatrix}
\]

\[
\|v\|_2 = \sqrt{9^2 + 3^2} = 3 \sqrt{10}, \quad u = \frac{v}{\|v\|_2} = \begin{pmatrix}\frac{9}{3\sqrt{10}} \\ \frac{3}{3\sqrt{10}}\end{pmatrix} = \begin{pmatrix}\frac{3}{\sqrt{10}} \\ \frac{1}{\sqrt{10}}\end{pmatrix}
\]

\[
u u^T = \begin{pmatrix} \frac{3}{\sqrt{10}} \\ \frac{1}{\sqrt{10}} \end{pmatrix} \begin{pmatrix} \frac{3}{\sqrt{10}} & \frac{1}{\sqrt{10}} \end{pmatrix} = \frac{1}{10} \begin{pmatrix}9 & 3 \\ 3 & 1 \end{pmatrix}
\]

\[
H = I - 2 u u^T = \begin{pmatrix}1 & 0 \\ 0 & 1\end{pmatrix} - 2 \cdot \frac{1}{10} \begin{pmatrix}9 & 3 \\ 3 & 1\end{pmatrix} = \begin{pmatrix} -0.8 & -0.6 \\ -0.6 & 0.8 \end{pmatrix}
\]

\[
R = H A = \begin{pmatrix} -0.8 & -0.6 \\ -0.6 & 0.8 \end{pmatrix} \begin{pmatrix}4 & 1 \\ 3 & 2 \end{pmatrix} = \begin{pmatrix} -5 & -2 \\ 0 & 2 \end{pmatrix}
\]

\[
Q = H^T = \begin{pmatrix} -0.8 & -0.6 \\ -0.6 & 0.8 \end{pmatrix}
\]

\[
A = Q R
\]

\newpage
\section*{Autovalori: Teorema di Gerschgorin e stima con norma}
\addcontentsline{toc}{section}{Autovalori: Teorema di Gerschgorin e stima con norma}

Sia \( A \in \mathbb{C}^{n \times n} \). Indichiamo con \( \mathrm{spec}(A) \) l'insieme degli autovalori di \( A \).

\subsection*{Teorema di Gerschgorin (cerchi riga)}
Ogni autovalore di una matrice quadrata \( A \in \mathbb{C}^{n \times n} \) è contenuto nell’unione di \( n \) cerchi detti \emph{cerchi riga di Gerschgorin}:
\[
\mathrm{spec}(A) \subseteq \bigcup_{i=1}^n C_i^{(R)}, \quad \text{dove} \quad 
C_i^{(R)} := \left\{ z \in \mathbb{C} : |z - a_{ii}| \leq \sum_{\substack{j=1 \\ j \neq i}}^n |a_{ij}| \right\}.
\]

\subsubsection*{Dimostrazione}
Sia \( \lambda \in \mathbb{C} \) un autovalore di \( A \), con autovettore associato \( x \in \mathbb{C}^n \setminus \{0\} \), cioè:
\[
A x = \lambda x.
\]
Sia \( k \in \{1,\dots,n\} \) tale che
\[
|x_k| = \max_{1 \leq i \leq n} |x_i|.
\]

Consideriamo la \( k \)-esima riga dell'equazione \( Ax = \lambda x \):
\[
\sum_{j=1}^n a_{kj} x_j = \lambda x_k.
\]
Riscriviamo il termine \( a_{kk} x_k \) separatamente:
\[
a_{kk} x_k + \sum_{\substack{j=1 \\ j \neq k}}^n a_{kj} x_j = \lambda x_k.
\]
Portiamo \( a_{kk} x_k \) a secondo membro:
\[
\sum_{\substack{j=1 \\ j \neq k}}^n a_{kj} x_j = (\lambda - a_{kk}) x_k.
\]
Dividiamo entrambi i membri per \( x_k \neq 0 \):
\[
\sum_{\substack{j=1 \\ j \neq k}}^n a_{kj} \frac{x_j}{x_k} = \lambda - a_{kk}.
\]
Ora passiamo al modulo:
\[
|\lambda - a_{kk}| = \left| \sum_{\substack{j=1 \\ j \neq k}}^n a_{kj} \frac{x_j}{x_k} \right| \leq \sum_{\substack{j=1 \\ j \neq k}}^n |a_{kj}| \cdot \left| \frac{x_j}{x_k} \right|.
\]
Ma per definizione di \( k \), si ha \( \left| \frac{x_j}{x_k} \right| \leq 1 \) per ogni \( j \) ($x_k = \max_{i} x_i$), quindi:nb
\[
|\lambda - a_{kk}| \leq \sum_{\substack{j=1 \\ j \neq k}}^n |a_{kj}|.
\]
Quindi \( \lambda \in C_k^{(R)} \), e poiché \( k \) era arbitrario (dipende solo da \( x \)), segue:
\[
\lambda \in \bigcup_{i=1}^n C_i^{(R)}.
\]
\qed


\subsection*{Corollario (cerchi colonna)}
Ogni autovalore di \( A \) appartiene anche all’unione dei cerchi colonna:
\[
\mathrm{spec}(A) \subseteq \bigcup_{j=1}^n C_j^{(C)}, \quad \text{dove} \quad 
C_j^{(C)} := \left\{ z \in \mathbb{C} : |z - a_{jj}| \leq \sum_{\substack{i=1\\ i \neq j}}^n |a_{ij}| \right\}.
\]

\subsubsection*{Dimostrazione (come corollario)}
Applichiamo il teorema di Gerschgorin alla matrice trasposta \( A^\top \).  
\( A \) e \( A^\top \) hanno gli stessi autovalori. I cerchi riga di \( A^\top \) sono esattamente i cerchi colonna di \( A \), poiché:
\[
(A^\top)_{ij} = A_{ji} \Rightarrow \sum_{j \neq i} |(A^\top)_{ij}| = \sum_{j \neq i} |a_{ji}| = \sum_{i \neq j} |a_{ij}|.
\]
Quindi l’inclusione degli autovalori nei cerchi colonna segue direttamente da quella nei cerchi riga di \( A^\top \).

\subsection*{Teorema di Hirsh (stima tramite norma 2)}
Sia \( \lambda \in \mathrm{spec}(A) \), allora
\[
|\lambda| \leq \|A\|_2, \quad \text{dove} \quad \|A\|_2 := \sup_{x \neq 0} \frac{\|Ax\|_2}{\|x\|_2}.
\]

\subsubsection*{Dimostrazione}
Sia \( A v = \lambda v \), con $\lambda\in spec(A)$ ovvero autovalore di A, v corrispondente autovettore normalizzato ovvero \( \|v\|_2 = 1 \). Allora:
\[ \underbrace{\|\lambda v\|_2}_{|\lambda|\cdot\|v\|_2=|\lambda|} = \|A v\|_2\Rightarrow
|\lambda| = \|A v\|_2 \leq \sup_{\|x\|_2 = 1} \|A x\|_2 = \|A\|_2.
\]
Quindi ogni autovalore di \( A \) ha modulo minore o uguale alla norma 2 di \( A \).


\section*{Metodo delle Potenze}
\addcontentsline{toc}{section}{Metodo delle Potenze}

\begin{flushleft}
\fbox{
\begin{minipage}{\linewidth}
\textsc{Input}: $A \in \mathbb{R}^{n \times n}$, $x_0 \in \mathbb{R}^n$, $it_{\text{max}}$, $\text{tol} > 0$

\item Inizializza: $\lambda = 0$, $k = 0$
\item $x_0 = \frac{x_0}{\|x_0\|_2}$
\item Per $k = 0$ a $it_{\text{max}}$:
\begin{itemize}
\item $x_{k+1} = A \cdot x_k$
\item $\lambda = x_k^T x_{k+1}$
\item $x_{k+1} = \frac{x_{k+1}}{\|x_{k+1}\|_2}$
\item Se $\|A x - \lambda x\|_2 \leq \text{tol} \cdot \|A\|_2$: termina
\end{itemize}

\textsc{Output}: $\lambda \approx \lambda_{\text{max}}$, $x \approx$ autovettore normalizzato tale che $A x \approx \lambda x$
\end{minipage}
}
\end{flushleft}

\vspace{1em}
\begin{lstlisting}[language=Matlab]
function [lambda, x] = metodo_potenze(A, x0, tol, itmax)
    lambda = 0;
    k = 0;
    xk = x0 / norm(x0);

    while k < itmax
        xkp1 = A * xk;                      % x_{k+1} = A * x_k
        lambda = xk' * xkp1;               % lambda = x_k^T * x_{k+1}
        xkp1 = xkp1 / norm(xkp1);          % normalizzazione x_{k+1}
        
        if norm(A * xkp1 - lambda * xkp1) <= tol * norm(A)
            break
        end

        xk = xkp1;                          % aggiorna x_k
        k = k + 1;
    end

    x = xkp1;
end
\end{lstlisting}
\newpage


\section*{Metodo della Potenza Inversa}
\addcontentsline{toc}{section}{Metodo della Potenza Inversa}

\begin{flushleft}
\fbox{
\begin{minipage}{\linewidth}
\textsc{Input}: $A \in \mathbb{R}^{n \times n}$ (invertibile), $x_0 \in \mathbb{R}^n$, $it_{\text{max}}$, $\text{tol} > 0$

\item Inizializza: $k = 0$
\item $x_0 = \dfrac{x_0}{\|x_0\|_2}$
\item $[L,U] = \text{FattorizzazioneLU}(A)$
\item Per $k = 0$ a $it_{\text{max}}$:
\begin{itemize}
    \item Risolvi $L U x_{k+1} = x_k$ per $x_{k+1}$
    \item $\mu = x_{k+1}^T x_k$
    \item $x_{k+1} = \dfrac{x_{k+1}}{\|x_{k+1}\|_2}$
    \item Se $\|A x_{k+1} - \mu x_{k+1}\|_2 \leq \text{tol} \cdot \|A\|_2$: termina
    \item $x_k = x_{k+1}$
\end{itemize}

\textsc{Output}: $\mu_k \approx \frac{1}{\lambda_{\min}}$, $x$
\end{minipage}
}
\end{flushleft}

\vspace{1em}
\begin{lstlisting}[language=Matlab]
function [mu, x] = potenza_inversa(A, x, tol, itmax)
    x = x / norm(x);          % Normalizza x0
    [L,U] = lu(A);            % Fattorizzazione LU
    k = 0;

    while k < itmax
        y = L \ x;            % Risolvi L*z = x
        x = U \ y;            % Risolvi U*x_{k+1} = z
        mu = x' * x;          % Stima autovalore inverso
        x = x / norm(x);      % Normalizza

        if norm(A * x - mu * x) <= tol * norm(A)
            break
        end

        k = k + 1;
    end
end
\end{lstlisting}
\newpage

\section*{Decomposizione ai Valori Singolari (SVD)}
\addcontentsline{toc}{section}{Decomposizione ai Valori Singolari (SVD)}

\begin{flushleft}
\fbox{
\begin{minipage}{\linewidth}
\textsc{Input}: $A \in \mathbb{R}^{m \times n}$

\begin{itemize}
\item Calcola $A^T A \in \mathbb{R}^{n \times n}$ (simmetrica semidefinita positiva)
\item Trova gli autovalori $\lambda_i$ e gli autovettori $v_i$ di $A^T A$
\item Ordina $\lambda_i$ in ordine decrescente: $\lambda_1 \geq \lambda_2 \geq \dots \geq \lambda_r > 0$
\item I valori singolari sono $\sigma_i = \sqrt{\lambda_i}$
\item Costruisci $V = [v_1, v_2, \dots, v_n]$
\item Calcola $u_i = \frac{1}{\sigma_i} A v_i$ per ogni $i = 1, \dots, r$
\item Completa $U$ in base ortonormale ($U \in \mathbb{R}^{m \times m}$)
\item Costruisci $\Sigma \in \mathbb{R}^{m \times n}$ tale che $\Sigma_{ii} = \sigma_i$, $i = 1, \dots, r$
\end{itemize}

\textsc{Output}: $U \in \mathbb{R}^{m \times m}$, $\Sigma \in \mathbb{R}^{m \times n}$, $V \in \mathbb{R}^{n \times n}$ tali che: $A = U \Sigma V^T$
\end{minipage}
}
\end{flushleft}

\vspace{1em}

\begin{lstlisting}[language=Matlab]
function [U, S, V] = single_value_decomposition(A)
    [m, n] = size(A); % dimensioni di A ∈ ℝ^{m×n}
    ATA = A'*A;
    % Calcola A^T A ∈ ℝ^{n×n}, matrice simmetrica e semidefinita positiva.
    
    [V,D] = eig(ATA);
    % Trova autovettori (colonne di V) e autovalori (diagonale di D) di A^T A.
    [s2,idx] = sort(diag(D),'descend');
    % Salva in s2 gli autovalori da D ordinati in modo decrescente; idx indica l’ordine.
    
    sigma = sqrt(s2);
    % Calcola valori singolari σ_i = √$\lambda$_i, con $\lambda$_i autovalori ordinati.
    V = V(:,idx);
    % Riordina le colonne di V seguendo l’ordinamento degli autovalori.
    
    r = sum(sigma > eps);
    % Calcola il rango numerico r, cioè quanti valori singolari sono significativi 
    % (maggiore di una piccola tolleranza ε).
    
    S = diag([sigma(1:r); zeros(min(m,n)-r,1)]);
    % Costruisce matrice diagonale S ∈ ℝ^{m×n} mettendo i primi r valori singolari sulla diagonale, zero il resto.
    
    U = A*V(:,1:r) ./ sigma(1:r)';
    % Calcola colonne di U come u_i = (A v_i) / σ_i,
    % cioè moltiplica A per i primi r autovettori di V e divide per il corrispondente valore singolare.
    
    if m > r, U = [U, null(U')]; end
    % Se m > r, completa U a base ortonormale aggiungendo vettori ortogonali a quelli già calcolati
    % (usando la funzione null su U^T).
end
\end{lstlisting}
\textbf{Esempio: Calcolo la fattorizzazione SVD di $A = \begin{bmatrix} 1 & 1 \\ 0 & 1 \end{bmatrix}$}
\addcontentsline{toc}{subsection}{Esempio}

\[
A = \begin{bmatrix} 1 & 1 \\ 0 & 1 \end{bmatrix} \implies A^T A =
\begin{bmatrix} 1 & 1 \\ 1 & 2 \end{bmatrix}
\]

\[
\det(A^T A - \lambda I) =
\begin{vmatrix} 1 - \lambda & 1 \\ 1 & 2 - \lambda \end{vmatrix} =
(1-\lambda)(2-\lambda) - 1 =
\lambda^2 - 3\lambda + 1
\]

\[
\lambda_{1,2} = \dfrac{3 \pm \sqrt{5}}{2} \implies
\sigma_{1,2} = \sqrt{\lambda_{1,2}}
\]

\[
\lambda_1 = \dfrac{3 + \sqrt{5}}{2} \approx 2.6180, \quad
\lambda_2 = \dfrac{3 - \sqrt{5}}{2} \approx 0.3819
\]

\[
\left(A^T A - \lambda_1 I\right)
\begin{bmatrix} x \\ y \end{bmatrix} =
\begin{bmatrix}
1 - \lambda_1 & 1 \\
1 & 2 - \lambda_1
\end{bmatrix}
\begin{bmatrix} x \\ y \end{bmatrix} =
\begin{bmatrix}
-1.6180 & 1 \\
1 & -0.6180
\end{bmatrix}
\begin{bmatrix} x \\ y \end{bmatrix} = 0
\]

\[
\Longleftrightarrow
\begin{cases}
-1.6180\,x + y = 0 \\
x - 0.6180\,y = 0
\end{cases}
\quad \implies \quad y = 1.6180\,x
\]

\[
v_1 = x \begin{bmatrix} 1 \\ 1.6180 \end{bmatrix}, \quad
\|v_1\|_2 = \sqrt{1^2 + (1.6180)^2} \approx 1.9021
\]

\[
v_1 = \frac{1}{\|v_1\|_2} \begin{bmatrix} 1 \\ 1.6180 \end{bmatrix} \approx
\begin{bmatrix} 0.5257 \\ 0.8507 \end{bmatrix}
\]

\[
\left(A^T A - \lambda_2 I\right)
\begin{bmatrix} x \\ y \end{bmatrix} =
\begin{bmatrix}
1 - \lambda_2 & 1 \\
1 & 2 - \lambda_2
\end{bmatrix}
\begin{bmatrix} x \\ y \end{bmatrix} =
\begin{bmatrix}
0.6180 & 1 \\
1 & 1.6180
\end{bmatrix}
\begin{bmatrix} x \\ y \end{bmatrix} = 0
\]

\[
\Longleftrightarrow
\begin{cases}
0.6180\,x + y = 0 \\
x + 1.6180\,y = 0
\end{cases}
\quad \implies \quad y = -0.6180\,x
\]

\[
v_2 = x \begin{bmatrix} 1 \\ -0.6180 \end{bmatrix}, \quad
\|v_2\|_2 = \sqrt{1^2 + (-0.6180)^2} \approx 1.1756
\]

\[
v_2 = \frac{1}{\|v_2\|_2} \begin{bmatrix} 1 \\ -0.6180 \end{bmatrix} \approx
\begin{bmatrix} 0.8507 \\ -0.5257 \end{bmatrix}
\]

\[
V = \begin{bmatrix} 0.5257 & 0.8507 \\ 0.8507 & -0.5257 \end{bmatrix}, \quad
V^T = \begin{bmatrix} 0.5257 & 0.8507 \\ 0.8507 & -0.5257 \end{bmatrix}^T
\]

\[
u_1 = \frac{1}{\sigma_1} A v_1 =
\frac{1}{1.6180}
\begin{bmatrix} 1 & 1 \\ 0 & 1 \end{bmatrix}
\begin{bmatrix} 0.5257 \\ 0.8507 \end{bmatrix} =
\frac{1}{1.6180}
\begin{bmatrix} 1.3764 \\ 0.8507 \end{bmatrix} \approx
\begin{bmatrix} 0.8510 \\ 0.5260 \end{bmatrix}
\]

\[
u_2 = \frac{1}{\sigma_2} A v_2 =
\frac{1}{0.6180}
\begin{bmatrix} 1 & 1 \\ 0 & 1 \end{bmatrix}
\begin{bmatrix} 0.8507 \\ -0.5257 \end{bmatrix} =
\frac{1}{0.6180}
\begin{bmatrix} 0.3250 \\ -0.5257 \end{bmatrix} \approx
\begin{bmatrix} 0.5260 \\ -0.8510 \end{bmatrix}
\]

\[
U = \begin{bmatrix} 0.8510 & 0.5260 \\ 0.5260 & -0.8510 \end{bmatrix}, \quad
\Sigma = \begin{bmatrix} 1.6180 & 0 \\ 0 & 0.6180 \end{bmatrix}, \quad
A = U \Sigma V^T
\]


\section*{Interpolazione Polinomiale}
\addcontentsline{toc}{section}{Interpolazione Polinomiale}

\begin{flushleft}
\fbox{
\begin{minipage}{\linewidth}
\textsc{Input}: intervallo $[a,b]$, grado $N$, funzione $f$

\begin{itemize}
  \item Definisci i nodi di interpolazione (Con la griglia di Gauss-Lobatto in questo caso)
\end{itemize}

\[
\mathbf{I} = \begin{bmatrix}
\frac{a + b - (b - a) \cos\left(\pi \frac{0}{N}\right)}{2} \\
\vdots \\
\frac{a + b - (b - a) \cos\left(\pi \frac{i}{N}\right)}{2} \\
\vdots \\
\frac{a + b - (b - a) \cos\left(\pi \frac{N}{N}\right)}{2}
\end{bmatrix}
\]

\begin{itemize}
  \item Calcola i valori $f(\mathbf{I})$
  \item Inizializza $p(x) = 0$
\end{itemize}

\begin{itemize}
  \item Per ogni $i = 0, \dots, N$:
    \begin{itemize}
      \item Calcola i polinomi base di Lagrange:
      \[
      l_i(x) = \prod_{\substack{j=0 \\ j \neq i}}^{N} \frac{x - I_j}{I_i - I_j}
      \]
      \item Costruisci il polinomio interpolatore:
      \[
      p(x) = p(x) + f(I_i) l_i(x)
      \]
    \end{itemize}
  \item Calcola il polinomio nodale:
  \[
  \omega_n(x) = \prod_{i=0}^N (x - I_i)
  \]
  \item Calcola la norma dell’errore usando la derivata $(N+1)$-esima di $f$ e $\omega_n(x)$
\end{itemize}

\textsc{Output}: norma dell’errore $\| \cdot \|_{\infty}$
\end{minipage}
}
\end{flushleft}

\vspace{1em}

\newpage
\begin{lstlisting}[language=Matlab]
function [Norm] = Interpolatore(a,b,N,f)
% Interpolazione polinomiale di grado N su [a,b] con funzione f

I = (a+b-(b-a)*cos(pi*(0:N)/N))/2; % nodi Gauss-Lobatto
fI = f(I);

corretto = linspace(a,b,15*N);
li = ones(N+1, length(corretto));
den = ones(N+1,1);

for i = 1:N+1
    div = I - I(i);
    div(i) = 1; % evitare divisione per zero
    den(i) = prod(div);
    
    % calcola il polinomio di Lagrange base l_i(x) su tutti i punti di corretto
    m = corretto - I([1:i-1,i+1:end])';
    li(i,:) = prod(m,1) ./ den(i);
end

out = fI * li; % polinomio interpolatore

hold off
figure(100)
plot(corretto, out, 'LineWidth', 2)
hold on
plot(corretto, f(corretto), 'LineWidth', 2)

nod = prod(corretto - I', 1); % polinomio nodale

figure(200)
plot(corretto, nod)

% calcolo errore
g = f;
for i = 1:N+1
    g = matlabFunction(diff(sym(g)));
end

try
    csi = max(g(corretto));
    Norm = norm(nod .* csi / factorial(N+1), inf);
catch
    Norm = 0;
end
end
\end{lstlisting}
\newpage
\subsection*{Stima dell'errore nell'interpolazione polinomiale}
\addcontentsline{toc}{subsection}{Stima dell'errore}
Siano dati $n+1$ nodi $x_0 < \dots < x_n \in \mathbb{R}$,  
sia $f:\mathbb{R}\to \mathbb{R}$ t.c. $f \in C^{n+1}([x_0,x_n])$.

Sia $P(x)$ il polinomio interpolatore di grado $\leq n$ tale che:
\[
P(x_i) = f(x_i) \quad \forall i \in \{0,\dots,n\}
\]
\[
E(x) := f(x) - P(x)
\Rightarrow E(x_i) = 0 \quad
\forall i \in \{0,\dots,n\}
\]
\[
\omega(x) := \prod_{i=0}^n (x - x_i)
\Rightarrow \exists g(x) \in C^1 \text{ tale che } E(x) = \omega(x) \cdot g(x)
\]
\[\text{Derivando n+1 volte }
E^{(n+1)}(x) = (f - P)^{(n+1)}(x) = f^{(n+1)}(x) - 0 = f^{(n+1)}(x)
\]
\[
(\omega \cdot g)^{(n+1)}(x) =
\sum_{k=0}^{n+1} \binom{n+1}{k} \cdot \omega^{(k)}(x) \cdot g^{(n+1-k)}(x)
\]

$\deg(\omega) = n+1 \Rightarrow \omega^{(n+1)}(x) = (n+1)!  
\Rightarrow  \omega^{(k)}(x) = 0 \textbf{ } \forall k > n+1$.
\[
E^{(n+1)}(x) = f^{(n+1)}(x) = (n+1)! \cdot g(x) + \sum_{k=0}^{n} \binom{n+1}{k} \cdot \omega^{(k)}(x) \cdot g^{(n+1-k)}(x)
\]
Per il teorema di Lagrange $\exists \, \xi \in [x_0,x_n]$ tale che:
\[
f^{(n+1)}(\xi) = (n+1)! \cdot g(\xi)
\Rightarrow
g(\xi) = \frac{f^{(n+1)}(\xi)}{(n+1)!}
\]
\[
f(x) - P(x) = \frac{f^{(n+1)}(\xi)}{(n+1)!} \cdot \prod_{i=0}^{n} (x - x_i),
\quad \xi \in [x_0,x_n]
\]

\newpage

\section*{Metodo del Punto Medio}
\addcontentsline{toc}{section}{Metodo del Punto Medio}

\begin{flushleft}
\fbox{
\begin{minipage}{\linewidth}
\textsc{Input}: intervallo $[a,b]$, funzione $f$, numero di sottointervalli $m$

\begin{itemize}
  \item Dividi $[a,b]$ in $m$ sottointervalli uguali con nodi:
  \[
  \mathbf{grid} = \left[ x_0, x_1, \dots, x_m \right], \quad x_i = a + i \cdot \frac{b-a}{m}
  \]
  \item Inizializza l’approssimazione integrale $I = 0$
  \item Calcola la lunghezza di ogni sottointervallo:
  \[
  len = \frac{b-a}{m}
  \]
  \item Per ogni sottointervallo $i = 1, \dots, m$:
  \begin{itemize}
    \item Calcola il punto medio:
    \[
    c_i = \frac{x_{i-1} + x_i}{2}
    \]
    \item Aggiorna l’integrale approssimato:
    \[
    I = I + f(c_i) \cdot len
    \]
  \end{itemize}
  \item Calcola l’errore come valore assoluto della differenza tra l’integrale esatto e $I$:
  \[
  error = \left| \int_a^b f(x) dx - I \right|
  \]
\end{itemize}

\textsc{Output}: approssimazione $I$ dell’integrale e errore $error$
\end{minipage}
}
\end{flushleft}

\vspace{1em}

\begin{lstlisting}[language=Matlab]
function [I,error] = Midpoint(a,b,f,m)
% Metodo del punto medio per l'integrazione numerica

grid = linspace(a,b,m+1);    % nodi della suddivisione
I = 0;
len = (b - a) / m;           % lunghezza sottointervalli

for i = 1:m
    c = (grid(i) + grid(i+1)) / 2; % punto medio del sottointervallo
    I = I + f(c) * len;             % somma rettangoli
end

error = abs(integral(f,a,b) - I); % errore assoluto rispetto a integrale esatto
end
\end{lstlisting}

\newpage
\subsection*{Stima dell'errore nel metodo del punto medio}
\addcontentsline{toc}{subsection}{Stima dell'errore}


\[
I = \int_a^b f(x)\,dx \approx (b - a) f(c), \quad c = \frac{a + b}{2}
\]
\[
E = \int_a^b f(x)\,dx - (b - a)f(c)
\]
Sviluppo di Taylor:$  \quad
f(x) = f(c) + f'(c)(x - c) + \frac{f''(\xi)}{2}(x - c)^2$ con $\xi \in [a,b]$
\[\Rightarrow
\int_a^b f(x)\,dx = f(c)(b - a) + f'(c)\underbrace{\int_a^b (x - c)\,dx}_{\frac{(b^2-a^2)}{2}-\frac{(b+a)(b-a)}{2}=0} + f''(\xi)\frac{(b - a)^3}{24}
\]

\[ \Rightarrow
E = -\frac{(b - a)^3}{24} f''(\xi)
\]
\textbf{Caso composto su \(m\) sottointervalli lunghi \(h = \frac{b - a}{m}\)}

Siano \(x_0 = a,\, x_1 = a + h,\, \dots,\, x_m = b\) i nodi. Per ogni \(i\), punto medio \(c_i = \frac{x_i + x_{i+1}}{2}\).

Approssimazione: \(I \approx \sum_{i = 0}^{m-1} h \cdot f(c_i)\)

Errore su ogni intervallo: \(E_i = -\frac{h^3}{24} f''(\xi_i)\), con \(\xi_i \in (x_i, x_{i+1})\)

Errore totale: \(E = \sum_{i=0}^{m-1} E_i = -\frac{h^3}{24} \sum_{i=0}^{m-1} f''(\xi_i)\)

\(M = \max_{x \in [a,b]} |f''(x)|\), \(m = \frac{b-a}{h}\) \Rightarrow
\[|E| \leq \frac{h^3}{24} \cdot m \cdot M = \frac{h^3}{24} \cdot \frac{b-a}{h} \cdot M = \frac{(b-a)h^2}{24} \cdot M\]

\section*{Metodo dei Trapezi}
\addcontentsline{toc}{section}{Metodo dei Trapezi}

\begin{flushleft}
\fbox{
\begin{minipage}{\linewidth}
\textsc{Input}: intervallo $[a,b]$, funzione $f$, numero di sottointervalli $m$

\begin{itemize}
  \item Suddividi l’intervallo $[a,b]$ in $m$ sottointervalli uguali:
  \[
  \mathbf{grid} = \left[ x_0, x_1, \dots, x_m \right], \quad x_i = a + i \frac{b-a}{m}
  \]
  \item Inizializza $I = 0$
  \item Calcola la lunghezza di ogni sottointervallo:
  \[
  len = \frac{b-a}{m}
  \]
  \item Per ogni sottointervallo $i = 1, \dots, m$:
  \begin{itemize}
    \item Calcola la media dei valori della funzione agli estremi:
    \[
    c_i = \frac{f(x_{i-1}) + f(x_i)}{2}
    \]
    \item Aggiorna la somma dell’integrale:
    \[
    I = I + c_i \cdot len
    \]
  \end{itemize}
  \item Calcola l’errore come:
  \[
  error = \left| \int_a^b f(x) dx - I \right|
  \]
\end{itemize}

\textsc{Output}: approssimazione $I$ dell’integrale, errore $error$
\end{minipage}
}
\end{flushleft}

\vspace{1em}

\begin{lstlisting}[language=Matlab]
function [I, error] = Trapezi(a,b,f,m)
% Metodo dei trapezi per l'integrazione numerica

grid = linspace(a,b,m+1);    % nodi della suddivisione
I = 0;
len = (b - a) / m;           % lunghezza sottointervalli

for i = 1:m
    c = (f(grid(i)) + f(grid(i+1))) / 2; % media dei valori della funzione
    I = I + c * len;                     % somma area trapezi
end

error = abs(integral(f,a,b) - I);       % errore assoluto
end
\end{lstlisting}

\newpage
\subsection*{Stima dell'errore nel metodo dei trapezi}
\addcontentsline{toc}{subsection}{Stima dell'errore}

\[
I = \int_a^b f(x)\,dx \approx \frac{b - a}{2} [f(a) + f(b)],\quad c= \frac{b+a}{2}
\]
\[
E = \int_a^b f(x)\,dx - \frac{b - a}{2} [f(a) + f(b)]
\]
\text{Sviluppiamo con Taylor su \(c = \frac{a + b}{2}\) per semplificare i conti:}

\[\begin{cases}
f(a) = f(c) + f'(c)(a-c) + \frac{f''(\xi)}{2}(a-c)^2 = f(c) - f'(c)\frac{b - a}{2} + \frac{f''(\xi)}{2} \left(\frac{b - a}{2}\right)^2\\
f(b) = f(c)+f'(c)(b-c)+\frac{f''(\xi)}{2}(b-c)^2=f(c) + f'(c)\frac{b - a}{2} + \frac{f''(\xi)}{2} \left(\frac{b - a}{2}\right)^2\end{cases}
\]

\[
\Rightarrow f(a) + f(b) = 2f(c) + f''(\xi) \cdot \frac{(b - a)^2}{4}
\]

\[
\frac{b - a}{2}[f(a) + f(b)] \approx (b - a)f(c) + \frac{f''(\xi)}{8}(b - a)^3
\]

\[
\int_a^b f(x)\,dx = f(c)(b - a) + f'(c)\underbrace{\int_a^b (x - c)\,dx}_{=0} + \frac{f''(\xi)}{2}\underbrace{\int_a^b (x - c)^2\,dx}_{\frac{(b-a)^3}{12}}
\]

\[\int_a^b f(x)\,dx = f(c)(b - a) + \frac{f''(\xi)}{2} \cdot \frac{(b-a)^3}{12} = f(c)(b - a) + \frac{f''(\xi)(b-a)^3}{24}
\]

Quindi
$E = \frac{f''(\xi)(b-a)^3}{24} - \frac{f''(\xi)(b-a)^3}{8} = f''(\xi)(b-a)^3\left(\frac{1}{24} - \frac{3}{24}\right) = -\frac{f''(\xi)(b-a)^3}{12}$

\vspace{1em}

\textbf{Caso composto su \(m\) sottointervalli lunghi \(h = \frac{b - a}{m}\)}

Siano \(x_0 = a,\, x_1 = a + h,\, \dots,\, x_m = b\) i nodi.
\[I \approx \frac{h}{2} \left[f(x_0) + 2\sum_{i=1}^{m-1} f(x_i) + f(x_m)\right]\]

Errore sull'i-esimo intervallo: \(E_i = -\frac{h^3}{12} f''(\xi_i)\), con \(\xi_i \in [x_i, x_{i+1}]\)

Errore totale:
\[
E = \sum_{i=0}^{m-1} E_i = -\frac{h^3}{12} \sum_{i=0}^{m-1} f''(\xi_i)
\]

$M = \max_{x \in [a,b]} |f''(x)| \Rightarrow
|E| \leq \frac{h^3}{12} \cdot m \cdot M = \frac{h^3}{12} \cdot \frac{b-a}{h} \cdot M = \frac{(b-a)h^2}{12} \cdot M$


\section*{Metodo di Cavalieri-Simpson}
\addcontentsline{toc}{section}{Metodo di Cavalieri-Simpson}

\begin{flushleft}
\fbox{
\begin{minipage}{\linewidth}
\textsc{Input}: intervallo $[a,b]$, funzione $f$, numero di sottointervalli $m$

\begin{itemize}
  \item Suddividi l’intervallo $[a,b]$ in $m$ sottointervalli:
  \[
  \mathbf{grid} = \left[ x_0, x_1, \dots, x_m \right], \quad x_i = a + i \frac{b-a}{m}
  \]
  \item Inizializza $I = 0$
  \item Calcola la lunghezza di ogni sottointervallo:
  \[
  len = \frac{b-a}{m}
  \]
  \item Per ogni sottointervallo $i = 1, \dots, m$:
  \begin{itemize}
    \item Calcola il punto medio:
    \[
    c_i = \frac{x_{i-1} + x_i}{2}
    \]
    \item Approssima l'integrale locale con la formula di Simpson:
    \[
    I = I + \frac{f(x_{i-1}) + 4f(c_i) + f(x_i)}{6} \cdot len
    \]
  \end{itemize}
  \item Calcola l’errore:
  \[
  error = \left| \int_a^b f(x)dx - I \right|
  \]
\end{itemize}

\textsc{Output}: approssimazione $I$ dell’integrale, errore $error$
\end{minipage}
}
\end{flushleft}

\vspace{1em}

\begin{lstlisting}[language=Matlab]
function [I, error] = Cavalieri_Simpson(a,b,f,m)
% Metodo di Cavalieri-Simpson per l'integrazione numerica

grid = linspace(a,b,m+1);    % nodi della suddivisione
I = 0;
len = (b - a) / m;           % lunghezza sottointervalli

for i = 1:m
    c = (grid(i) + grid(i+1)) / 2; % punto medio del sottointervallo
    I = I + (f(grid(i)) + 4*f(c) + f(grid(i+1))) * len / 6;
end

error = abs(integral(f,a,b) - I); % errore assoluto
end
\end{lstlisting}

\newpage
\subsection*{Stima dell'errore nel metodo di Cavalieri-Simpson}
\addcontentsline{toc}{subsection}{Stima dell'errore}

\[
I = \int_a^b f(x)\,dx \approx \frac{b - a}{6} [f(a) + 4f(c) + f(b)],\quad c= \frac{b+a}{2}
\]
\[
E = \int_a^b f(x)\,dx - \frac{b - a}{6} [f(a) + 4f(c) + f(b)]
\]
\text{Sviluppiamo con Taylor su \(c = \frac{a + b}{2}\), con \(h = \frac{b - a}{2}\):}

\[\begin{cases}
f(a) = f(c) - f'(c)h + \frac{f''(c)}{2}h^2 - \frac{f^{(3)}(c)}{6}h^3 + \frac{f^{(4)}(\xi_1)}{24}h^4\\
f(b) = f(c) + f'(c)h + \frac{f''(c)}{2}h^2 + \frac{f^{(3)}(c)}{6}h^3 + \frac{f^{(4)}(\xi_2)}{24}h^4
\end{cases}
\]

\[
\Rightarrow f(a) + f(b) = 2f(c) + f''(c)h^2 + \frac{f^{(4)}(\xi)}{12}h^4,\quad \text{per un certo } \xi \in (a,b)
\]

\[
f(a) + 4f(c) + f(b) = 6f(c) + f''(c)h^2 + \frac{f^{(4)}(\xi)}{12}h^4
\]

\[
\frac{b - a}{6} [f(a) + 4f(c) + f(b)] = 2h f(c) + \frac{f''(c)}{6}h^3 + \frac{f^{(4)}(\xi)}{72}h^5
\]

\[
\int_a^b f(x)\,dx = 2h f(c) + \frac{f''(c)}{6}h^3 + \frac{f^{(4)}(\xi)}{90}h^5
\]

\[
E = \frac{f^{(4)}(\xi)}{90}h^5 - \frac{f^{(4)}(\xi)}{72}h^5 = f^{(4)}(\xi)h^5\left(\frac{1}{90} - \frac{1}{72}\right) = -\frac{f^{(4)}(\xi)}{2160}(b - a)^5
\]

\textbf{Caso composto su \(m\) sottointervalli lunghi \(h = \frac{b - a}{m}\), con \(m\) pari}

Siano \(x_0 = a,\, x_1 = a + h,\, \dots,\, x_m = b\) i nodi.

Approssimazione:
\[
I \approx \sum_{i=1}^m \frac{h}{6} [f(x_{i-1}) + 4f(c_i) + f(x_i)],\quad c_i = \frac{x_{i-1} + x_i}{2}
\]

Errore su ogni intervallo:
\[
E_i = -\frac{h^5}{2880} f^{(4)}(\xi_i),\quad \xi_i \in [x_{i-1}, x_i]
\]

Errore totale:
\[
E = \sum_{i=1}^m E_i = -\frac{h^5}{2880} \sum_{i=1}^m f^{(4)}(\xi_i)
\]

Se \(M = \max_{x \in [a,b]} |f^{(4)}(x)|\), allora:
\[
|E| \leq \frac{h^5}{2880} \cdot m \cdot M = \frac{h^5}{2880} \cdot \frac{b-a}{h} \cdot M = \frac{(b-a)h^4}{2880} \cdot M
\]



\section*{Confronto degli Errori dei Metodi di Integrazione}
\addcontentsline{toc}{section}{Confronto degli Errori dei Metodi di Integrazione}

\begin{flushleft}
\fbox{
\begin{minipage}{\linewidth}
\textsc{Input}: funzione $f$, intervallo $[a,b] = [0,1]$, massimo numero di sottointervalli $n = 100$

\begin{itemize}
  \item Per ogni $i = 1, \dots, n$:
  \begin{itemize}
    \item Calcola errore del metodo del punto medio con $i$ sottointervalli
    \item Calcola errore del metodo dei trapezi con $i$ sottointervalli
    \item Calcola errore del metodo di Cavalieri-Simpson con $i$ sottointervalli
  \end{itemize}
  \item Salva ciascun errore in un vettore
  \item Plotta in scala semilogaritmica $\log_{10}(\text{errore})$ rispetto al numero di sottointervalli
\end{itemize}

\textsc{Output}: grafico con curve di errore per i tre metodi
\end{minipage}
}
\end{flushleft}

\vspace{1em}

\begin{lstlisting}[language=Matlab]
clear
close all

% Plotter degli errori
n = 100;
a = 0;
b = 1;

functions = {@(x) sin(x).*cos(x), @(x) exp(x), @(x) exp(sin(x + 2*x.^2)), @(x) atan(log(abs(x)+1))};
f = functions{1};

Mid = zeros(1,n);
Cav = zeros(1,n);
Trap = zeros(1,n);

for i = 1:n
    [~, e] = Midpoint(a,b,f,i);
    Mid(i) = e;

    [~, e] = Trapezi(a,b,f,i);
    Trap(i) = e;

    [~, e] = Cavalieri_Simpson(a,b,f,i);
    Cav(i) = e;
end

figure(1)
semilogy(1:n, Mid, 'r', 1:n, Trap, 'g', 1:n, Cav, 'b')
legend('Midpoint', 'Trapezi', 'Cavalieri Simpson')
\end{lstlisting}

\section*{Interpolazione Polinomiale Composita}
\addcontentsline{toc}{section}{Interpolazione Polinomiale Composita}

\begin{flushleft}
\fbox{
\begin{minipage}{\linewidth}
\textsc{Input}: intervallo $[a,b]$, funzione $f$, numero di sottointervalli $m$, grado $N$

\begin{itemize}
  \item Suddividi $[a,b]$ in $m$ sottointervalli: $x_i = a + i \frac{b-a}{m}$
  \item Per ogni sottointervallo $[x_i, x_{i+1}]$:
    \begin{itemize}
      \item Applica interpolazione polinomiale di grado $N$ con griglia uniforme
      \item Valuta il polinomio in 9 punti equispaziati
      \item Inserisci i risultati nel vettore finale
    \end{itemize}
  \item Calcola l’errore come $\|f(x) - \text{interpolato}(x)\|_{\infty}$
\end{itemize}

\textsc{Output}: funzione interpolata e errore massimo
\end{minipage}
}
\end{flushleft}

\vspace{1em}

\begin{lstlisting}[language=Matlab]
function [out]= Interpolatore(a,b,N,f) % il polinomio approssimatore è di grado N
% ovvero il numero di punti che scegliamo -1
I=linspace(a,b,N+1); % Griglia normale, non Gauss-Lobatto
fI=f(I);
corretto=linspace(a,b,10);   % MODIFICA: qui non riscaliamo tutto per N perchè non ci serve
li=1;
den=1;
for i= 1:N+1
    div=I-I(i);
    div(i)=1;
    den=den.*div;
    m=corretto-(ones(N+1,1)*I(i)); 
    m(i,:)=1;                             
    li=li.*m;
end
%ora plottiamo i risultati
out=fI*(li./den');

function [errore] = Pol_int_comp(a,b,f,m,deg)
close
grid = linspace(a,b,m+1);
simul= zeros(1,9*m);

for i = 1:m
    out=Interpolatore(grid(i),grid(i+1),deg,f);
    out=out(1, 1:9);
    simul(1+(9*(i-1)): 9*i)=out;
end

grid= linspace(a,b, 9*m);
figure(100)
plot(grid, simul)
hold on

errore = norm(abs(simul-f(grid)),'inf');
\end{lstlisting}


\section*{Equazioni Differenziali Ordinarie (ODE)}
\addcontentsline{toc}{section}{Equazioni Differenziali Ordinarie (ODE)}

\begin{flushleft}
Una equazione differenziale ordinaria (ODE) ha la forma:
\begin{equation*}
y'(t) = f(t, y(t)), \quad y(t_0) = y_0
\end{equation*}

\textbf{Significato:}
\begin{itemize}
  \item \(t\): variabile indipendente (es. tempo)
  \item \(y(t)\): variabile dipendente (es. posizione, temperatura...)
  \item \(f(t, y(t))\): regola che descrive come varia \(y(t)\)
\end{itemize}

L'obiettivo \`e trovare la funzione \(y(t)\) che soddisfa l'equazione e il valore iniziale.

\textbf{Problema di Cauchy:} trovare \(y(t)\) tale che:
\begin{equation*}
y'(t) = f(t, y(t)), \quad y(t_0) = y_0
\end{equation*}

\subsection*{Problema di Cauchy: Esempio Numerico}

Consideriamo il problema:

\[
\frac{dy}{dt} = -2y, \qquad y(0) = 1
\]

\vspace{0.5em}
\textbf{Cosa vogliamo fare:} vogliamo trovare una funzione \( y(t) \) che soddisfa questa equazione differenziale, cioè:

\begin{itemize}
    \item \( y'(t) = -2y(t) \), ovvero sappiamo come varia nel tempo
    \item \( y(0) = 1 \), ovvero conosciamo il valore iniziale
\end{itemize}

\vspace{0.5em}
\textbf{Input che diamo a un metodo numerico:}

Non conosciamo \( y(t) \), ma conosciamo la funzione derivata:

\[
f(t,y) = -2y
\]

Questa funzione prende in input:
\begin{itemize}
    \item un tempo \( t \)
    \item un valore approssimato \( y \)
\end{itemize}
e restituisce la derivata \( \frac{dy}{dt} \) in quel punto.

\vspace{0.5em}
\textbf{Output che vogliamo:}

Vogliamo ottenere una lista di punti:

\[
(t_0, y_0), (t_1, y_1), \dots, (t_n, y_n)
\]

ovvero:
\begin{itemize}
    \item una lista di tempi \( t_i \)
    \item e i valori approssimati \( y(t_i) \) della soluzione
\end{itemize}

\section*{Metodo di Eulero Esplicito (Forward Euler)}
\addcontentsline{toc}{section}{Metodo di Eulero Esplicito (Forward Euler)}

\fbox{
\begin{minipage}{\linewidth}
\textsc{Input}: \(t_0\), \(y_0\), passo \(h\), numero di passi \(N\), funzione \(f\)

\begin{itemize}
  \item \textbf{Per ogni} \(n = 0, \dots, N-1\):
  \begin{itemize}
    \item \(t_{n+1} = t_n + h\)
    \item \(y_{n+1} = y_n + h \cdot f(t_n, y_n)\)
  \end{itemize}
\end{itemize}

\textsc{Output}: valori approssimati \(y_0, y_1, \dots, y_N\), errore assoluto (se soluzione esatta nota), grafico soluzione approssimata vs esatta, ordine di convergenza
\end{minipage}}

\vspace{1em}
\textbf{Codice MATLAB}
\begin{lstlisting}[language=Matlab]
function [t, y, err] = Eulero(t0, y0, h, N, f, y_exact)
    t = t0:h:(t0 + N*h);
    y = zeros(1,N+1);
    y(1) = y0;
    for n = 1:N
        y(n+1) = y(n) + h * f(t(n), y(n));
    end
    if nargin == 6
        err = abs(y - y_exact(t));
        figure; plot(t, y, 'b', t, y_exact(t), 'r--');
        legend('Numerica','Esatta'); title('Eulero Esplicito');
    else
        err = NaN;
    end
end
\end{lstlisting}

\vspace{1em}

\newpage
\section*{Metodo di Eulero Implicito (Backward Euler)}
\addcontentsline{toc}{section}{Metodo di Eulero Implicito (Backward Euler)}

\fbox{
\begin{minipage}{\linewidth}
\textsc{Input}: \(t_0\), \(y_0\), passo \(h\), numero di passi \(N\), funzione \(f\)

\begin{itemize}
  \item \textbf{Per ogni} \(n = 0, \dots, N-1\):
  \begin{itemize}
    \item \(t_{n+1} = t_n + h\)
    \item risolvere: \(y_{n+1} = y_n + h \cdot f(t_{n+1}, y_{n+1})\)
  \end{itemize}
\end{itemize}

\textsc{Output}: valori approssimati \(y_0, y_1, \dots, y_N\), errore assoluto (se soluzione esatta nota), grafico soluzione approssimata vs esatta, ordine di convergenza
\end{minipage}}

\vspace{1em}
\textbf{Codice MATLAB}
\begin{lstlisting}[language=Matlab]
function [t, y, err] = EuleroImp(t0, y0, h, N, f, y_exact)
    t = t0:h:(t0 + N*h);
    y = zeros(1,N+1);
    y(1) = y0;
    for n = 1:N
        g = @(yn1) yn1 - y(n) - h * f(t(n+1), yn1);
        y(n+1) = fsolve(g, y(n));
    end
    if nargin == 6
        err = abs(y - y_exact(t));
        figure; plot(t, y, 'b', t, y_exact(t), 'r--');
        legend('Numerica','Esatta'); title('Eulero Implicito');
    else
        err = NaN;
    end
end
\end{lstlisting}

\vspace{1em}
\newpage
\section*{Metodo di Crank-Nicolson}
\addcontentsline{toc}{section}{Metodo di Crank-Nicolson}

\begin{flushleft}
\fbox{
\begin{minipage}{\linewidth}
\textsc{Input}: passo temporale \(\Delta t\), numero di iterazioni \(N\), funzione \(f(t,y)\), valore iniziale \(y_0\)

\begin{itemize}
\item Per \(j = 0, \dots, N-1\):
\[
y_{j+1} = y_j + \frac{\Delta t}{2} \big( f(t_j, y_j) + f(t_{j+1}, y_{j+1}) \big)
\]
\end{itemize}

\textsc{Output}: valori approssimati \(y_j\) per \(j=1,\dots,N\)
\end{minipage}
}
\end{flushleft}

\vspace{1em}

\begin{lstlisting}[language=Matlab]
function y = crankNicolson(f, t0, y0, dt, N)
    y = zeros(1, N+1);
    t = t0 + (0:N)*dt;
    y(1) = y0;
    for j = 1:N
        % Risolvo y_{j+1} implicitamente con fzero o altra tecnica
        fun = @(Y) Y - y(j) - dt/2*(f(t(j), y(j)) + f(t(j+1), Y));
        y(j+1) = fzero(fun, y(j)); % o altro metodo di risoluzione
    end
end
\end{lstlisting}

\vspace{1em}


\newpage

\section*{Metodo di Heun}
\addcontentsline{toc}{section}{Metodo di Heun}

\begin{flushleft}
\fbox{
\begin{minipage}{\linewidth}
\textsc{Input}: passo \(h\), numero di passi \(N\), funzione \(f(t,y)\), condizioni iniziali \(t_0, y_0\)

\begin{itemize}
\item Per \(j = 0, \dots, N-1\):
\[
y_{j+1} = y_j + \frac{h}{2} \left( f(t_j, y_j) + f\big(t_{j+1}, y_j + h f(t_j, y_j) \big) \right)
\]
\end{itemize}

\textsc{Output}: valori approssimati \(y_j\) per \(j=1,\dots,N\)
\end{minipage}
}
\end{flushleft}

\vspace{1em}

\textbf{Codice MATLAB}

\begin{lstlisting}[language=Matlab]
function [t, y, err] = Heun(t0, y0, h, N, f, y_exact)
    t = t0:h:(t0 + N*h);
    y = zeros(1, N+1);
    y(1) = y0;
    for j = 1:N
        k1 = f(t(j), y(j));
        k2 = f(t(j+1), y(j) + h*k1);
        y(j+1) = y(j) + h/2 * (k1 + k2);
    end
    if nargin == 6
        err = abs(y - y_exact(t));
        figure; plot(t, y, 'b', t, y_exact(t), 'r--');
        legend('Numerica','Esatta'); title('Heun');
    else
        err = NaN;
    end
end
\end{lstlisting}

\newpage

\section*{Metodo di Runge-Kutta del 4\textdegree ordine (RK4)}
\addcontentsline{toc}{section}{Metodo di Runge-Kutta del 4\textdegree ordine}

\fbox{
\begin{minipage}{\linewidth}
\textsc{Input}: \(t_0\), \(y_0\), passo \(h\), numero di passi \(N\), funzione \(f\)

\begin{itemize}
  \item \textbf{Per ogni} \(n = 0, \dots, N-1\):
  \begin{itemize}
    \item \(t_{n+1} = t_n + h\)
    \item \(k_1 = f(t_n, y_n)\)
    \item \(k_2 = f(t_n + \frac{h}{2}, y_n + \frac{h}{2} k_1)\)
    \item \(k_3 = f(t_n + \frac{h}{2}, y_n + \frac{h}{2} k_2)\)
    \item \(k_4 = f(t_n + h, y_n + h \cdot k_3)\)
    \item \(y_{n+1} = y_n + \frac{h}{6} (k_1 + 2k_2 + 2k_3 + k_4)\)
  \end{itemize}
\end{itemize}

\textsc{Output}: valori approssimati, errore, grafico, ordine
\end{minipage}}

\vspace{1em}
\textbf{Codice MATLAB}
\begin{lstlisting}[language=Matlab]
function [t, y, err] = RK4(t0, y0, h, N, f, y_exact)
    t = t0:h:(t0 + N*h);
    y = zeros(1,N+1);
    y(1) = y0;
    for n = 1:N
        k1 = f(t(n), y(n));
        k2 = f(t(n) + h/2, y(n) + h/2 * k1);
        k3 = f(t(n) + h/2, y(n) + h/2 * k2);
        k4 = f(t(n) + h, y(n) + h * k3);
        y(n+1) = y(n) + h/6 * (k1 + 2*k2 + 2*k3 + k4);
    end
    if nargin == 6
        err = abs(y - y_exact(t));
        figure; plot(t, y, 'b', t, y_exact(t), 'r--');
        legend('Numerica','Esatta'); title('RK4');
    else
        err = NaN;
    end
end
\end{lstlisting}


\end{flushleft}


\newpage
\section*{Test completo degli algoritmi per ODE}
\addcontentsline{toc}{section}{Test degli algoritmi per ODE}

Sia dato il problema generale:  
\[
\dot{y} = f(t,y), \quad y(t_0) = y_0
\]
dove \(y \in \mathbb{R}^m\).  
Il dominio temporale è \([t_0, T]\) con passo \(h\) e \(N = \frac{T - t_0}{h}\).

Definiamo il vettore tempo  
\[
t_n = t_0 + n h, \quad n = 0, \ldots, N.
\]

\vspace{1em}

\noindent
\textbf{Parametri comuni in MATLAB:}
\begin{lstlisting}
T = 5;        % tempo finale
h = 0.1;      % passo
t0 = 0;       % tempo iniziale
N = floor((T - t0)/h);
t_grid = t0:h:(t0 + N*h);
\end{lstlisting}

\vspace{1em}

\noindent
\textbf{Definizione metodi da testare:}
\begin{lstlisting}
methods = {@Eulero, @EuleroImp, @Heun, @RK4};
method_names = {'Eulero Esplicito', 'Eulero Implicito', 'Heun (RK2)', 'Runge-Kutta 4'};
\end{lstlisting}

\section*{(a) Problema scalare:}  
\[
\dot{y} = - y \log y, \quad y(0) = 0.5
\]
Soluzione esatta:  
\[
y(t) = e^{-e^{\log(\log(2)) - t}}
\]

\noindent
\textbf{Codice MATLAB:}
\begin{lstlisting}[language=Matlab]
f_a = @(t,y) -y.*log(y);
y0_a = 0.5;
y_exact_a = @(t) exp(-exp(log(log(2)) - t));

fprintf('Test (a)\n');
for i = 1:length(methods)
    [t, y, err] = methods{i}(t0, y0_a, h, N, f_a, y_exact_a);
    figure(1); hold on;
    plot(t, y, 'DisplayName', method_names{i});
    figure(2); hold on;
    plot(t, err, 'DisplayName', method_names{i});
end
figure(1);
plot(t_grid, y_exact_a(t_grid), 'k--', 'DisplayName', 'Esatta');
title('Soluzioni Test (a)');
xlabel('t'); ylabel('y');
legend('show'); hold off;

figure(2);
title('Errore assoluto Test (a)');
xlabel('t'); ylabel('|Errore|');
legend('show'); hold off;
\end{lstlisting}

\vspace{2em}

\section*{(b) Problema scalare:}  
\[
\dot{y} = - e^{-(t + y)}, \quad y(0) = 1
\]
Soluzione esatta:  
\[
y(t) = \log\bigl(e + e^{-t} - 1\bigr)
\]

\noindent
\textbf{Codice MATLAB:}
\begin{lstlisting}[language=Matlab]
f_b = @(t,y) -exp(-(t + y));
y0_b = 1;
y_exact_b = @(t) log(exp(1) + exp(-t) - 1);

fprintf('Test (b)\n');
for i = 1:length(methods)
    [t, y, err] = methods{i}(t0, y0_b, h, N, f_b, y_exact_b);
    figure(3); hold on;
    plot(t, y, 'DisplayName', method_names{i});
    figure(4); hold on;
    plot(t, err, 'DisplayName', method_names{i});
end
figure(3);
plot(t_grid, y_exact_b(t_grid), 'k--', 'DisplayName', 'Esatta');
title('Soluzioni Test (b)');
xlabel('t'); ylabel('y');
legend('show'); hold off;

figure(4);
title('Errore assoluto Test (b)');
xlabel('t'); ylabel('|Errore|');
legend('show'); hold off;
\end{lstlisting}
\newpage
\section*{(c) Problema scalare:}  
\[
\dot{y} = y (1 - y), \quad y(0) = 0.5
\]
Soluzione esatta:  
\[
y(t) = \frac{e^t}{1 + e^t}
\]

\noindent
\textbf{Codice MATLAB:}
\begin{lstlisting}[language=Matlab]
f_c = @(t,y) y.*(1 - y);
y0_c = 0.5;
y_exact_c = @(t) exp(t) ./ (1 + exp(t));

fprintf('Test (c)\n');
for i = 1:length(methods)
    [t, y, err] = methods{i}(t0, y0_c, h, N, f_c, y_exact_c);
    figure(5); hold on;
    plot(t, y, 'DisplayName', method_names{i});
    figure(6); hold on;
    plot(t, err, 'DisplayName', method_names{i});
end
figure(5);
plot(t_grid, y_exact_c(t_grid), 'k--', 'DisplayName', 'Esatta');
title('Soluzioni Test (c)');
xlabel('t'); ylabel('y');
legend('show'); hold off;

figure(6);
title('Errore assoluto Test (c)');
xlabel('t'); ylabel('|Errore|');
legend('show'); hold off;
\end{lstlisting}

\vspace{2em}

\newpage
\section*{(d) Problema scalare:}  
\[
\dot{y} = 16 y (1 - y), \quad y(0) = \frac{1}{1024}
\]
Soluzione esatta:  
\[
y(t) = \frac{e^{16t - \log(1023)}}{1 + e^{16t - \log(1023)}}
\]

\noindent
\textbf{Codice MATLAB:}
\begin{lstlisting}
f_d = @(t,y) 16*y.*(1 - y);
y0_d = 1/1024;
y_exact_d = @(t) exp(16*t - log(1023)) ./ (1 + exp(16*t - log(1023)));

fprintf('Test (d)\n');
for i = 1:length(methods)
    [t, y, err] = methods{i}(t0, y0_d, h, N, f_d, y_exact_d);
    figure(7); hold on;
    plot(t, y, 'DisplayName', method_names{i});
    figure(8); hold on;
    plot(t, err, 'DisplayName', method_names{i});
end
figure(7);
plot(t_grid, y_exact_d(t_grid), 'k--', 'DisplayName', 'Esatta');
title('Soluzioni Test (d)');
xlabel('t'); ylabel('y');
legend('show'); hold off;

figure(8);
title('Errore assoluto Test (d)');
xlabel('t'); ylabel('|Errore|');
legend('show'); hold off;
\end{lstlisting}

\newpage
\section*{(e) Oscillatore armonico (sistema di ordine 2):}  
\[
\begin{cases}
\dot{y}_1 = y_2, \\
\dot{y}_2 = -\omega^2 y_1,
\end{cases}
\quad y_1(0) = x_0, \quad y_2(0) = y_0 \omega
\]

Soluzione esatta:  
\[
y_1(t) = x_0 \cos(\omega t) + y_0 \sin(\omega t)
\]

\noindent
\textbf{Parametri e definizione MATLAB:}
\begin{lstlisting}
omega = 2;
f_e = @(t,Y) [Y(2); -omega^2*Y(1)];
y0_e = [1; 0]; % x0 = 1, y0 = 0
y_exact_e = @(t) [1]*cos(omega*t) + [0]*sin(omega*t);
% Il vettore soluzione è solo y1(t)! y2(t) si può calcolare ma ci interessa y1
\end{lstlisting}

\vspace{1em}

\noindent
\textbf{Risolvere con tutti i metodi (notare che y0 è vettore 2D):}
\begin{lstlisting}
fprintf('Test (e) Oscillatore armonico\n');
for i = 1:length(methods)
    [t, Y, err] = methods{i}(t0, y0_e, h, N, f_e, []);
    figure(9); hold on;
    plot(t, Y(1,:), 'DisplayName', method_names{i}); % y1(t)
end
% soluzione esatta y1(t)
plot(t_grid, y_exact_e(t_grid), 'k--', 'DisplayName', 'Esatta');
title('Soluzioni Test (e) Oscillatore armonico');
xlabel('t'); ylabel('y_1(t)');
legend('show'); hold off;
\end{lstlisting}

\section*{(f) Lotka-Volterra (sistema di ordine 2):}  
\[
\begin{cases}
\dot{y}_1 = \alpha y_1 - \beta y_1 y_2, \\
\dot{y}_2 = -\gamma y_2 + \delta y_1 y_2,
\end{cases}
\quad y_1(0) = 5, \quad y_2(0) = 3,
\]
con \(\alpha=1, \beta=0.2, \gamma=1.5, \delta=0.75\).

\noindent
\textbf{Definizione MATLAB:}
\begin{lstlisting}
alpha = 1; beta = 0.2; gamma = 1.5; delta = 0.75;
f_f = @(t,Y) [alpha*Y(1) - beta*Y(1)*Y(2); -gamma*Y(2) + delta*Y(1)*Y(2)];
y0_f = [5; 3];

fprintf('Test (f) Lotka-Volterra\n');
for i = 1:length(methods)
    [t, Y, ~] = methods{i}(t0, y0_f, h, N, f_f, []);
    figure(10); hold on;
    plot(t, Y(1,:), 'DisplayName', [method_names{i} ' preda']);
    plot(t, Y(2,:), 'DisplayName', [method_names{i} ' predatore']);
end
title('Soluzioni Test (f) Lotka-Volterra');
xlabel('t'); ylabel('Popolazioni');
legend('show'); hold off;
\end{lstlisting}

\section*{(g) Problema matriciale:}  
\[
\dot{Y} = A Y, \quad Y(0) = Y_0,
\]
dove  
\[
A = \begin{bmatrix} 0 & 1 \\ -4 & 0 \end{bmatrix}, \quad Y_0 = \begin{bmatrix} 1 \\ 0 \end{bmatrix}.
\]

\noindent
Soluzione esatta:  
\[
Y(t) = \begin{bmatrix}
\cos(2t) & \frac{1}{2} \sin(2t) \\
-2 \sin(2t) & \cos(2t)
\end{bmatrix} Y_0.
\]

\noindent
\textbf{Definizione MATLAB:}
\begin{lstlisting}
A = [0 1; -4 0];
f_g = @(t,Y) A*Y;
y0_g = [1; 0];

fprintf('Test (g) matrice\n');
for i = 1:length(methods)
    [t, Y, ~] = methods{i}(t0, y0_g, h, N, f_g, []);
    figure(11); hold on;
    plot(t, Y(1,:), 'DisplayName', method_names{i});
    plot(t, Y(2,:), 'DisplayName', method_names{i});
end
% soluzione esatta y1(t)
Y_exact_g = @(t) [cos(2*t); -2*sin(2*t)];
plot(t_grid, Y_exact_g(t_grid), 'k--', 'DisplayName', 'Esatta');
title('Soluzioni Test (g) matrice');
xlabel('t'); ylabel('Y(t)');
legend('show'); hold off;
\end{lstlisting}
\section*{Filo elastico}
\addcontentsline{toc}{section}{Filo elastico}
Il problema seguente rappresenta il comportamento di un filo elastico (o una corda) fissato alle estremità e sottoposto a una forza esterna. Questo problema ha molte applicazioni:
\begin{itemize}
  \item Una corda tesa che vibra
  \item Un filo elettrico o un cavo sospeso tra due punti
  \item Una molla unidimensionale sottoposta a forze
\end{itemize}

L'equazione differenziale che descrive il fenomeno è:
\[
- \frac{d^2 u}{dx^2} = f(x) \quad \text{in } (0,1)
\]
con condizioni al bordo:
\[
u(0) = 0, \quad u(1) = 0
\]

Dove:
\begin{itemize}
  \item $u(x)$ rappresenta lo spostamento verticale del filo dalla posizione di equilibrio
  \item $f(x)$ rappresenta la forza esterna applicata per unità di lunghezza
\end{itemize}

Fisicamente, il termine $-\frac{d^2 u}{dx^2}$ rappresenta la tensione interna del filo, che in equilibrio deve bilanciare la forza esterna $f(x)$. In altre parole:
\[
- \frac{d^2 u}{dx^2} = f(x) \quad \Rightarrow \quad \text{la curvatura del filo è proporzionale alla forza applicata (con segno opposto)}
\]

La funzione $u(x)$ descrive la deformazione del filo. Se $\frac{d^2 u}{dx^2} \approx 0$, allora il filo è quasi rettilineo (curvatura nulla). Se $\frac{d^2 u}{dx^2}$ è grande, la curvatura è accentuata.

Per risolvere numericamente questo problema, si discretizza l'equazione trasformandola da continua a discreta:

\begin{enumerate}
  \item Si divide l'intervallo $(0,1)$ in $n+1$ segmenti uguali, ottenendo $n$ punti interni.
  \item Il passo di discretizzazione è $h = \frac{1}{n+1}$.
  \item Si definiscono $u_0, u_1, \dots, u_{n+1}$ i valori della funzione nei punti discreti.
  \item Si approssima la derivata seconda nel punto $x_i$ con la formula:
  \[
  \frac{d^2 u}{dx^2} \bigg|_{x_i} \approx \frac{u_{i+1} - 2u_i + u_{i-1}}{h^2}
  \]
  Questa formula deriva dallo sviluppo di Taylor:
  \[
  u(x_{i+1}=x_i +h) = u(x_i) +h u'(x_i) + \frac{h^2}{2} u''(x_i) + o(h^3)
  \]
  Sommando $u(x_i+h) + u(x_i-h)$:
  \[
  u_{i+1} + u_{i-1} = 2u_i + h^2 u''(x_i) + o(h^3) \quad \Rightarrow \quad u''(x_i) \approx \frac{u_{i+1} - 2u_i + u_{i-1}}{h^2}
  \]
\end{enumerate}

Sostituendo nella PDE originale si ottiene:
\[
- \frac{u_{i+1} - 2u_i + u_{i-1}}{h^2} = f(x_i)
\quad \Rightarrow \quad
u_{i+1} - 2u_i + u_{i-1} = h^2 f(x_i)
\]

Questa discretizzazione porta a un sistema lineare del tipo:
\[
A \mathbf{u} = \mathbf{b}
\]
dove:
\begin{itemize}
  \item $A = A_h$ è una matrice tridiagonale $n \times n$
  \item $\mathbf{u} = (u_1, u_2, \dots, u_n)^T$ è il vettore incognito delle deformazioni
  \item $\mathbf{b} = (h^2 f(x_1), h^2 f(x_2), \dots, h^2 f(x_n))^T$ è il vettore dei termini noti
\end{itemize}

La matrice $A$ ha la forma:
\[
A = \frac{1}{h^2}
\begin{bmatrix}
2 & -1 & 0 & \cdots & 0 \\
-1 & 2 & -1 & \ddots & \vdots \\
0 & -1 & 2 & \ddots & 0 \\
\vdots & \ddots & \ddots & \ddots & -1 \\
0 & \cdots & 0 & -1 & 2
\end{bmatrix}
\]

È tridiagonale perché ogni equazione coinvolge solo $u_{i-1}$, $u_i$, $u_{i+1}$. Ogni riga del sistema lineare corrisponde a:
\[
- u_{i-1} + 2u_i - u_{i+1} = h^2 f(x_i)
\]

\begin{lstlisting}[language=Matlab]
clear; close all;

% Funzione per risolvere il problema: discretizzazione e risoluzione
function [x, u, u_analitica] = risolvi_filo(n, f_handle, f_analitica)
    h = 1/n; N = n-1;           % Passo della griglia e numero di incognite
    x = linspace(0, 1, n+1);    % Griglia di punti [x0, x1, ..., xn]
    A = sparse(N, N);           % Matrice sparsa per il sistema
    for i = 1:N
        A(i,i) = 2/h^2;         % Diagonale principale
        if i > 1, A(i,i-1) = -1/h^2; end  % Legami sotto-diagonale
        if i < N, A(i,i+1) = -1/h^2; end  % Legami sopra-diagonale
    end
    b = f_handle(x(2:n))';      % Vettore dei termini noti [f(x1), ..., f(xn-1)]
    U = A\b;                    % Risoluzione del sistema lineare
    u = [0; U; 0];              % Soluzione completa (includendo i bordi u(0) = 0, u(1) = 0)
    u_analitica = isempty(f_analitica) ? [] : f_analitica(x); % Soluzione analitica (se disponibile)
end

% Caso 1: f(x) = -1
f1 = @(x) -1; u1_analitica = @(x) (x-x.^2)/2;  % Soluzione analitica u(x) = (x - x^2) / 2
n_values = [10, 20, 50, 100];  % Diverse risoluzioni per testare la convergenza

% Grafico della soluzione numerica per vari n
figure; hold on;
for n = n_values
    [x, u, u_analitica] = risolvi_filo(n, f1, u1_analitica);  % Calcola la soluzione numerica
    plot(x, u, 'o-', 'DisplayName', sprintf('n = %d', n));    % Mostra la soluzione numerica
    errore = max(abs(u - u_analitica));  % Calcola l'errore massimo tra numerica e analitica
    fprintf('f(x) = -1, n = %d, Errore massimo = %e\n', n, errore);  % Mostra errore
end
% Soluzione analitica per confronto
plot(linspace(0, 1, 1000), u1_analitica(linspace(0, 1, 1000)), 'k-', 'LineWidth', 1.5, 'DisplayName', 'Soluzione analitica');
title('Soluzione con f(x) = -1'); xlabel('x'); ylabel('u(x)'); legend; grid on;

% Caso 2: f(x) = sin(pi*x)
f2 = @(x) sin(pi*x); u2_analitica = @(x) sin(pi*x)/(pi^2);  % Soluzione analitica u(x) = sin(pi*x) / pi^2
figure; hold on;
for n = n_values
    [x, u, u_analitica] = risolvi_filo(n, f2, u2_analitica);  % Calcola la soluzione numerica
    plot(x, u, 'o-', 'DisplayName', sprintf('n = %d', n));    % Mostra la soluzione numerica
    errore = max(abs(u - u_analitica));  % Calcola l'errore massimo
    fprintf('f(x) = sin(pi*x), n = %d, Errore massimo = %e\n', n, errore);
end
% Soluzione analitica per confronto
plot(linspace(0, 1, 1000), u2_analitica(linspace(0, 1, 1000)), 'k-', 'LineWidth', 1.5, 'DisplayName', 'Soluzione analitica');
title('Soluzione con f(x) = sin(\pi x)'); xlabel('x'); ylabel('u(x)'); legend; grid on;

% Caso 3: f(x) = 10*sin(5*pi*x)
f3 = @(x) 10*sin(5*pi*x); u3_analitica = @(x) 10*sin(5*pi*x)/(25*pi^2);  % Soluzione analitica u(x) = 10*sin(5*pi*x)/(25*pi^2)
figure; hold on;
for n = n_values
    [x, u, u_analitica] = risolvi_filo(n, f3, u3_analitica);  % Calcola la soluzione numerica
    plot(x, u, 'o-', 'DisplayName', sprintf('n = %d', n));    % Mostra la soluzione numerica
    errore = max(abs(u - u_analitica));  % Calcola l'errore massimo
    fprintf('f(x) = 10*sin(5*pi*x), n = %d, Errore massimo = %e\n', n, errore);
end
% Soluzione analitica per confronto
plot(linspace(0, 1, 1000), u3_analitica(linspace(0, 1, 1000)), 'k-', 'LineWidth', 1.5, 'DisplayName', 'Soluzione analitica');
title('Soluzione con f(x) = 10*sin(5\pi x)'); xlabel('x'); ylabel('u(x)'); legend; grid on;

% Analisi della convergenza: calcolo dell'errore per diversi valori di n
n_conv = [10, 20, 40, 80, 160];  % Diversi valori di n per testare la convergenza
errors1 = zeros(length(n_conv), 1); errors2 = zeros(length(n_conv), 1); errors3 = zeros(length(n_conv), 1);
for i = 1:length(n_conv)
    n = n_conv(i);
    % Calcolo dell'errore massimo per ciascun caso
    errors1(i) = max(abs(risolvi_filo(n, f1, u1_analitica) - u1_analitica(n)));
    errors2(i) = max(abs(risolvi_filo(n, f2, u2_analitica) - u2_analitica(n)));
    errors3(i) = max(abs(risolvi_filo(n, f3, u3_analitica) - u3_analitica(n)));
end

% Grafico della convergenza
figure; loglog(n_conv, errors1, 'bo-', n_conv, errors2, 'ro-', n_conv, errors3, 'go-', n_conv, (1./n_conv).^2, 'k--');
legend('f(x) = -1', 'f(x) = sin(\pi x)', 'f(x) = 10*sin(5\pi x)', 'O(h^2)'); % Linea di convergenza teorica
xlabel('Numero di intervalli (n)'); ylabel('Errore massimo'); title('Convergenza del metodo'); grid on;
\end{lstlisting}
\newpage
\section*{Il sistema di Lotka-Volterra: modello preda-predatore}
\addcontentsline{toc}{section}{Il sistema di Lotka-Volterra: modello preda-predatore}

Il sistema di Lotka-Volterra è un modello matematico che descrive l’interazione tra due specie animali: una \emph{preda} e un \emph{predatore}. Lo scopo è capire come varia nel tempo la quantità di prede e predatori in un ecosistema.

Il modello si basa su due equazioni differenziali ordinarie che rappresentano l’evoluzione nel tempo delle popolazioni:

\[
\begin{cases}
\frac{dx}{dt} = \alpha x - \beta x y \\
\frac{dy}{dt} = \delta x y - \gamma y
\end{cases}
\]

dove:
\begin{itemize}
    \item \(x(t)\) è il numero di prede al tempo \(t\),
    \item \(y(t)\) è il numero di predatori al tempo \(t\),
    \item \(\alpha, \beta, \gamma, \delta > 0\) sono parametri che descrivono il comportamento delle specie.
\end{itemize}

Queste equazioni si interpretano così:

\begin{itemize}
    \item \(\frac{dx}{dt} = \alpha x - \beta x y\):
    \begin{itemize}
        \item Il termine \(\alpha x\) rappresenta la crescita naturale delle prede, che aumentano proporzionalmente alla loro quantità con tasso \(\alpha\).
        \item Il termine \(- \beta x y\) rappresenta la diminuzione delle prede dovuta alla predazione: più prede e più predatori ci sono, più le prede vengono mangiate.
    \end{itemize}
    
    \item \(\frac{dy}{dt} = \delta x y - \gamma y\):
    \begin{itemize}
        \item Il termine \(\delta x y\) rappresenta la crescita dei predatori, che dipende da quante prede riescono a catturare (più prede, più cibo, più predatori nascono).
        \item Il termine \(- \gamma y\) rappresenta la mortalità naturale dei predatori, che diminuiscono proporzionalmente alla loro quantità con tasso \(\gamma\).
    \end{itemize}
\end{itemize}

Perché usare equazioni differenziali?

\[
\frac{dx}{dt}, \frac{dy}{dt}
\]

sono derivate rispetto al tempo e indicano la variazione istantanea della popolazione. Le equazioni differenziali modellano quindi il cambiamento continuo nel tempo, il modo più naturale per descrivere fenomeni biologici in evoluzione.

Il modello di Lotka-Volterra è molto semplice ma cattura bene il comportamento ciclico osservato in molti sistemi reali, dove la popolazione delle prede cresce, poi aumenta anche quella dei predatori che però, crescendo troppo, fanno diminuire le prede, causando poi un calo dei predatori e ricominciando il ciclo.

Questo sistema è uno strumento base per studiare le dinamiche di ecosistemi con interazioni tra specie, e le sue soluzioni offrono informazioni utili sia in biologia che in ecologia.
\newpage
I parametri \(\alpha, \beta, \gamma, \delta\) nel sistema di Lotka-Volterra sono fondamentali perché definiscono il comportamento delle popolazioni:

\begin{itemize}
    \item \(\alpha\) è il tasso di crescita naturale delle prede in assenza di predatori. Un valore più alto implica una crescita più rapida delle prede.
    \item \(\beta\) misura l’efficacia della predazione, cioè la velocità con cui i predatori catturano le prede. Un valore alto di \(\beta\) riduce rapidamente la popolazione delle prede.
    \item \(\gamma\) rappresenta il tasso di mortalità naturale dei predatori in assenza di prede.
    \item \(\delta\) indica l’efficienza con cui i predatori trasformano le prede catturate in nuovi predatori.
\end{itemize}

\vspace{1em}

\textbf{Esempio numerico semplice:}

Supponiamo di avere

\[
\alpha = 1.1, \quad \beta = 0.4, \quad \gamma = 0.4, \quad \delta = 0.1,
\]

con condizioni iniziali

\[
x(0) = 40 \quad \text{(prede)}, \quad y(0) = 9 \quad \text{(predatori)}.
\]

Questo significa che:

\begin{itemize}
    \item Le prede crescono velocemente senza predatori (\(\alpha=1.1\)).
    \item I predatori sono abbastanza efficaci nella caccia (\(\beta=0.4\)).
    \item I predatori muoiono lentamente se non trovano prede (\(\gamma=0.4\)).
    \item La riproduzione dei predatori è più lenta rispetto alla crescita delle prede (\(\delta=0.1\)).
\end{itemize}

Simulando il sistema con questi valori, si osserva una dinamica oscillante, in cui la popolazione delle prede cresce prima e quella dei predatori segue con un ritardo, dando luogo a cicli regolari di crescita e decrescita.

\vspace{1em}

\textbf{Applicazioni del modello Lotka-Volterra:}

\begin{itemize}
    \item \textit{Ecologia:} Per studiare la dinamica di ecosistemi con interazioni predatore-preda.
    \item \textit{Epidemiologia:} Per modellare la diffusione di malattie tra popolazioni suscettibili e infette.
    \item \textit{Economia:} Per rappresentare sistemi competitivi o di cooperazione tra agenti economici.
    \item \textit{Chimica:} Per descrivere reazioni chimiche autocatalitiche o dinamiche di popolazioni molecolari.
\end{itemize}

Il modello, pur nella sua semplicità, fornisce intuizioni utili sulla complessità dei sistemi interattivi e serve come base per modelli più sofisticati in diversi campi.

\newpage
\section*{Esercizi}
\addcontentsline{toc}{section}{Esercizi}

\fbox{%
Sia
\[
A =
\begin{bmatrix}
-2 & 4 & 0 & 0 \\
3 & 1 & 0 & 1 \\
2 & -1 & 2 & 1 \\
-1 & 0 & 0 & 1
\end{bmatrix}
\]
Dire se $\lambda = 7 + i$ può essere un autovalore di $A$.%
}

\vspace{1em}

\noindent
Calcoliamo i dischi di Gershgorin riga per riga:

\begin{itemize}
  \item Riga 1: centro $a_{11} = -2$, raggio $R_1 = |4| + |0| + |0| = 4$ \\
  $\Rightarrow D_1 = \{ z \in \mathbb{C} : |z + 2| \le 4 \}$
  
  \item Riga 2: centro $a_{22} = 1$, raggio $R_2 = |3| + |0| + |1| = 4$ \\
  $\Rightarrow D_2 = \{ z \in \mathbb{C} : |z - 1| \le 4 \}$
  
  \item Riga 3: centro $a_{33} = 2$, raggio $R_3 = |2| + |{-1}| + |1| = 4$ \\
  $\Rightarrow D_3 = \{ z \in \mathbb{C} : |z - 2| \le 4 \}$
  
  \item Riga 4: centro $a_{44} = 1$, raggio $R_4 = |{-1}| + |0| + |0| = 1$ \\
  $\Rightarrow D_4 = \{ z \in \mathbb{C} : |z - 1| \le 1 \}$
\end{itemize}

\vspace{1em}

\noindent
Verifichiamo se $\lambda = 7 + i$ appartiene a qualcuno dei dischi:

\[
\begin{aligned}
|7 + i - (-2)| &= \sqrt{(9)^2 + (1)^2} = \sqrt{82} > 4 \\
|7 + i - 1| &= \sqrt{(6)^2 + (1)^2} = \sqrt{37} > 4 \\
|7 + i - 2| &= \sqrt{(5)^2 + (1)^2} = \sqrt{26} > 4 \\
|7 + i - 1| &= \sqrt{(6)^2 + (1)^2} = \sqrt{37} > 1
\end{aligned}
\]

\vspace{1em}

\noindent
Poiché $\lambda = 7 + i \notin \bigcup_i D_i \Rightarrow \lambda \notin \text{spec}(A)$, possiamo escludere $\lambda$ come autovalore.

\vspace{1em}

\noindent
\textbf{Conclusione:} $\lambda = 7 + i$ non può essere un autovalore di $A$ (escluso a priori tramite Gershgorin).


\vspace{1em}
\noindent
\fbox{%
\begin{minipage}{\widthof{Siano $A = \begin{bmatrix} 1 & 2 & 3 \\ 1 & 0 & 0 \\ 2 & 4 & 2 \end{bmatrix}, \quad L = \begin{bmatrix} 1 & 0 & 0 \\ 1 & 1 & 0 \\ 2 & 1 & 1 \end{bmatrix}, \quad U = \begin{bmatrix} 1 & 2 & 3 \\ 0 & c & 2 \\ 0 & 0 & c \end{bmatrix}$}}
Siano
$
A =
\begin{bmatrix}
1 & 2 & 3 \\
1 & 0 & 0 \\
2 & 4 & 2
\end{bmatrix}, \quad
L =
\begin{bmatrix}
1 & 0 & 0 \\
1 & 1 & 0 \\
2 & 1 & 1
\end{bmatrix}, \quad
U =
\begin{bmatrix}
1 & 2 & 3 \\
0 & c & 2 \\
0 & 0 & c
\end{bmatrix}
$\\[0.5em]
Trovare il parametro $c$, se esiste, tale che $A = LU$.
\end{minipage}%
}


\vspace{1em}
\[
L =
\begin{bmatrix}
\colorbox{red!20}{1} & \colorbox{red!20}{0} & \colorbox{red!20}{0} \\
\colorbox{blue!20}{1} & \colorbox{blue!20}{1} & \colorbox{blue!20}{0} \\
\colorbox{green!20}{2} & \colorbox{green!20}{1} & \colorbox{green!20}{1}
\end{bmatrix}, \quad
U =
\begin{bmatrix}
\colorbox{violet!20}{1} & \colorbox{yellow!20}{2} & \colorbox{cyan!20}{3} \\
\colorbox{violet!20}{0} & \colorbox{yellow!20}{c} & \colorbox{cyan!20}{2} \\
\colorbox{violet!20}{0} & \colorbox{yellow!20}{0} & \colorbox{cyan!20}{c}
\end{bmatrix}
\]

\vspace{1em}

Riga 1 del prodotto \( LU \):
\[
\begin{aligned}
(LU)_{1,1} &= \colorbox{red!20}{1} \cdot \colorbox{violet!20}{1} + \colorbox{red!20}{0} \cdot \colorbox{violet!20}{0} + \colorbox{red!20}{0} \cdot \colorbox{violet!20}{0} = 1 \\
(LU)_{1,2} &= \colorbox{red!20}{1} \cdot \colorbox{yellow!20}{2} + \colorbox{red!20}{0} \cdot \colorbox{yellow!20}{c} + \colorbox{red!20}{0} \cdot \colorbox{yellow!20}{0} = 2 \\
(LU)_{1,3} &= \colorbox{red!20}{1} \cdot \colorbox{cyan!20}{3} + \colorbox{red!20}{0} \cdot \colorbox{cyan!20}{2} + \colorbox{red!20}{0} \cdot \colorbox{cyan!20}{c} = 3 \\
\end{aligned}
\]

\vspace{1em}

Riga 2 del prodotto \( LU \):
\[
\begin{aligned}
(LU)_{2,1} &= \colorbox{blue!20}{1} \cdot \colorbox{violet!20}{1} + \colorbox{blue!20}{1} \cdot \colorbox{violet!20}{0} + \colorbox{blue!20}{0} \cdot \colorbox{violet!20}{0} = 1 \\
(LU)_{2,2} &= \colorbox{blue!20}{1} \cdot \colorbox{yellow!20}{2} + \colorbox{blue!20}{1} \cdot \colorbox{yellow!20}{c} + \colorbox{blue!20}{0} \cdot \colorbox{yellow!20}{0} = 2 + c \\
(LU)_{2,3} &= \colorbox{blue!20}{1} \cdot \colorbox{cyan!20}{3} + \colorbox{blue!20}{1} \cdot \colorbox{cyan!20}{2} + \colorbox{blue!20}{0} \cdot \colorbox{cyan!20}{c} = 5 \\
\end{aligned}
\]

\vspace{1em}

Riga 3 del prodotto \( LU \):
\[
\begin{aligned}
(LU)_{3,1} &= \colorbox{green!20}{2} \cdot \colorbox{violet!20}{1} + \colorbox{green!20}{1} \cdot \colorbox{violet!20}{0} + \colorbox{green!20}{1} \cdot \colorbox{violet!20}{0} = 2 \\
(LU)_{3,2} &= \colorbox{green!20}{2} \cdot \colorbox{yellow!20}{2} + \colorbox{green!20}{1} \cdot \colorbox{yellow!20}{c} + \colorbox{green!20}{1} \cdot \colorbox{yellow!20}{0} = 4 + c \\
(LU)_{3,3} &= \colorbox{green!20}{2} \cdot \colorbox{cyan!20}{3} + \colorbox{green!20}{1} \cdot \colorbox{cyan!20}{2} + \colorbox{green!20}{1} \cdot \colorbox{cyan!20}{c} = 8 + c \\
\end{aligned}
\]
\vspace{1em}
\[
LU =
\begin{bmatrix}
1 & 2 & 3 \\
1 & c + 2 & 5 \\
2 & c + 4 & c + 8
\end{bmatrix}
=
\begin{bmatrix}
1 & 2 & 3 \\
1 & 0 & 0 \\
2 & 4 & 2
\end{bmatrix}
\Rightarrow 5 = 0 \Rightarrow   A \not= LU \text{   } \forall c \in \mathbb{R}
\]


\newpage
\noindent
\fbox{%
\begin{minipage}{\textwidth}
Sia $v \in \mathbb{R}^n$. Consideriamo la matrice $V = vv^T$. Quali delle seguenti affermazioni sono corrette?\\[0.5em]
A) $V$ è invertibile\\
B) $V$ è una matrice simmetrica\\
C) Non è possibile determinare il rango di $V$ senza fare calcoli\\
D) $V$ è una matrice di rango 1
\end{minipage}%
}

\vspace{1em}
\begin{tabular}{@{}l@{\hspace{1em}}p{0.85\textwidth}}
A) &
$v = \begin{bmatrix}
0 \\
0 \\
\vdots \\
0
\end{bmatrix}
\Rightarrow V = vv^T =
\begin{bmatrix}
0 & \cdots & 0 \\
\vdots & \ddots & \vdots \\
0 & \cdots & 0
\end{bmatrix}
\Rightarrow \det(V) = 0
\Rightarrow \textit{V non è invertibile.}$

In generale, tranne nel caso in cui $n = 1$ e $v \ne 0$, la matrice $V = vv^T$ ha rango 1, quindi non è invertibile (anche perchè l'i-esima riga i = $v_i\cdot v^T$ ovvero tutte le righe sono combinazioni lineari di v)
\\[1em]

B) &
$V^T = (vv^T)^T = (v^T)^Tv^T = vv^T = V \Rightarrow$ \textit{$V$ è una matrice simmetrica.}
\\[1em]

C) &
$Vx = (vv^T)x = v(v^T x) = (v^T x)v \quad \forall x \in \mathbb{R}^n
\Rightarrow Im(V) \subseteq \underbrace{span(v)}_{\mathclap{\{\alpha v \forall \alpha \in \mathbb{R}\}}}
\Rightarrow rank(V) \leq 1$

\\[1em]

D) &
Se $v \ne \begin{bmatrix}
0 \\
0 \\
\vdots \\
0
\end{bmatrix}$ allora $\operatorname{rank}(V) = 1$;

Se $v = \begin{bmatrix}
0 \\
0 \\
\vdots \\
0
\end{bmatrix}$ allora $\operatorname{rank}(V) = 0$.
\end{tabular}

\vspace{1em}
\noindent
\fbox{%
\begin{minipage}{\textwidth}
Sia $A \in \mathbb{R}^{n \times n}$ con fattorizzazione SVD $A = U \Sigma V^T$. Scrivere la fattorizzazione SVD di $A A^T$ e $A^T A$. È possibile calcolare facilmente i valori singolari di $A A^T$?
\end{minipage}%
}

\vspace{2em}
\noindent

$A A^T = U \Sigma V^T V \Sigma^T U^T = U \Sigma \Sigma^T U^T$\\
\noindent

$A^T A = V \Sigma^T U^T U \Sigma V^T = V \Sigma^T \Sigma V^T$

\vspace{1em}

\noindent
I valori singolari di \(A\) sono le radici quadrate degli autovalori di \(A^T A\) o \(A A^T\),
quindi è possibile calcolarli facilmente.

\newpage
\noindent
\fbox{%
\begin{minipage}{\textwidth}
Siano dati i nodi $x_0 = -1$, $x_1 = 0$, $x_2 = 1$ e $x_3 = 1$, e la funzione $f(x) = x^4 - x^2$. Interpolare i dati assegnati con un polinomio $P$ di grado adatto. Calcolare la formula dell'errore $P(x) - f(x)$. Aggiungere ora il punto $x_4 = 1$. Si ottiene il nuovo polinomio interpolatore $Q(x)$. Qual è la stima dell'errore $Q(x) - f(x)$, senza calcolare esplicitamente $Q(x)$? Riportare $P(x)$ e le due stime di errore.
\end{minipage}%
}

\vspace{1em}
\noindent

\[
P(x) = \sum_{i=0}^3 f(x_i) L_i(x), \quad L_i(x) = \prod_{j=0, j \neq i}}^3 \frac{x - x_j}{x_i - x_j}
\]

\[
f(x_0)=0, \quad f(x_1)=0, \quad f(x_2)=0, \quad f(x_3)=0
\Rightarrow
 P(x) = 0 \]

\[
\Rightarrow E(x) = -f(x) = - (x^4 - x^2)
\]

L'errore dell'interpolazione polinomiale di grado \(3\) su nodi \(x_0, \ldots, x_3\) è dato da:
\[
E(x) = \frac{f^{(4)}(\xi)}{4!} \prod_{i=0}^3 (x - x_i), \quad \xi \in [x_0, x_3].
\]
\[ \text{Nel nostro caso }
f(x) = x^4 - x^2
\Rightarrow 
f^{(4)}(x) = 24.
\Rightarrow
E(x) = \frac{24}{4!} \prod_{i=0}^3 (x - x_i) = \prod_{i=0}^3 (x - x_i),
\]

I nodi sono \(x_0 = -1, x_1 = 0, x_2 = 1, x_3 = 1\), quindi

\[
\prod_{i=0}^3 (x - x_i) = (x + 1) \cdot x \cdot (x - 1)^2.
\]

\[
P(x) - f(x) = (x+1) \cdot x \cdot (x-1)^2.
\]

Aggiungendo \( x_4 = 1 \), otteniamo il polinomio interpolatore \( Q(x) \) di grado \(\leq 4\) sui nodi
\[
x_0 = -1, \quad x_1 = 0, \quad x_2 = 1, \quad x_3 = 1, \quad x_4 = 1.
\]

L'errore dell'interpolazione polinomiale di grado \(4\) su nodi \(x_0, \ldots, x_4\) è dato da:\[
E(x) = Q(x) - f(x) = \frac{f^{(5)}(\zeta)}{5!} \prod_{i=0}^4 (x - x_i), \quad \zeta \in [x_0, x_4].
\]
\[
f^{(5)}(x) = 0 \text{   } \forall x \in \mathbb{R}\implies E(x) = 0.
\]

\vspace{1em}
\noindent
\fbox{%
\begin{minipage}{\textwidth}
Sia $A$ una matrice tridiagonale $n \times n$ definita da:
\[
A =
\begin{bmatrix}
c & 1 & 0 & \cdots & 0 \\
1 & c & 1 & \ddots & \vdots \\
0 & 1 & \ddots & \ddots & 0 \\
\vdots & \ddots & \ddots & c & 1 \\
0 & \cdots & 0 & 1 & c
\end{bmatrix}
\]
\begin{enumerate}
  \item Determinare i valori di $c$ per i quali il metodo di Jacobi converge per A.
  \item Scrivere una funzione che, dato un vettore $b$, il parametro $c$ e le dimensioni $n$ di $A$, applica il metodo di Jacobi per trovare la soluzione del sistema $Ax = b$. La funzione deve ritornare errore se $c$ non garantisce la convergenza del metodo di Jacobi. L'iterazione si interrompe se $\|Ax^{(k)} - b\|_2 < 10^{-5}$, o dopo 1000 iterazioni.
  \item Scrivere il programma principale che legge i parametri $c$ ed $n$ da tastiera (o li inizializza), riempie il vettore $b$ di numeri pseudocasuali distribuiti uniformemente su $(-1,4)$, chiama la funzione e scrive la soluzione trovata, il numero di iterazioni necessario ad arrivare a convergenza ed eventuali messaggi di errore.
\end{enumerate}
\end{minipage}%
}

\vspace{1em}
\noindent

\begin{enumerate}
  \item \textbf{Condizione di convergenza per Jacobi:} $a_{ii}\geq\sum_{j=1}^na_{ij}\forall i \in \{1 \dots n \} \land \exists i \in \{1 \dots n \}$ \text{ t.c. }$a_{ii}>\sum_{j=1}^na_{ij}\Rightarrow $\text{|c|>2}

  \item \textbf{Funzione metodo di Jacobi (pseudocodice):}
\begin{verbatim}
  funzione Ax(x, c, n):
    r = vettore di dimensione n
    r[1] = c * x[1] + x[2]
    per i = 2 fino a n-1:
        r[i] = x[i-1] + c * x[i] + x[i+1]
    r[n] = x[n-1] + c * x[n]
    ritorna r
  funzione Jacobi(b, c, n):
    se |c| <= 2 allora
        ritorna "Errore: metodo non converge"
    x = vettore zero di dimensione n
    x_new, k = x, 0
    mentre k <= 1000:
        per i da 1 a n:
        somma = 0
        se i > 1 allora somma = somma + x[i - 1]
        se i < n allora somma = somma + x[i + 1]
        x_new[i] = (b[i] - somma) / c
        residuo = norma_2(Ax(x_new, c, n) - b)
        se residuo < 1e-5 allora
            ritorna x_new, k+1
        x = x_new, k = k + 1
    ritorna x_new, k
  \end{verbatim}
  \item \textbf{Programma principale (pseudocodice):}
  \begin{verbatim}
  leggi o inizializza c, n
  genera vettore b con valori uniformi su (-1, 4)
  risultato = Jacobi(b, c, n)
  se risultato è messaggio di errore:
    stampa messaggio di errore
  altrimenti
    stampa soluzione x, numero iterazioni k
  \end{verbatim}
\end{enumerate}

\newpage
\noindent
\fbox{%
\begin{minipage}{\textwidth}
Siano dati i nodi $x_0 = -1$, $x_1 = -\frac{1}{2}$, $x_2 = 0$, $x_3 = \frac{1}{2}$, e la funzione $f(x) = x^4 - x^2$. Interpolare i dati assegnati con un polinomio $P$ di grado adatto. Calcolare la formula dell'errore $P(x) - f(x)$. Aggiungere ora il punto $x_4 = 1$. Si ottiene il nuovo polinomio interpolatore $Q(x)$. Qual è la stima dell'errore $Q(x) - f(x)$, senza calcolare esplicitamente $Q(x)$? Riportare $P(x)$ e le due stime di errore.
\end{minipage}%
}

\vspace{1em}
\noindent

\[
P(x) = \sum_{i=0}^3 f(x_i) L_i(x), \quad L_i(x) = \prod_{\substack{j=0 \\ j \neq i}}^3 \frac{x - x_j}{x_i - x_j}
\]

\[
\begin{aligned}
f(x_0) &= f(-1) = 1 - 1 = 0, \\
f(x_1) &= f(-\tfrac{1}{2}) = \tfrac{1}{16} - \tfrac{1}{4} = -\tfrac{3}{16}, \\
f(x_2) &= f(0) = 0, \\
f(x_3) &= f(\tfrac{1}{2}) = \tfrac{1}{16} - \tfrac{1}{4} = -\tfrac{3}{16},
\end{aligned}
\Rightarrow P(x) \text{ è un polinomio di grado } \leq 3
\]

L'errore dell'interpolazione polinomiale su 4 nodi distinti è:
\[
E(x) = f(x) - P(x) = \frac{f^{(4)}(\xi)}{4!} \prod_{i=0}^3 (x - x_i), \quad \xi \in [x_0, x_3]
\]

\[
f(x) = x^4 - x^2 \Rightarrow f^{(4)}(x) = 24, \quad \forall x
\Rightarrow
E(x) = \frac{24}{24} \prod_{i=0}^3 (x - x_i) = \prod_{i=0}^3 (x - x_i)
\]

\[
\text{Con } x_0 = -1,\ x_1 = -\tfrac{1}{2},\ x_2 = 0,\ x_3 = \tfrac{1}{2}, \Rightarrow
E(x) = (x + 1)(x + \tfrac{1}{2})x(x - \tfrac{1}{2})
\]

\vspace{1em}
\noindent
Aggiungendo \( x_4 = 1 \), otteniamo il polinomio interpolatore \( Q(x) \) di grado \( \leq 4 \) sui nodi:
\[
x_0 = -1, \quad x_1 = -\tfrac{1}{2}, \quad x_2 = 0, \quad x_3 = \tfrac{1}{2}, \quad x_4 = 1
\]

L'errore dell'interpolazione polinomiale su 5 nodi è:
\[
E(x) = f(x) - Q(x) = \frac{f^{(5)}(\zeta)}{5!} \prod_{i=0}^4 (x - x_i), \quad \zeta \in [x_0, x_4]
\]

\[
f(x) = x^4 - x^2 \Rightarrow f^{(5)}(x) = 0 \Rightarrow E(x) = 0
\]

\vspace{1em}
\noindent
Conclusione:
\[
P(x) = \text{Interpolazione Lagrange su 4 punti} \Rightarrow E(x) = (x + 1)(x + \tfrac{1}{2})x(x - \tfrac{1}{2})
\]
\[
Q(x) = \text{Interpolazione su 5 punti} \Rightarrow f \in \mathbb{P}_4 \Rightarrow E(x) = 0
\]

\newpage
\noindent
\fbox{%
\begin{minipage}{\textwidth}
Data la funzione $f(x) = e^x - 14$. Costruire un'iterazione di punto fisso per determinare la soluzione di $f(x) = 0$, del tipo $g(x) = x + \lambda(e^x - 14)$. Trovare $\lambda$ in modo tale che l'iterazione proposta sia effettivamente convergente. Notare che $\ln 14 \approx 2{,}63$.
\end{minipage}%
}
\vspace{1em}

$f(x) = e^x - 14 = 0 \Rightarrow x^* = \ln(14) \approx 2,63 \Rightarrow$ Scelgo come intervallo (a,b) = (2,3)
$\Rightarrow f \in C^1(a,b)$ 

$g(x) = x + \lambda(e^x - 14)$

$|g'(x)| = |1 + \lambda e^x| \leq |1+\lambda e^2| \leq |1+\lambda 2^2| < 1 \iff -1<1+\lambda 2^2<1 \quad \forall x \in (2,3)$

$\implies \begin{cases}-\frac{2}{2^2}<\lambda\\\lambda<0\end{cases}\iff-\frac{1}{2}<\lambda<0\implies\lambda=-\frac{1}{4}$ ad esempio

$g(x) = x - \frac{1}{4}(e^x - 14)$

\vspace{2em}
\noindent
\textbf{Iterazione di Punto Fisso}
\begin{verbatim}
g(x) = x - 1/4(e^x - 14)
x0 = 2.5
tol = 1e-15 * modulo di (g(x0))
max_iter = 1000
it = 0
finché it < max_iter e modulo di (g(x0) - x0) > tol
    x0 = g(x0)
    it = it + 1
se valore assoluto di (g(x0) - x0) <= tol
    convergenza in it iterazioni
    soluzione approssimata: x0
altrimenti
    non convergente entro max_iter iterazioni
\end{verbatim}

 \newpage
\noindent
\fbox{%
\begin{minipage}{\textwidth}
Data la matrice $Q = \begin{bmatrix} 1 & 0 & 0 \\ 0 & 2 & c \\ 0 & \frac{1}{3} & 2 \end{bmatrix}$ determinare $c$, se esiste, affinché $Q$ sia ortogonale.
\end{minipage}%
}

\vspace{1em}
Q è ortogonale $\iff Q^T Q = I$
\[
Q^T =
\begin{bmatrix}
\colorbox{red!20}{$1$} & \colorbox{red!20}{$0$} & \colorbox{red!20}{$0$} \\
\colorbox{blue!20}{$0$} & \colorbox{blue!20}{$2$} & \colorbox{blue!20}{$\frac{1}{3}$} \\
\colorbox{green!20}{$0$} & \colorbox{green!20}{$c$} & \colorbox{green!20}{$2$}
\end{bmatrix}, \quad
Q =
\begin{bmatrix}
\colorbox{violet!20}{$1$} & \colorbox{yellow!20}{$0$} & \colorbox{cyan!20}{$0$} \\
\colorbox{violet!20}{$0$} & \colorbox{yellow!20}{$2$} & \colorbox{cyan!20}{$c$} \\
\colorbox{violet!20}{$0$} & \colorbox{yellow!20}{$\frac{1}{3}$} & \colorbox{cyan!20}{$2$}
\end{bmatrix}
\]
\vspace{1em}
\noindent
Riga 1 del prodotto $Q^T Q$:
\[
\begin{aligned}
(Q^T Q)_{1,1} &= \colorbox{red!20}{$1$} \cdot \colorbox{violet!20}{$1$} + \colorbox{red!20}{$0$} \cdot \colorbox{violet!20}{$0$} + \colorbox{red!20}{$0$} \cdot \colorbox{violet!20}{$0$} = 1 \\
(Q^T Q)_{1,2} &= \colorbox{red!20}{$1$} \cdot \colorbox{yellow!20}{$0$} + \colorbox{red!20}{$0$} \cdot \colorbox{yellow!20}{$2$} + \colorbox{red!20}{$0$} \cdot \colorbox{yellow!20}{$\frac{1}{3}$} = 0 \\
(Q^T Q)_{1,3} &= \colorbox{red!20}{$1$} \cdot \colorbox{cyan!20}{$0$} + \colorbox{red!20}{$0$} \cdot \colorbox{cyan!20}{$c$} + \colorbox{red!20}{$0$} \cdot \colorbox{cyan!20}{$2$} = 0
\end{aligned}
\]
\vspace{1em}
\noindent
Riga 2 del prodotto $Q^T Q$:
\[
\begin{aligned}
(Q^T Q)_{2,1} &= \colorbox{blue!20}{$0$} \cdot \colorbox{violet!20}{$1$} + \colorbox{blue!20}{$2$} \cdot \colorbox{violet!20}{$0$} + \colorbox{blue!20}{$\frac{1}{3}$} \cdot \colorbox{violet!20}{$0$} = 0 \\
(Q^T Q)_{2,2} &= \colorbox{blue!20}{$0$} \cdot \colorbox{yellow!20}{$0$} + \colorbox{blue!20}{$2$} \cdot \colorbox{yellow!20}{$2$} + \colorbox{blue!20}{$\frac{1}{3}$} \cdot \colorbox{yellow!20}{$\frac{1}{3}$} = 4 + \frac{1}{9} = \frac{37}{9} \\
(Q^T Q)_{2,3} &= \colorbox{blue!20}{$0$} \cdot \colorbox{cyan!20}{$0$} + \colorbox{blue!20}{$2$} \cdot \colorbox{cyan!20}{$c$} + \colorbox{blue!20}{$\frac{1}{3}$} \cdot \colorbox{cyan!20}{$2$} = 2c + \frac{2}{3}
\end{aligned}
\]
\vspace{1em}
\noindent
Riga 3 del prodotto $Q^T Q$:
\[
\begin{aligned}
(Q^T Q)_{3,1} &= \colorbox{green!20}{$0$} \cdot \colorbox{violet!20}{$1$} + \colorbox{green!20}{$c$} \cdot \colorbox{violet!20}{$0$} + \colorbox{green!20}{$2$} \cdot \colorbox{violet!20}{$0$} = 0 \\
(Q^T Q)_{3,2} &= \colorbox{green!20}{$0$} \cdot \colorbox{yellow!20}{$0$} + \colorbox{green!20}{$c$} \cdot \colorbox{yellow!20}{$2$} + \colorbox{green!20}{$2$} \cdot \colorbox{yellow!20}{$\frac{1}{3}$} = 2c + \frac{2}{3} \\
(Q^T Q)_{3,3} &= \colorbox{green!20}{$0$} \cdot \colorbox{cyan!20}{$0$} + \colorbox{green!20}{$c$} \cdot \colorbox{cyan!20}{$c$} + \colorbox{green!20}{$2$} \cdot \colorbox{cyan!20}{$2$} = c^2 + 4
\end{aligned}
\]
\vspace{1em}
\[
Q^T Q =
\begin{bmatrix}
1 & 0 & 0 \\
0 & \frac{37}{9} & 2c + \frac{2}{3} \\
0 & 2c + \frac{2}{3} & c^2 + 4
\end{bmatrix}
\overbrace{=}^{\text{Imponendo Q ortogonale}}
\begin{bmatrix}
1 & 0 & 0 \\
0 & 1 & 0 \\
0 & 0 & 1
\end{bmatrix}
\iff
\frac{37}{9} = 1 \Rightarrow
\nexists c \in \mathbb{R} : Q \text{ ortogonale}
\]

\newpage
\noindent
\fbox{%
\begin{minipage}{\textwidth}
È dato il sistema lineare $Ax = b$, con $A = \begin{bmatrix} 2 & -1 & 0 \\ -1 & 2 & -1 \\ 0 & -1 & 2 \end{bmatrix}$, $b = \begin{bmatrix} 2 \\ -1 \\ 1 \end{bmatrix}$. Partendo dalla stima iniziale $x_0 = [0, 0, 0]^T$, calcolare la prima iterazione $x_1$ ottenuta applicando il metodo del gradiente al sistema assegnato.
\end{minipage}%
}

\vspace{1em}

Metodo del gradiente: $x_{k+1} = x_k + \alpha_k r_k$ con $r_k = b - Ax_k$, $\alpha_k = \frac{r_k^T r_k}{r_k^T A r_k},x_0 = \begin{bmatrix} 0 \\ 0 \\ 0 \end{bmatrix}$

\[
Ax_0 = 
\begin{bmatrix}
\colorbox{red!20}{$2$} & \colorbox{red!20}{$-1$} & \colorbox{red!20}{$0$} \\
\colorbox{blue!20}{$-1$} & \colorbox{blue!20}{$2$} & \colorbox{blue!20}{$-1$} \\
\colorbox{green!20}{$0$} & \colorbox{green!20}{$-1$} & \colorbox{green!20}{$2$}
\end{bmatrix}
\begin{bmatrix}
\colorbox{yellow!20}{$0$} \\
\colorbox{yellow!20}{$0$} \\
\colorbox{yellow!20}{$0$}
\end{bmatrix}
=
\begin{bmatrix}
\colorbox{red!20}{$2$} \cdot \colorbox{yellow!20}{$0$} + \colorbox{red!20}{$-1$} \cdot \colorbox{yellow!20}{$0$} + \colorbox{red!20}{$0$} \cdot \colorbox{yellow!20}{$0$} \\
\colorbox{blue!20}{$-1$} \cdot \colorbox{yellow!20}{$0$} + \colorbox{blue!20}{$2$} \cdot \colorbox{yellow!20}{$0$} + \colorbox{blue!20}{$-1$} \cdot \colorbox{yellow!20}{$0$} \\
\colorbox{green!20}{$0$} \cdot \colorbox{yellow!20}{$0$} + \colorbox{green!20}{$-1$} \cdot \colorbox{yellow!20}{$0$} + \colorbox{green!20}{$2$} \cdot \colorbox{yellow!20}{$0$}
\end{bmatrix}
=
\begin{bmatrix}
0 \\ 0 \\ 0
\end{bmatrix}
\]
\[
r_0 = b - Ax_0 = \begin{bmatrix} 2 \\ -1 \\ 1 \end{bmatrix} - \begin{bmatrix} 0 \\ 0 \\ 0 \end{bmatrix} = \begin{bmatrix} 2 \\ -1 \\ 1 \end{bmatrix}
\]

\[
r_0^T r_0 = \begin{bmatrix} 2 & -1 & 1 \end{bmatrix} \begin{bmatrix} 2 \\ -1 \\ 1 \end{bmatrix} = 4 + 1 + 1 = 6
\]

\[
Ar_0 = 
\begin{bmatrix}
\colorbox{red!20}{$2$} & \colorbox{red!20}{$-1$} & \colorbox{red!20}{$0$} \\
\colorbox{blue!20}{$-1$} & \colorbox{blue!20}{$2$} & \colorbox{blue!20}{$-1$} \\
\colorbox{green!20}{$0$} & \colorbox{green!20}{$-1$} & \colorbox{green!20}{$2$}
\end{bmatrix}
\begin{bmatrix}
\colorbox{violet!20}{$2$} \\
\colorbox{violet!20}{$-1$} \\
\colorbox{violet!20}{$1$}
\end{bmatrix}
=
\begin{bmatrix}
\colorbox{red!20}{$2$} \cdot \colorbox{violet!20}{$2$} + \colorbox{red!20}{$-1$} \cdot \colorbox{violet!20}{$-1$} + \colorbox{red!20}{$0$} \cdot \colorbox{violet!20}{$1$} \\
\colorbox{blue!20}{$-1$} \cdot \colorbox{violet!20}{$2$} + \colorbox{blue!20}{$2$} \cdot \colorbox{violet!20}{$-1$} + \colorbox{blue!20}{$-1$} \cdot \colorbox{violet!20}{$1$} \\
\colorbox{green!20}{$0$} \cdot \colorbox{violet!20}{$2$} + \colorbox{green!20}{$-1$} \cdot \colorbox{violet!20}{$-1$} + \colorbox{green!20}{$2$} \cdot \colorbox{violet!20}{$1$}
\end{bmatrix}
=
\begin{bmatrix}
5 \\ -5 \\ 3
\end{bmatrix}
\]

\[
r_0^T A r_0 = \begin{bmatrix} 2 & -1 & 1 \end{bmatrix} \begin{bmatrix} 5 \\ -5 \\ 3 \end{bmatrix} = 10 + 5 + 3 = 18
\]

\[
\alpha_0 = \frac{r_0^T r_0}{r_0^T A r_0} = \frac{6}{18} = \frac{1}{3}
\]

\[
x_1 = x_0 + \alpha_0 r_0 = \begin{bmatrix} 0 \\ 0 \\ 0 \end{bmatrix} + \frac{1}{3} \begin{bmatrix} 2 \\ -1 \\ 1 \end{bmatrix} = \begin{bmatrix} \frac{2}{3} \\ -\frac{1}{3} \\ \frac{1}{3} \end{bmatrix}
\]
\newpage
\noindent
\fbox{%
\begin{minipage}{\textwidth}

Dato $n \in \mathbb{N} \text{ scrivere funzioni per}$:

\begin{enumerate}
  \item Calcolare $A_n = I - 2 \dfrac{v v^T}{v^T v}$, con $v \in \mathbb{R}^n $ t.c. $v_i = i \quad \forall i \in \{1,\dots,n\}$
  \item Trovare $\lambda$, stima dell'autovalore di modulo massimo di $A_n$ con metodo delle potenze
  \item Verificare che $\lambda$ sia un autovalore di $A_n$
\end{enumerate}
\end{minipage}%
}

\vspace{1em}
\textbf{Funzione Mat(n):}
\begin{itemize}
  \item \texttt{v = [1,...,n]}
  \item \texttt{A = I - 2 $\cdot$ ($v \cdot v^t) / (v^t \cdot v)$} \Rightarrow   $a_{ij} = \begin{cases} \frac{1-2(v_i \cdot v_j)}{\sum_{i=1}^{n}v_i^2} i=j\\ 
   \frac{-2(v_i \cdot v_j)}{\sum_{i=1}^{n}v_i^2} i\not=j
  
  \end{cases}$
  \item \texttt{return A}
\end{itemize}

\vspace{0.5em}
\textbf{Funzione Potenze(A, tol):}
\begin{itemize}
  \item \texttt{x = vettore casuale normalizzato}
  \item $\lambda_0 = 0$
  \item \texttt{finchè $||Ax-\lambda x||_2$ < Tol:}
  \begin{itemize}
    \item \texttt{y = A} $\cdot$ \texttt{x}
    \item $\lambda = x^t \cdot y$
    \item \texttt{x =} $\dfrac{y}{\lVert y \rVert_2}$
  \end{itemize}
  \item \texttt{return} $\lambda$
\end{itemize}

\vspace{0.5em}
\textbf{Funzione Verifica(A, $\lambda$, x, tol):}
\begin{itemize}
  \item \texttt{se $||Ax-\lambda x||_2 < ||A||_2 \cdot$ tol:}
  \begin{itemize}
    \item \texttt{return OK}
  \end{itemize}
  \item \texttt{altrimenti:}
  \begin{itemize}
    \item \texttt{return ERRORE}
  \end{itemize}
\end{itemize}

\textbf{Pseudocodice}
\begin{itemize}
  \item \texttt{n = input()}
  \item \texttt{A = Mat(n)}
  \item \texttt{tol = 10\textsuperscript{-6}}
  \item \texttt{$\lambda$ = Potenze(A, tol)}
  \item \texttt{Verifica(A, $\lambda$, x, tol)}
\end{itemize}

\newpage
\noindent

\fbox{%
\begin{minipage}{\textwidth}
Sia \( I = \int_{-1}^{3} f(x)\,dx \) con: $f(x) = \begin{cases}
    x^2 & \text{se } x \leq 1 \\
    (x - 2)^2 & \text{se } x > 1
  \end{cases}$


Suggerire una divisione del dominio di integrazione e indicare una formula di quadratura che permettano di calcolare \( I \) esattamente. Qual \`e il numero minimo di valutazioni funzionali necessarie per calcolare \( I \) esattamente?

Calcolare l'approssimazione a \( I \) ottenuta utilizzando la formula dei trapezi composta su nodi equispaziati con spaziatura \( h = 1 \).

Calcolare l'errore commesso utilizzando il risultato del punto precedente rispetto alla soluzione esatta (ottenuta risolvendo l'integrale analiticamente).

Scrivere la stima dell'errore relativa al metodo dei trapezi separatamente per i due sottointervalli \([-1,1]\) e \([1,3]\).

Perch\'e per utilizzare la stima di errore \`e necessario dividere l'intervallo di integrazione \([-1,3]\) in due parti?
\end{minipage}%
}

\vspace{1em}

\noindent

Divido in \([-1,1]\), \([1,3]\). In entrambi \(f\) \`e di grado 2. Uso Simpson (esatta per grado \(\leq 2\)), richiede 3 nodi per intervallo \(\Rightarrow 5\) valutazioni (nodo condiviso).

\[
I = \int_{-1}^{1} x^2\,dx + \int_{1}^{3} (x-2)^2\,dx = \left[ \frac{x^3}{3} \right]_{-1}^{1} + \left[ \frac{(x-2)^3}{3} \right]_{1}^{3} = \frac{2}{3} + \frac{2}{3} = \frac{4}{3}
\]

Trapezi composta con \(h = 1\): nodi \(x \in \{ -1, 0, 1, 2, 3\}\)
\[
T = \sum_{i=1}^nh\cdot\frac{f(x_i)-f(x_{i+1})}{2}=\frac{1}{2} [f(-1) + 2f(0) + 2f(1) + 2f(2) + f(3)] = \frac{1}{2} [1 + 0 + 2 + 0 + 1] = \frac{1}{2} \cdot 4 = 2
\]

Errore assoluto: \(\left|2 - \frac{4}{3}\right| = \frac{2}{3}\)

Errore trapezi su \([-1,1]\):
\[
E_1 = -\frac{(1 - (-1))^3}{12} f''(\xi_1) = -\frac{8}{12} \cdot 2 = -\frac{4}{3}
\]

Errore su \([1,3]\):
\[
E_2 = -\frac{(3 - 1)^3}{12} f''(\xi_2) = -\frac{8}{12} \cdot 2 = -\frac{4}{3}
\]

Totale stimato: \(E = E_1 + E_2 = -\frac{8}{3}\)

Dividere \([-1,3]\) \`e necessario perch\'e \(f(x)\) \`e definita a tratti: non \`e polinomio unico su tutto l'intervallo.

\newpage
\fbox{
\parbox{0.97\textwidth}{
Dato un sistema lineare $Ax = b$, si chiede di:

1. Scrivere una funzione che risolva un sistema lineare $Ux = b$ dove $U$ è triangolare superiore.\\
2. Scrivere una funzione che risolva un sistema lineare $Lx = b$ dove $L$ è triangolare inferiore.\\
3. Scrivere una funzione che, dati in input una matrice $A$, un vettore $b$ e una tolleranza $Tol$, restituisca un vettore $x$ approssimato tale che $\|x^{(k+1)} - x^{(k)}\| < Tol$, usando il metodo di Gauss-Seidel simmetrico.

Il metodo di Gauss-Seidel simmetrico è un metodo iterativo composto da due fasi:

Fase 1: si scrive $A = E + F$ dove
\[
E_{ij} = \begin{cases}
A_{ij} & \text{se } i \ge j \\
0 & \text{altrimenti}
\end{cases}, \quad
F_{ij} = \begin{cases}
A_{ij} & \text{se } i < j \\
0 & \text{altrimenti}
\end{cases}
\]
e si risolve il sistema $E x^{(k+1/2)} = b - F x^{(k)}$ per $x^{(k+1/2)}$.

Fase 2: si scrive $A = E + F$ con nuova definizione:
\[
E_{ij} = \begin{cases}
A_{ij} & \text{se } i > j \\
0 & \text{altrimenti}
\end{cases}, \quad
F_{ij} = \begin{cases}
A_{ij} & \text{se } i \le j \\
0 & \text{altrimenti}
\end{cases}
\]
e si risolve il sistema $F x^{(k+1)} = b - E x^{(k+1/2)}$ per $x^{(k+1)}$.
}
}
\begin{flushleft}
\textbf{Risolutore per sistemi triangolari inferiori (forward substitution):}
\begin{tabbing}
\hspace{1.5cm} \= \hspace{10cm} \= \kill
\text{Input:} matrice triangolare inferiore \( L \), vettore \( b \in \mathbb{R}^n \) \\
\text{for} \( i = 1 \) \text{to} \( n \): \\
\> \( x_i = \left(b_i - \sum_{j=1}^{i-1} L_{i,j} x_j\right) / L_{i,i} \) \\
\text{Output:} vettore \( x \in \mathbb{R}^n \) soluzione di \( Lx = b \) 
\end{tabbing}
\end{flushleft}
\begin{flushleft}
\textbf{Risolutore per sistemi triangolari superiori (back substitution):}
\begin{tabbing}
\hspace{1.5cm} \= \hspace{10cm} \= \kill
\text{Input:} matrice triangolare superiore \( U \), vettore \( b \in \mathbb{R}^n \) \\
\text{for} \( i = n \) \text{to} \( 1 \) \text{step} \( -1 \): \\
\> \( x_i = \left(b_i - \sum_{j=i+1}^{n} U_{i,j} x_j\right) / U_{i,i} \) \\
\text{Output:} vettore soluzione \( x \in \mathbb{R}^n \) di \( Ux = b \) 
\end{tabbing}
\end{flushleft}
\begin{flushleft}
\textbf{Metodo di Gauss-Seidel simmetrico:}
\begin{tabbing}
\hspace{1.5cm} \= \hspace{10cm} \= \kill
\text{Input:} A, b, Tol, \( x^{(0)} \) \\
Decomponi \( A = E_1 + F_1 \), con \( E_1 = \{A_{ij} \mid i \ge j\} \), \( F_1 = \{A_{ij} \mid i < j\} \) \\
Decomponi \( A = E_2 + F_2 \), con \( E_2 = \{A_{ij} \mid i > j\} \), \( F_2 = \{A_{ij} \mid i \le j\} \) \\
\( k = 0 \) \\
\text{while} \( \|x^{(k+1)} - x^{(k)}\| \ge \text{Tol} \): \\
\> \( y = b - F_1 x^{(k)} \) \\
\> risolvi \( E_1 \hat{x}^{(k+1)} = y \) (risolutore inferiore) \\
\> \( z = b - E_2 \hat{x}^{(k+1)} \) \\
\> risolvi \( F_2 x^{(k+1)} = z \) (risolutore superiore) \\
\> \( k = k + 1 \) \\
\textbf{Output:} \( x^{(k+1)} \), numero iterazioni \( = k \)
\end{tabbing}
\end{flushleft}

\newpage
\thispagestyle{empty}

\fbox{%
\begin{minipage}{\textwidth}
Considera l'equazione differenziale di secondo ordine
\[
y''(t) + \omega^2 y(t) = 0, \quad y(0) = a, \quad y'(0) = b \quad \text{ con }\omega > 0
\]

\begin{enumerate}
  \item Riscrivi il problema come sistema del primo ordine \( y' = f(y) \), separando la parte lineare e quella non lineare se presente.
  
  \item Applica il metodo di Eulero esplicito:
  \begin{enumerate}
    \item Scrivi lo schema numerico per il sistema.
    \item Calcola un passo con \(\Delta t = 0.1\), \(\omega=2\), \(a=1\), \(b=0\).
  \end{enumerate}
  
  \item Studia la stabilità:
  \begin{enumerate}
    \item Scrivi l’equazione test associata.
    \item Definisci la regione di stabilità assoluta del metodo di Eulero esplicito.
    \item Determina se il metodo è stabile per il problema dato.
  \end{enumerate}
  
  \item Parla di convergenza e errore di troncamento:
  \begin{enumerate}
    \item Definisci l’errore di troncamento locale e il suo ordine per Eulero esplicito.
    \item Definisci la convergenza e la relazione tra consistenza e stabilità.
  \end{enumerate}
  
  \item Metodo di Eulero implicito:
  \begin{enumerate}
    \item Scrivi lo schema implicito per il sistema.
    \item Forma il sistema non lineare \( G(y_{i+1}) = 0 \).
    \item Calcola lo Jacobiano di \( G \).
    \item Scrivi la formula di Newton per risolvere il sistema a ogni passo.
  \end{enumerate}
  
  \item Metodo di Runge-Kutta esplicito di ordine 2:
  \begin{enumerate}
    \item Scrivi lo schema generale.
    \item Applica il primo passo con i dati iniziali.
    \item Verifica l’ordine tramite lo sviluppo di Taylor.
  \end{enumerate}
  
  \item Passo adattivo:
  \begin{enumerate}
    \item Spiega come scegliere il passo in modo adattivo.
    \item Dai un criterio basato sull’errore locale stimato (metodi embedded).
  \end{enumerate}
  
  \item Sistema di Lotka-Volterra:
  \[
  \begin{cases}
  y_1' = \alpha y_1 - \beta y_1 y_2 \\
  y_2' = -\gamma y_2 + \delta y_1 y_2
  \end{cases}
  \]
  \begin{enumerate}
    \item Scrivi \( y' = f(y) \).
    \item Implementa Eulero esplicito con \(\Delta t=0.1\), \(T=30\).
    \item Traccia \( y_1 \) vs \( y_2 \).
    \item Spiega come stimare errore numerico se soluzione esatta non nota.
    \item Calcola lo Jacobiano del sistema.
  \end{enumerate}
\end{enumerate}
\end{minipage}%
}

\vspace{1em}
\noindent

\begin{enumerate}
\item \textbf{Riscrittura come sistema:}

Ponendo \( y_1 = y \), \( y_2 = y' \) si ha

\[
\begin{cases}
y_1' = y_2 \\
y_2' = -\omega^2 y_1
\end{cases}
\quad \Rightarrow \quad
\frac{d}{dt} \begin{bmatrix} y_1 \\ y_2 \end{bmatrix} 
= \underbrace{\begin{bmatrix} 0 & 1 \\ -\omega^2 & 0 \end{bmatrix}}_{A}
\begin{bmatrix} y_1 \\ y_2 \end{bmatrix}.
\]

Qui la parte lineare è tutta contenuta in \( A y \) e non ci sono termini non lineari.

\item \textbf{Metodo di Eulero esplicito:}

(a) Schema

\[
y_{i+1} = y_i + \Delta t \, A y_i.
\]

Esplicitamente

\[
\begin{bmatrix} y_1^{i+1} \\ y_2^{i+1} \end{bmatrix}
= \begin{bmatrix} y_1^i \\ y_2^i \end{bmatrix} + \Delta t
\begin{bmatrix} 0 & 1 \\ -\omega^2 & 0 \end{bmatrix}
\begin{bmatrix} y_1^i \\ y_2^i \end{bmatrix}.
\]

(b) Calcolo con \(\Delta t=0.1\), \(\omega=2\), \(y_1^0=1\), \(y_2^0=0\):

\[
y^1 = \begin{bmatrix} 1 \\ 0 \end{bmatrix} + 0.1
\begin{bmatrix} 0 & 1 \\ -4 & 0 \end{bmatrix}
\begin{bmatrix} 1 \\ 0 \end{bmatrix}
= \begin{bmatrix} 1 \\ 0 \end{bmatrix} + 0.1 \begin{bmatrix} 0 \\ -4 \end{bmatrix} 
= \begin{bmatrix} 1 \\ -0.4 \end{bmatrix}.
\]

\item \textbf{Stabilità:}

Consideriamo l’equazione test lineare

\[
y' = \lambda y \quad \text{con } \lambda \in \mathbb{C}.
\]

La soluzione esatta è

\[
y(t) = y(0) e^{\lambda t} \Rightarrow |y(t)| = |y(0)| e^{\Re(\lambda) t}
\]

$\begin{cases}
\Re(\lambda) < 0 \Rightarrow \text{soluzione decrescente (problema stabile)}
\\
\Re(\lambda) > 0 \Rightarrow \text{soluzione divergente (problema instabile)}\end{cases}$

Applichiamo ora il metodo di Eulero esplicito all’equazione test:

\[
y_{n+1} = y_n + \Delta t \lambda y_n = (1 + \lambda \Delta t) y_n
\Rightarrow
y_n = (1 + \lambda \Delta t)^n y_0.
\]

Affinché il metodo sia stabile, imponiamo

\[
|y_n| \leq |y_0| \Rightarrow |1 + \lambda \Delta t| \leq 1.
\]

La regione di stabilità del metodo è quindi:

\[
R = \left\{ z \in \mathbb{C} : |1 + z| < 1 \right\}, \quad \text{con } z = \lambda \Delta t.
\]

Nel nostro problema il sistema è:

\[
\frac{d}{dt} \begin{bmatrix} y_1 \\ y_2 \end{bmatrix}
= A \begin{bmatrix} y_1 \\ y_2 \end{bmatrix}, \quad
A = \begin{bmatrix} 0 & 1 \\ -\omega^2 & 0 \end{bmatrix}
\Rightarrow \text{autovalori: } \lambda = \pm i\omega.
\]

Con \( \omega = 2 \) e \( \Delta t = 0.1 \):

\[
z = \lambda \Delta t = \pm i \cdot 0.2 \Rightarrow |1 + z| = \sqrt{1^2 + 0.2^2} = \sqrt{1.04} > 1.
\]

Quindi:

\[
z \notin R \Rightarrow \text{il metodo di Eulero esplicito è instabile per questo problema.}
\]


\item \textbf{Convergenza ed errore di troncamento:}

(a) Errore di troncamento locale (ETL):

\[
\tau(\Delta t) = \frac{1}{\Delta t} \left| y(t_{i+1}) - y(t_i) - \Delta t f(y(t_i)) \right| = O(\Delta t).
\]

Per Eulero esplicito è di ordine 1.

(b) Convergenza: il metodo converge se

\[
\lim_{\Delta t \to 0} \max_i | y_i - y(t_i) | = 0.
\]

Teorema fondamentale (Lax equivalence): 

\[
\text{Consistenza} + \text{Stabilità} \implies \text{Convergenza}.
\]



\item \textbf{Metodo di Eulero implicito:}

(a) Schema

\[
y_{i+1} = y_i + \Delta t \, A y_{i+1}.
\]

(b) Sistema non lineare da risolvere

\[
G(y_{i+1}) := y_{i+1} - y_i - \Delta t A y_{i+1} = 0.
\]

(c) Jacobiano

\[
J = \frac{\partial G}{\partial y_{i+1}} = I - \Delta t A.
\]

(d) Formula di Newton

Dato \( y_{i+1}^{(k)} \), aggiorno con

\[
y_{i+1}^{(k+1)} = y_{i+1}^{(k)} - J^{-1} G(y_{i+1}^{(k)}).
\]



\item \textbf{Metodo di Runge-Kutta esplicito ordine 2:}

(a) Schema generico

\[
\begin{cases}
k_1 = f(t_i, y_i) \\
k_2 = f(t_i + \alpha \Delta t, y_i + \beta \Delta t k_1) \\
y_{i+1} = y_i + \Delta t (c_1 k_1 + c_2 k_2),
\end{cases}
\]

con coefficenti scelti per ordine 2, ad esempio Heun:

\[
\alpha = 1, \quad \beta = 1, \quad c_1 = \frac{1}{2}, \quad c_2 = \frac{1}{2}.
\]

(b) Primo passo con dati iniziali:

\[
k_1 = A y_0, \quad k_2 = A (y_0 + \Delta t k_1).
\]

(c) Verifica ordine tramite sviluppo di Taylor è standard e conferma ordine 2.



\item \textbf{Passo adattivo:}

(a) Si sceglie \(\Delta t\) variabile per mantenere l’errore locale entro una tolleranza \(\varepsilon\).

(b) Metodo embedded (es. RK45) fornisce due stime \( y_{i+1}^{(p)} \), \( y_{i+1}^{(p-1)} \). La differenza

\[
e_i = \| y_{i+1}^{(p)} - y_{i+1}^{(p-1)} \|
\]

stima l’errore locale e permette di adattare \(\Delta t\):

\[
\Delta t_{\text{new}} = \Delta t \left( \frac{\varepsilon}{e_i} \right)^{1/(p+1)}.
\]


\item \textbf{Sistema Lotka-Volterra:}

(a) Scrittura vettoriale

\[
y = \begin{bmatrix} y_1 \\ y_2 \end{bmatrix}, \quad
f(y) = \begin{bmatrix} \alpha y_1 - \beta y_1 y_2 \\ -\gamma y_2 + \delta y_1 y_2 \end{bmatrix}.
\]

(b) Applica Eulero esplicito

\[
y_{i+1} = y_i + \Delta t f(y_i).
\]

(c) Grafico fase \( (y_1, y_2) \).

(d) Se soluzione esatta non nota, si può stimare errore confrontando soluzioni con passi \( \Delta t \) e \( \Delta t/2 \) (metodo della doppia risoluzione).

(e) Jacobiano di \( f \):

\[
J_f(y) = 
\begin{bmatrix}
\alpha - \beta y_2 & -\beta y_1 \\
\delta y_2 & -\gamma + \delta y_1
\end{bmatrix}.
\]

\end{enumerate}

\fbox{
\parbox{0.95\textwidth}{È data la matrice:

$A = \frac{vv^\top}{\|v\|^2} \qquad \text{con } v \in \mathbb{R}^n, \ v \neq 0$


\newline

Scegliere tutte le affermazioni che sono vere tra:

\begin{itemize}
    \item[\boxed{\text{A}}] A è invertibile
    \item[\boxed{\text{B}}] \(\forall u \in \mathbb{R}^n, u \ne 0 \Rightarrow Au\) è un vettore parallelo a \(v\)
    \item[\boxed{\text{C}}] \(\forall u \in \mathbb{R}^n, u \ne 0 \Rightarrow Au\) è un vettore parallelo a \(u\)
    \item[\boxed{\text{D}}] A è simmetrica
\end{itemize}}}

\begin{itemize}
    \item[\boxed{\text{A}}] \textbf{Falsa} \\
    La matrice \( A = \frac{vv^\top}{\|v\|^2} \) è un prodotto esterno diviso per uno scalare:
    \[
    \operatorname{rank}(A) = \operatorname{rank}(vv^\top) = 1 \quad (\text{poiché } v \ne 0)
    \]
    Una matrice \( n \times n \) è invertibile solo se ha rango pieno (\( \operatorname{rank} = n \))\\
    \(\boxed{\text{A non è invertibile 
    $\forall n \in \mathbb{N}\backslash$\{1\}}}\)

    \item[\boxed{\text{B}}] \textbf{Vera} \\
    Calcolo:
    \[
    Au = \frac{vv^\top}{\|v\|^2} u = \left( \frac{v^\top u}{\|v\|^2} \right) v
    \]
    Quindi:
    \[
    \boxed{Au = \lambda v \quad \text{con } \lambda = \frac{v^\top u}{\|v\|^2}}
    \Rightarrow Au \parallel v
    \]

    \item[\boxed{\text{C}}] \textbf{Falsa} \\
    Abbiamo appena visto che:
    \[
    Au = \lambda v \Rightarrow Au \in \operatorname{span}(v)
    \]
    Quindi in generale:
    \[
    Au \not\in \operatorname{span}(u) \quad \text{(tranne se } u \parallel v)
    \]
    \(\boxed{Au \parallel v, \text{ non } u}\)

    \item[\boxed{\text{D}}] \textbf{Vera} \\
    La trasposizione dà:
    \[
    A^\top = \left( \frac{vv^\top}{\|v\|^2} \right)^\top = \frac{(vv^\top)^\top}{\|v\|^2} = \frac{vv^\top}{\|v\|^2} = A
    \]
    \(\boxed{A \text{ è simmetrica}}\)
\end{itemize}

\fbox{
\parbox{0.95\textwidth}{
Utilizzare lo spazio sottostante per dimostrare che la norma 1 di un vettore è equivalente alla sua norma infinito, per ogni vettore \(u \in \mathbb{R}^n\). Più precisamente, dimostrare che:

\[
\frac{1}{n} \|u\|_1 \leq \|u\|_\infty \leq \|u\|_1, \quad \text{dove } \|u\|_1 = \sum_{i=1}^n |u_i|, \quad \|u\|_\infty = \max_i |u_i|
\]
}
}

\begin{itemize}
    \boxed{\text{Disuguaglianza sinistra}}
    Poiché \( \|u\|_\infty = \max_i |u_i| \), allora:
    \[
    |u_i| \leq \|u\|_\infty \quad \forall i \Rightarrow \sum_{i=1}^n |u_i| \leq \sum_{i=1}^n \|u\|_\infty = n\|u\|_\infty
    \Rightarrow \|u\|_1 \leq n\|u\|_\infty
    \Rightarrow \|u\|_\infty \geq \frac{1}{n}\|u\|_1
    \]

    \boxed{\text{Disuguaglianza destra}}
    Sia \( j \in \{1,\dots,n\} \) tale che \( |u_j| = \|u\|_\infty \), allora:
    \[
    \|u\|_1 = \sum_{i=1}^n |u_i| \geq |u_j| = \|u\|_\infty
    \Rightarrow \|u\|_\infty \leq \|u\|_1
    \]
\end{itemize}

\fbox{
\parbox{0.95\textwidth}{
Supponiamo di conoscere la SVD di una matrice \( A = U\Sigma V^\top \), e di volerla utilizzare per risolvere il sistema lineare \( Ax = b \). Quale fra i seguenti è il modo corretto di risolvere il sistema?

\begin{itemize}
    \item[\boxed{\text{A}}] Risolvo prima \( Uy = b \), poi \( \Sigma z = y \) e infine \( x = Vz \)
    \item[\boxed{\text{B}}] Risolvo prima \( Uy = b \), poi \( \Sigma y = z \) e infine \( x = Vb \)
    \item[\boxed{\text{C}}] Risolvo prima \( y = V^\top x \), poi \( z = \Sigma y \) e infine \( Uy = b \)
    \item[\boxed{\text{D}}] Risolvo prima \( Uy = b \), poi \( \Sigma z = y \) e infine \( V^\top x = z \)
\end{itemize}
}
}

\begin{itemize}
    \boxed{\text{Risposta corretta: A}} \\
    Se \( A = U\Sigma V^\top \), allora:

    \[
    Ax = b \Rightarrow U\underbrace{\Sigma V^\top x}_y = b
    \Rightarrow \Sigma\underbrace{ V^\top x}_z = y \Rightarrow  V^\top x=z\]

    Quindi siccome V è ortogonale ($V=V^T$) risolviamo
    \[
    Uy = b \Rightarrow \Sigma z = y \Rightarrow x = Vz
    \]
\end{itemize}

\newpage
\fbox{
\parbox{0.95\textwidth}{
Ogni metodo di Runge-Kutta del secondo ordine a 2 stadi può essere scritto sotto forma di tableau di Butcher nel modo seguente:

\[
\begin{array}{c|cc}
0 & 0 & 0 \\
\alpha & \alpha & 0 \\
\hline
 & 1 - \frac{1}{2\alpha} & \frac{1}{2\alpha} \\
\end{array}
\]

dove \(\alpha\) è un numero reale. Verificare che la regione di stabilità assoluta del metodo non dipende da \(\alpha\), e dunque tutti i metodi Runge-Kutta a due stadi del secondo ordine hanno la stessa regione di stabilità. Utilizzare il box sottostante per riportare i calcoli effettuati per raggiungere questa conclusione.
}
}

\vspace{1em}


\textbf{Obiettivo:} calcolare la \textit{funzione di stabilità assoluta} \( R(z) \), dove \( z = \lambda \Delta t \), per il metodo Runge-Kutta dato.

\textbf{Equazione test:} \( y' = \lambda y \Rightarrow f(t, y) = \lambda y \)

\textbf{Notazione:}
\[
y_{n+1} = y_n + \Delta t \left( b_1 k_1 + b_2 k_2 \right)
\]
\[
\begin{cases}
k_1 = f(t_n, y_n) = \lambda y_n \\
k_2 = f(t_n + \alpha \Delta t, y_n + \alpha \Delta t \cdot k_1)
= \lambda \left( y_n + \alpha \Delta t \cdot \lambda y_n \right)
= \lambda y_n (1 + \alpha \Delta t \lambda)
\end{cases}
\]

\textbf{Coefficienti:} \( b_1 = 1 - \frac{1}{2\alpha}, \quad b_2 = \frac{1}{2\alpha} \)

\textbf{Sviluppo:}
\[
y_{n+1} = y_n + \Delta t \left[
\left( 1 - \frac{1}{2\alpha} \right) \cdot \lambda y_n
+ \frac{1}{2\alpha} \cdot \lambda y_n (1 + \alpha \Delta t \lambda)
\right]
\]

\textbf{Raccolgo:}
\[
= y_n + \Delta t \cdot \lambda y_n \left[
1 - \frac{1}{2\alpha} + \frac{1}{2\alpha}(1 + \alpha \Delta t \lambda)
\right]
\]

\textbf{Espando e semplifico:}
\[
= y_n + \Delta t \cdot \lambda y_n \left[
1 - \frac{1}{2\alpha} + \frac{1}{2\alpha} + \frac{\Delta t \lambda}{2}
\right]
= y_n \left( 1 + \lambda \Delta t + \frac{1}{2} (\lambda \Delta t)^2 \right)
\]

\textbf{Funzione di stabilità:}
\[
\boxed{R(z) = 1 + z + \frac{1}{2} z^2 \quad \text{con } z = \lambda \Delta t}
\]

\textbf{Conclusione:}
La funzione \( R(z) \) non dipende da \( \alpha \), quindi:

\[
\boxed{
\text{Tutti i metodi RK a 2 stadi di ordine 2 hanno la stessa regione di stabilità}
}
\]

\textbf{Regione di stabilità assoluta:}
\[
\boxed{
\left\{ z \in \mathbb{C} \ \middle| \ \left| 1 + z + \frac{1}{2}z^2 \right| < 1 \right\}
}
\]


\paragraph*{Riepilogo sulle ODE}

Un’equazione differenziale ordinaria (ODE) è un’equazione nella forma:

\[
\frac{dy}{dt} = f(t, y)
\]

dove:
\begin{itemize}
    \item \( y(t) \) è la funzione incognita (da trovare),
    \item \( f(t, y) \) è una funzione nota che specifica il comportamento del sistema.
\end{itemize}

Esempio semplice: \( \frac{dy}{dt} = \lambda y \), detta \textbf{equazione test lineare}. La soluzione esatta è:

\[
y(t) = y(0) e^{\lambda t}
\]

\paragraph{Obiettivo dei metodi numerici}

In pratica, la soluzione esatta \( y(t) \) è spesso difficile o impossibile da calcolare. Si vuole quindi calcolare approssimazioni \( y_0, y_1, y_2, \dots \) in punti \( t_0, t_1, t_2, \dots \) dove:

\[
t_n = t_0 + n \Delta t
\]

\paragraph{Metodo di Eulero (Runge-Kutta di ordine 1)}

Il metodo più semplice è:

\[
y_{n+1} = y_n + \Delta t \cdot f(t_n, y_n)
\]

Questo è il metodo di Eulero, anche chiamato Runge-Kutta di ordine 1.

\paragraph{Metodi di Runge-Kutta (RK)}

I metodi Runge-Kutta sono una famiglia di metodi con la seguente forma generale:

\[
\begin{cases}
k_1 = f(t_n, y_n) \\
k_2 = f(t_n + c_2 \Delta t, y_n + a_{21} \Delta t \cdot k_1) \\
\vdots \\
k_s = f(t_n + c_s \Delta t, y_n + \sum_{j=1}^{s-1} a_{sj} \Delta t \cdot k_j) \\
\\
y_{n+1} = y_n + \Delta t \cdot \sum_{i=1}^s b_i k_i
\end{cases}
\]

\paragraph{Tableau di Butcher}

Tutti i coefficienti si scrivono in una tabella detta \textbf{tableau di Butcher}:

\[
\begin{array}{c|ccc}
c_1 & a_{11} & \cdots & a_{1s} \\
\vdots & \vdots & & \vdots \\
c_s & a_{s1} & \cdots & a_{ss} \\
\hline
& b_1 & \cdots & b_s
\end{array}
\]

\paragraph{Metodo Runge-Kutta a 2 stadi di ordine 2}

Un metodo Runge-Kutta del secondo ordine con 2 stadi ha il seguente tableau:

\[
\begin{array}{c|cc}
0 & 0 & 0 \\
\alpha & \alpha & 0 \\
\hline
& 1 - \frac{1}{2\alpha} & \frac{1}{2\alpha}
\end{array}
\]

\noindent
dove \( \alpha \in \mathbb{R} \) è un parametro.

\paragraph{Equazione test: \( y' = \lambda y \)}

Per studiare la stabilità, si usa l’\textbf{equazione test}:

\[
\frac{dy}{dt} = \lambda y
\quad \Rightarrow \quad f(t, y) = \lambda y
\]

\paragraph{Cos'è la funzione di stabilità \( R(z) \)}

Applichiamo un metodo RK all’equazione test:

\[
y_{n+1} = R(z) \cdot y_n
\quad \text{dove} \quad z := \lambda \Delta t
\]

\( R(z) \) è una funzione razionale che descrive come il metodo propaga \( y_n \) nel tempo. Serve a studiare la \textbf{stabilità} del metodo.

\textbf{Stabilità assoluta} significa che, se \( |R(z)| < 1 \), allora la soluzione numerica non esplode nel tempo (si smorza)

\paragraph{Calcolo di \( R(z) \) per il metodo dato}

Dato il tableau con parametro \( \alpha \), i coefficienti sono:

\[
b_1 = 1 - \frac{1}{2\alpha}, \quad b_2 = \frac{1}{2\alpha}
\]

\[
k_1 = \lambda y_n
\quad,\quad
k_2 = \lambda \left( y_n + \alpha \Delta t \cdot \lambda y_n \right)
= \lambda y_n (1 + \alpha z)
\]

\[
y_{n+1} = y_n + \Delta t \cdot \left[
b_1 k_1 + b_2 k_2
\right]
= y_n + z \cdot \left[
b_1 + b_2(1 + \alpha z)
\right] \cdot y_n
\]

Sostituiamo i coefficienti:

\[
R(z) = 1 + z \left( 1 - \frac{1}{2\alpha} + \frac{1}{2\alpha} (1 + \alpha z) \right)
= 1 + z + \frac{1}{2} z^2
\]

\paragraph{Conclusione}

\[
\boxed{
R(z) = 1 + z + \frac{1}{2}z^2
\quad \text{non dipende da } \alpha
}
\]

\[
\Rightarrow \text{Tutti i metodi RK di ordine 2 a 2 stadi hanno la stessa funzione di stabilità}
\]

\[
\text{Regione di stabilità } = \left\{ z \in \mathbb{C} \mid |R(z)| < 1 \right\}
= \left\{ z \in \mathbb{C} \mid \left|1 + z + \frac{1}{2}z^2 \right| < 1 \right\}
\]


\newpage
\fbox{%
\parbox{\textwidth}{%
\textbf{Esercizio:} \\
Calcolare il polinomio che interpola i punti \((0,2), (1,2), (2,4)\). Verificare che si ottiene lo stesso risultato sia usando la matrice di Vandermonde che usando i polinomi di Lagrange.
}%
}

\vspace{0.5em}

\textbf{Soluzione:}

Abbiamo i punti:  
\[
(x_0, y_0) = (0,2), \quad (x_1, y_1) = (1,2), \quad (x_2, y_2) = (2,4)
\]

\bigskip

\textbf{1. Metodo della matrice di Vandermonde}

Cerchiamo un polinomio:
\[
p(x) = a_0 + a_1 x + a_2 x^2
\]

che soddisfa:
\[
\begin{cases}
p(0) = a_0 = 2 \\
p(1) = a_0 + a_1 + a_2 = 2 \\
p(2) = a_0 + 2a_1 + 4a_2 = 4
\end{cases}
\]

In forma matriciale:
\[
\begin{bmatrix}
1 & 0 & 0 \\
1 & 1 & 1 \\
1 & 2 & 4 \\
\end{bmatrix}
\begin{bmatrix}
a_0 \\ a_1 \\ a_2
\end{bmatrix}
=
\begin{bmatrix}
2 \\ 2 \\ 4
\end{bmatrix}
\]

Risolvendo:
\[
a_0 = 2
\]
\[
a_1 + a_2 = 0 \implies a_1 = -a_2
\]
\[
2a_1 + 4a_2 = 2 \implies 2(-a_2) + 4a_2 = 2 \implies 2a_2 = 2 \implies a_2 = 1 \implies a_1 = -1
\]

Quindi:
\[
p(x) = 2 - x + x^2
\]

\bigskip

\textbf{2. Metodo dei polinomi di Lagrange}

Definiamo:
\[
p(x) = y_0 L_0(x) + y_1 L_1(x) + y_2 L_2(x)
\]
con
\[
L_0(x) = \frac{(x - x_1)(x - x_2)}{(x_0 - x_1)(x_0 - x_2)} = \frac{(x - 1)(x - 2)}{(0-1)(0-2)} = \frac{(x-1)(x-2)}{2}
\]
\[
L_1(x) = \frac{(x - x_0)(x - x_2)}{(x_1 - x_0)(x_1 - x_2)} = \frac{x(x - 2)}{1 \cdot (-1)} = -x(x-2)
\]
\[
L_2(x) = \frac{(x - x_0)(x - x_1)}{(x_2 - x_0)(x_2 - x_1)} = \frac{x(x-1)}{2 \cdot 1} = \frac{x(x-1)}{2}
\]

Quindi:
\[
p(x) = 2 \cdot \frac{(x-1)(x-2)}{2} + 2 \cdot [-x(x-2)] + 4 \cdot \frac{x(x-1)}{2}
\]

Calcoliamo ogni termine:
\[
2 \cdot \frac{(x-1)(x-2)}{2} = (x-1)(x-2) = x^2 - 3x + 2
\]
\[
2 \cdot [-x(x-2)] = -2x^2 + 4x
\]
\[
4 \cdot \frac{x(x-1)}{2} = 2x^2 - 2x
\]

Sommiamo:
\[
(x^2 - 3x + 2) + (-2x^2 + 4x) + (2x^2 - 2x) = x^2 - x + 2
\]

\[
\boxed{
\text{Entrambi i metodi danno: } p(x) = x^2 - x + 2
}
\]

\vspace{2em}

\fbox{%
\parbox{\textwidth}{%
\textbf{Esercizio:} \\
Data una matrice \( A \in \mathbb{R}^{n \times n} \) di numeri reali, scrivere una funzione che riceve \( A \) e ritorna due array \( E \in \mathbb{R}^k \) e \( \text{Ind} \in \mathbb{R}^{k \times 2} \) contenenti solo gli elementi di \( A \) con modulo maggiore di una soglia \( \text{tol} \), e i corrispondenti indici. La relazione è che per ogni elemento significativo \( a_{ij} \), esiste un indice \( k \) tale che \( E_k = a_{ij} \) e \( \text{Ind}_k = (i,j) \). Inoltre, bisogna ritornare \( n \) per poter ricostruire la matrice originale da questi dati.
}%
}

\vspace{1em}
\textbf{Pseudocodice:}
\vspace{1em}

\begin{cases}
k := 0 \\
E := \emptyset \\
\text{Ind} := \emptyset \\
\text{per } i = 1, \dots, n \\
\quad \text{per } j = 1, \dots, n \\
\quad\quad \text{se } |a_{ij}| > \text{tol} \\
\quad\quad\quad k := k + 1 \\
\quad\quad\quad E_k := a_{ij} \\
\quad\quad\quad \text{Ind}_k := (i, j) \\
\text{ritorna } E, \text{Ind}, n
\end{cases}

\vspace{1em}
\textbf{Relazione formale:}

\[
|a_{ij}| > \text{tol} \implies \exists k \in \{1,\dots,k\} : E_k = a_{ij}, \quad \text{Ind}_k = (i,j)
\]

\[
A_{ij} = 
\begin{cases}
E_k & \text{se } \exists k: \text{Ind}_k = (i,j) \\
0 & \text{altrimenti}
\end{cases}
\]
\newpage
\textbf{Codice MATLAB:}

\begin{verbatim}
function [E, Ind, n] = sparsify(A, tol)
    n = size(A, 1);           % dimensione matrice quadrata
    k = 0;                    % contatore elementi significativi
    E = [];
    Ind = [];
    for i = 1:n
        for j = 1:n
            if abs(A(i,j)) > tol
                k = k + 1;
                E(k) = A(i,j);
                Ind(k,:) = [i, j];
            end
        end
    end
end
\end{verbatim}

\bigskip

\textbf{Funzione per ricostruire la matrice originale:}

\begin{verbatim}
function A = reconstruct(E, Ind, n)
    A = zeros(n);
    for k = 1:length(E)
        i = Ind(k,1);
        j = Ind(k,2);
        A(i,j) = E(k);
    end
end
\end{verbatim}


\newpage
\section*{Crediti}
%\addcontentsline{toc}{section}{Crediti}

Alcuni codici presenti in questa dispensa sono stati scritti dai miei colleghi Davide Le Bone e Giovanni Adelfio, alcuni contenuti per la sezione "Richiami di Algebra Lineare" sono presi dalla lezione di Salvo Romeo https://www.youtube.com/watch?v=fy8TUPrPnD0.

Questa dispensa si basa sugli argomenti trattati nel corso \textit{Metodi Numerici} tenuto dai professori Alessandro Alla e Gabriella Puppo presso l'Università La Sapienza (Roma), nel Corso di Laurea Scienze Matematiche per l'Intelligenza Artificiale nell'Anno Accademico 2024-2025.

\end{document}


\end{document}